\documentclass[10pt,oneside,openany]{book}
\usepackage[a4paper,top=22mm,bottom=23mm,left=23mm,right=23mm,headheight=15pt,headsep=7mm]{geometry}
\usepackage[T1]{fontenc}
\usepackage[utf8]{inputenc}
\usepackage{lmodern}
\usepackage{amsmath,amssymb,amsthm,mathtools}
\usepackage{microtype}
\usepackage{parskip}
\usepackage{enumitem}
\usepackage{xcolor}
\definecolor{Ink}{HTML}{172C45}
\definecolor{Accent}{HTML}{14686B}
\definecolor{RuleGray}{HTML}{C8D0D6}
\usepackage{graphicx,booktabs,longtable,array,calc,multirow}
\usepackage{needspace}
\usepackage{fancyvrb}
\usepackage{fancyhdr}
\usepackage{titlesec}
\usepackage{etoolbox}
\usepackage{xurl}
\usepackage[colorlinks=true,linkcolor=Accent,citecolor=Accent,urlcolor=Accent,bookmarksnumbered=true,pdfpagelabels=true,hypertexnames=false]{hyperref}
\usepackage{bookmark}
\hypersetup{pdftitle={Full BFGS on Convex Nonsmooth Objectives: Consolidated Proofs and Verification},pdfauthor={Working research consolidation},pdfsubject={Ordinary full BFGS, Armijo--weak-Wolfe, positive and negative convergence results},pdfkeywords={BFGS, nonsmooth convex optimization, convergence, counterexample, affine criticality}}
\providecommand{\tightlist}{\setlength{\itemsep}{0pt}\setlength{\parskip}{0pt}}
\providecommand{\passthrough}[1]{#1}
\providecommand{\pandocbounded}[1]{#1}
\newcommand{\ModuleTag}{I}
\renewcommand{\thechapter}{\ModuleTag}
\renewcommand{\thesection}{\ModuleTag.\arabic{section}}
\renewcommand{\thesubsection}{\thesection.\arabic{subsection}}
\renewcommand{\theequation}{\ModuleTag.\arabic{equation}}
\newtheorem{theorem}{Theorem}[chapter]
\newtheorem{lemma}{Lemma}[chapter]
\newtheorem{corollary}{Corollary}[chapter]
\newtheorem{proposition}{Proposition}[chapter]
\theoremstyle{remark}\newtheorem{remark}{Remark}[chapter]
\renewcommand{\thetheorem}{\ModuleTag.\arabic{theorem}}
\renewcommand{\thelemma}{\ModuleTag.\arabic{lemma}}
\renewcommand{\thecorollary}{\ModuleTag.\arabic{corollary}}
\renewcommand{\theproposition}{\ModuleTag.\arabic{proposition}}
\renewcommand{\theremark}{\ModuleTag.\arabic{remark}}
\setlength{\emergencystretch}{2em}
\setlength{\parindent}{0pt}
\setlength{\parskip}{5pt plus 1pt minus 1pt}
\setlength{\tabcolsep}{5pt}
\renewcommand{\arraystretch}{1.12}
\setlist{itemsep=2pt,topsep=4pt,parsep=0pt}
\setcounter{tocdepth}{1}
\setcounter{secnumdepth}{2}
\raggedbottom
\allowdisplaybreaks[2]
\titleformat{\chapter}[display]{\normalfont\color{Ink}}{\Large\sffamily\bfseries MODULE \ModuleTag}{5pt}{\huge\bfseries}
\titlespacing*{\chapter}{0pt}{0pt}{18pt}
\titleformat{\section}{\large\bfseries\color{Ink}}{\thesection}{.6em}{}
\titleformat{\subsection}{\normalsize\bfseries\color{Ink}}{\thesubsection}{.6em}{}
\titlespacing*{\section}{0pt}{15pt}{6pt}
\titlespacing*{\subsection}{0pt}{10pt}{4pt}
\pagestyle{fancy}
\fancyhf{}
\fancyhead[L]{\small\sffamily Full BFGS on convex nonsmooth objectives}
\fancyhead[R]{\small\sffamily Working proof record}
\fancyfoot[L]{\footnotesize\sffamily 10 September 2026}
\fancyfoot[R]{\thepage}
\renewcommand{\headrulewidth}{.3pt}
\fancypagestyle{plain}{\fancyhf{}\fancyfoot[L]{\footnotesize\sffamily Working proof record\enspace|\enspace10 September 2026}\fancyfoot[R]{\thepage}\renewcommand{\headrulewidth}{0pt}}
\newcommand{\StartModule}[3]{\cleardoublepage\renewcommand{\ModuleTag}{#1}\chapter{#2}\label{#1:chapter}\begin{small}\color{Accent}#3\end{small}\par\vspace{3pt}{\color{RuleGray}\hrule}\vspace{9pt}}
\begin{document}
\begin{titlepage}
\color{Ink}
{\large\sffamily\bfseries RESEARCH CONSOLIDATION}\par
\vspace{22mm}
{\Huge\bfseries Full BFGS on\\[3mm]convex nonsmooth\\[3mm]objectives}\par
\vspace{10mm}
{\Large Complete retained proofs, negative constructions,\\verification, and the remaining boundary}\par
\vspace{15mm}
{\color{Accent}\rule{\textwidth}{1.2pt}}\par
\vspace{7mm}
\begin{minipage}{.94\textwidth}
\large
Ordinary full BFGS. Fixed Armijo--weak-Wolfe conditions.\\
Objective-level convergence criteria. Deterministic counterexamples.\\
Planar affine criticality and arbitrary-dimensional fixed-normal geometry.
\end{minipage}
\vfill
\begin{minipage}{.94\textwidth}
\normalsize
\textbf{Status.} A consolidated working mathematical record, not an independently reviewed publication or a priority certification. The proof modules retain distinct assumptions. Exact finite certificates, numerical regressions, and analytic infinite-sequence arguments are identified separately.\par
\vspace{5mm}
\textbf{Fresh reproducibility audit:} 19 requested jobs passed.\par
\textbf{Record closed:} 10 September 2026.\par
\textbf{Companion material:} editable sources, original snapshots, verification scripts, historical and fresh logs, hashes, and supersession notes.
\end{minipage}
\end{titlepage}
\frontmatter
\chapter*{How to read this record}
\addcontentsline{toc}{chapter}{How to read this record}
This edition replaces the earlier inventory with a proof-bearing consolidation through the fixed-normal multidimensional continuation. Module I is the statement and supersession map. Modules M--U present the retained positive framework and its geometric applications; Modules D and L give the two incomparable negative constructions. Module E is the fresh verification record.

\textbf{The mathematical scope does not change between examples:} the algorithm is ordinary full BFGS, not a limited-memory or modified substitute. Positive statements concern infinite admissible differentiable-iterate runs unless stated otherwise. The counterexamples fix a particular unit-first weak-Wolfe doubling/bisection implementation and verify every oracle decision.

\textbf{Not one maximal function class.} The strongest results are not all ordered by implication. Compact-level geometry does not subsume noncompact homogeneous behavior. Positive lower curvature and higher branch differentiability are separate hypotheses. Planar curved geometry does not encompass arbitrary dimension, and fixed normals do not encompass rotating normals. The report explicitly retains these distinctions.

\textbf{Detailed proofs, with shared algebra centralized.} The common principal-minor and energy foundation is proved once in Module M. Source proof modules then refer to that foundation rather than repeating the appendix. The selected modules preserve their original internal section and theorem ordering with a module prefix; occasional numbering gaps correspond to omitted superseded or historical material. The complete original snapshots remain available for comparison.

\textbf{Disclosed consolidation deductions.} First, the older general two-branch scalar-ridge result is a corollary of the later planar theorem; the implication is proved in Module I, and its independent source proof is archived rather than repeated as a principal result. Second, the new superquadratic reciprocal estimate allows an $O(N^2)$ reference-energy correction. The same ordered-tail bound also strengthens the recorded conservative rates to little-$o$ corollaries in Module F. No unproved theorem is created by combining the two different local counterexamples.

\textbf{Verification is not peer review.} All 19 selected jobs passed in a fresh isolated run. This establishes reproducibility of their finite calculations. The root certificates, continuum geometric arguments, ordered limits and stable-manifold existence steps have different roles; their limits are detailed in Module E. Neither the results nor their novelty have been independently certified.

\begin{center}
\begin{tabular}{p{.16\textwidth}p{.74\textwidth}}
\toprule
Read first & What it contains\\\midrule
I & The strongest statements, exact quantifiers, supersessions and unresolved cases.\\
M, T, R & Shared BFGS algebra, objective certificates and radial interpretation.\\
C, U & New geometric local/cluster theorems: planar and fixed-normal multidimensional.\\
D, L & Exact deterministic failure and persistence with positive lower curvature.\\
E & Fresh audit, arithmetic classification, commands and limitations.\\\bottomrule
\end{tabular}
\end{center}
\tableofcontents
\mainmatter
\StartModule{I}{Strongest results and logical map}{Editorial synthesis of the source record. Full proofs follow; source snapshots and explicit supersession notes are included in the companion package.}

% BEGIN original chapters/inventory.tex
\textbf{Record closed: 10 September 2026.} This edition includes the deterministic-bracketing, positive-lower-curvature, planar affine-critical-cluster, and arbitrary-dimensional fixed-normal continuations. It supersedes the earlier theorem inventory, not the archived source proofs.

\section*{Status and reading convention}\label{status-and-reading-convention}

These are mathematical claims developed in the supplied working reports. They have not received independent mathematical review or an exhaustive publication-priority audit. This consolidation distinguishes complete working arguments, computer-assisted algebraic certificates, finite algebraic regressions, numerical experiments, and unresolved conjectures. A successful verifier is not treated as a proof of an infinite-sequence theorem or of a geometric joining argument.

The accompanying PDF contains the retained proof sections, not only this inventory. Its proof modules preserve source terminology and internal numbering, with a module prefix to make references unique. Shared BFGS algebra is proved once in Module M. Where a source proof is condensed or replaced by an implication, that change is identified. Original source snapshots remain in \path{provenance/} for comparison; obsolete statements in those snapshots are not the operative conclusions of this edition.

No new literature survey was performed for this consolidation. Published results mentioned in the source reports remain background, and are not represented as discoveries of this investigation.

\section{Algorithm, quantifiers, and notions of convergence}\label{algorithm-quantifiers-and-notions-of-convergence}

The algorithm is \textbf{ordinary full BFGS} in finite-dimensional Euclidean space:

\[
 g_k=\nabla F(x_k),\qquad s_k=-\alpha_kH_kg_k,\qquad
 y_k=g_{k+1}-g_k,\qquad H_0\succ0,
\]

\[
 H_{k+1}=\left(I-\frac{s_ky_k^T}{y_k^Ts_k}\right)H_k
 \left(I-\frac{y_ks_k^T}{y_k^Ts_k}\right)
 +\frac{s_ks_k^T}{y_k^Ts_k}.
\]

An \emph{infinite admissible run} has nonzero gradients at differentiable accepted iterates, positive lengths, and fixed Armijo--weak-Wolfe constants \(0<c_1<c_2<1\):

\[
 F(x_{k+1})\le F(x_k)+c_1g_k^Ts_k,\qquad
 g_{k+1}^Ts_k\ge c_2g_k^Ts_k.
\]

No damping, restart, skipped update, memory truncation, proximal operation, imposed spectral limit, or assumed iterate convergence is introduced. Positive theorems quantify over all such infinite runs unless a more specific implementation is stated.

Three conclusions must remain distinct:

\[
 F(x_k)\to F_*;\qquad \operatorname{dist}(x_k,X_*)\to0;\qquad x_k\to x_*.
\]

On a compact sublevel, objective convergence implies distance to the solution set tends to zero. Uniqueness then gives whole-sequence convergence. Affine-fiber nonuniqueness has a separate lifting theorem; arbitrary nonunique solution sets do not inherit that conclusion.

A zero gradient at a differentiable point is successful termination. A ray with no differentiable acceptable point is not an infinite admissible run. The positive statements must not silently be strengthened into universal all-ray existence or finite oracle complexity.

\section{The negative results that must both be retained}\label{the-negative-results-that-must-both-be-retained}

\subsection{N1. Deterministic local failure with globally C2,1 convex branches}\label{n1.-deterministic-local-failure-with-globally-c21-convex-branches}

There is one fixed coercive \(F:\mathbb R^2\to\mathbb R\), with a unique nonsmooth minimizer at zero, of the form

\[
 F=\max_{0\le i\le3}f_i,\qquad f_i\in C^{2,1}(\mathbb R^2),\qquad
 0\preceq\nabla^2f_i\preceq64I.
\]

Every branch Hessian is globally Lipschitz. With the fixed initialization \(H_0=I\), a sequence of bad starts approaches the minimizer. The \textbf{actual deterministic unit-first weak-Wolfe doubling/bisection search} produces nonstationary converged limits:

\[
 x_0^{(j)}\to0,\qquad x_k^{(j)}\to p_j\ne0,\qquad
 F(x_k^{(j)})\downarrow F(p_j)>0,\qquad0\notin\partial F(p_j).
\]

The same objective, starts, and trial decisions work for every fixed \(0<c_1\le1/3\), \(c_1<c_2<1\). Every accepted and rejected query is differentiable. Accepted lengths repeat \(1,1/4,1/16\), with at most five trials per iteration. The entire query trajectory can remain in an arbitrarily small neighborhood of the true minimizer. Initial objective gaps also tend to zero.

This is a failure of a deterministic neighborhood guarantee, not a claim that every close start fails. It is proved by an exactly isolated algebraic cycle, exact matrix closure, a scalar convex interpolation, all-trial containment, and scaling induction. \textbf{Proof: Module D.}

\subsection{N2. Deterministic local failure with positive lower curvature}\label{n2.-deterministic-local-failure-with-positive-lower-curvature}

There is also a strongly convex finite objective, with a unique nonsmooth minimizer and compact sublevels, such that

\[
 0<\mu I\preceq\nabla^2F(x)\preceq LI
\]

wherever its classical Hessian exists, yet the same specified deterministic implementation admits bad starts arbitrarily close to the minimizer. A single fixed affine change of coordinates makes the common initial matrix \(I\).

One version is a maximum of four globally \(C^{1,1}\), uniformly strongly convex, gradient-Lipschitz branches. The individual branches \textbf{are not twice differentiable at the accumulation point}. Another version has one constant SPD classical Hessian wherever it exists, but a countably piecewise-affine nonsmooth component rather than finitely many smooth branches.

The source uses \(c_1=10^{-4}\), \(c_2=1/2\). Its bad seed is an existence object on a certified hyperbolic stable manifold, not a displayed rounded point claimed to lie exactly on that manifold. \textbf{Proof: Module L.}

\textbf{N1 and N2 are not merged into a fictitious stronger theorem.} N1 has more branch differentiability and no uniform positive lower curvature. N2 has positive lower curvature and less differentiability at the accumulation point. Neither proves local failure for a finite maximum of uniformly strongly convex \(C^2\) branches through the optimizer.

\subsection{N3. Simple global counterexamples retained as corollaries}\label{n3.-simple-global-counterexamples-retained-as-corollaries}

The deterministic seed gives failure on a coercive finite maximum of four affine functions, with a unique minimizer. This retains global failure in the finite polyhedral and semialgebraic class. It is not superseded by an accumulation-based local example.

Adding a common positive quadratic to the three-piece seed yields

\[
 Q_\mu(x)=\frac\mu2\|x\|^2+\max_{i=0,1,2}g_i^Tx.
\]

For every \(\mu>0\), suitable initial states produce deterministic nonstationary convergence to zero. Each branch is quadratic, analytic, and has Hessian exactly \(\mu I\). The true minimizer lies elsewhere. This is \textbf{global failure}, and is consistent with local convergence near the true minimizer. The stable-manifold argument is required; the unperturbed orbit cannot simply be reused. \textbf{Proofs: D.3.1 and Module L.}

\section{General positive mechanisms: retain distinct hypotheses}\label{general-positive-mechanisms-retain-distinct-hypotheses}

\subsection{P1. Objective-level transported certificates}\label{p1.-objective-level-transported-certificates}

For differentiable endpoints define \(g=\nabla F(x)\), \(y=\nabla F(z)-g\), \(s=z-x\), \(d=-g^Ts\), \(\Delta=F(x)-F(z)\), and \(t=y^Ts\). An eligible pair satisfies

\[d>0,\qquad c_1d\le\Delta\le d,\qquad t\ge(1-c_2)d.\]

On each compact positive objective band, it suffices that a finite-valued PSD field \(W\) satisfy

\[g^TW(x)g\ge m>0,\qquad y^TW(z)y\le Ct,\]

\[\Pi_sW(z)\Pi_s\preceq\Pi_sW(x)\Pi_s,\]

for every eligible pair with sufficiently small \(d\). The latter two constants and the field may depend on the band and on the fixed line-search parameters. No continuity, uniform boundedness, or positive definiteness of \(W\) is required.

Then every admissible run on the compact sublevel has value and solution-set convergence; uniqueness gives iterate convergence. A local structural neighborhood gives a trapped low-sublevel basin. A condition on one-sided shells around possible positive limiting levels suffices as well.

The normalization with exact projected nonexpansion is equivalent to allowing \((1+Ld)\): multiply \(W\) by \(\exp((L/c_1)(F-F_*))\). \textbf{Proof: Module T; normalization and geometry: Module G.}

RS is the radial specialization \(W(x)=(x-x_*)(x-x_*)^T\). Its equivalent inward Taylor condition is useful for checking functions, but is not a separate strongest certificate theorem. The inclusion of RS is strict. The transported condition itself is not necessary for actual convergence.

\subsection{P2. Moving frames with controlled cumulative transport}\label{p2.-moving-frames-with-controlled-cumulative-transport}

A frame satisfying \(V_{k+1}=V_kT_k+s_kb_k^T\) has an exact BFGS comparison-energy identity. Uniform visibility, bounded secant correction, and bounded cumulative transport products and inverses exclude positive-level stagnation. This trajectory mechanism need not produce a single position-only certificate. \textbf{Proof: Module M.}

\subsection{P3. Imperfect transport in fixed smooth directions}\label{p3.-imperfect-transport-in-fixed-smooth-directions}

If a fixed subspace \(S\) satisfies \(\|P_Sy_k\|^2\le L_Sy_k^Ts_k\), the actual update gives \(\operatorname{tr}(P_SB_k)=O(k)\). Frame residuals are then permitted:

\[V_{k+1}=V_kT_k+s_kb_k^T+R_k,\]

\[\|T_k\|^2\le1+Cd_k,\qquad \operatorname{range}R_k\subset S,
\qquad\|R_k\|_F^2\le Cd_k.\]

Together with uniform visibility and bounded secant correction, the resulting \(O(N^2)\) error budget is defeated by the superquadratic reciprocal theorem below. This extends the proof mechanism; \textbf{a strict inclusion between the existential function classes of P1 and P3 has not been proved}. \textbf{Proof: Module M; objective applications: Module U.}

\section{Strongest concrete local and global positive families}\label{strongest-concrete-local-and-global-positive-families}

\subsection{P4. Radially regular and finite-partition geometry}\label{p4.-radially-regular-and-finite-partition-geometry}

On a compact positive band, RS is equivalent to uniformly quadratic inward first-order modeling error:

\[
 0\le F(x-\theta r)-F(x)+\theta\nabla F(x)^Tr
 \le(M_\eta/2)\theta^2,\qquad r=x-x_*.
\]

The constants may deteriorate as the band approaches the minimum. This yields local convergence for a finite convex function whose neighborhood of a unique minimizer has a \textbf{finite polyhedral partition with \(C^{1,1}\) extensions on each cell}. Individual cell extensions need not be convex outside their cells; the full objective is convex.

Consequences include finite convex PLQ objectives on polyhedral cells; smooth convex plus finite polyhedral terms; appropriate absolute-value, hinge, quantile, fused and generalized \(\ell_1\) terms; and local mixtures of Euclidean group residuals. A finite maximum of arbitrary quadratics is not automatically PLQ.

Centered local representations

\[F(x_*+z)=\psi(z)+\gamma(z)\]

are covered when \(\psi\) is convex \(C^{1,1}_{\rm loc}\) and \(\gamma\) is finite convex degree-one homogeneous. A global representation on a compact initial sublevel gives global value convergence; uniqueness gives iterate convergence. Neither original summand must separately be minimized at the center.

Also retained are finite sums of powered nonnegative gauges, and increasing convex smooth transforms of coercive gauges. Different powers, including \(1<p<2\), and coupled linear maps are permitted under the stated source hypotheses. General spectral regularization at a nonzero rank-deficient solution is not automatic. \textbf{Proofs: Modules R and F.}

For smooth-plus-polyhedral objectives, retain the explicit sufficient basin threshold obtained by minimizing the objective over the active regions of pieces inactive at the optimizer. For nonunique affine fibers, retain the exact linear-image lifting formula. These are not replaced by a generic uniqueness statement.

\subsection{P5. Planar affine-critical cluster points}\label{p5.-planar-affine-critical-cluster-points}

Let \(F:\mathbb R^2\to\mathbb R\) be finite convex, have an attained minimum and a compact sublevel, and have strictly convex positive sublevels. Away from affine-critical points, assume either local convex \(C^2\) smoothness or a local maximum of two convex \(C^2\) branches whose active gradients are linearly independent.

Then every infinite admissible run has an affine-critical cluster point. Adding no false affine criticality on the local basin gives convergence. For finite smooth maxima, affine independence and strictly positive optimal barycentric weights at a unique optimizer make the required local geometry automatic in the plane. Strictly convex low contours then suffice; uniformly strongly convex branches are a stronger convenient assumption. \textbf{Complete proof: Module C.}

This does not cover every planar finite maximum: rank-degenerate ties and flat positive contours are outside the stated theorem. Nor does it prove the arbitrary-dimensional rotating-normal conjecture.

\subsection{An explicit supersession established in this consolidation}\label{an-explicit-supersession-established-in-this-consolidation}

The older theorem for

\[F(u,v)=\max\{a(u)+pv,b(u)-qv\},\qquad p,q>0,\]

with convex \(C^2\) profiles, coercive \(\phi=(qa+pb)/(p+q)\) having unique minimizer \(u_*\), and \(a'',b''>0\) off \(u_*\), is a \textbf{corollary of P5}, not an independent strongest theorem.

Here is the implication. The profiles are strictly convex. If two different points of a positive sublevel have different horizontal coordinates, strict profile convexity puts the interior of their segment strictly below the level. If their horizontal coordinates agree, their vertical coordinates differ; the two nonzero normal slopes again make every strict convex combination satisfy both inequalities strictly. Thus every positive sublevel is strictly convex. At a tie, the determinant of the active gradients is \(-qa'(u)-pb'(u)=-(p+q)\phi'(u)\). It is nonzero except at \(u_*\), where the tied point is precisely the optimizer. Smooth points are already regular. Hence P5 applies globally, and its localization argument gives the old local result.

This is a disclosed deduction of a relationship between already recorded theorems. The original independent moving-ridge proof and its verification remain in the provenance archive, but are not repeated as a standalone principal result.

\subsection{P6. Arbitrarily many parabolic branches, including flat contours}\label{p6.-arbitrarily-many-parabolic-branches-including-flat-contours}

For

\[F(u,v)=\max_{i\le m}\{a_i u^2+b_iv\},\qquad a_i>0,
\quad\min_i b_i<0<\max_i b_i,\]

all admissible full-BFGS runs converge globally to zero. Branches with \(b_i=0\) are allowed. Such branches can create flat contour segments, so P5 does \textbf{not} subsume the whole family. Its explicit certificate also proves strict inclusion beyond eligible-pair RS. \textbf{Full proof: Module Q.}

\subsection{P7. Arbitrary-dimensional fixed-normal geometry}\label{p7.-arbitrary-dimensional-fixed-normal-geometry}

Let \(P=\{x:a^Tx=b\}\) be an exceptional hyperplane through a unique minimizer, with \(S=a^\perp\). On a neighborhood of a compact initial sublevel, assume strong convexity, a single-valued Lipschitz tangential subgradient projection \(P_S\partial F(x)=\{A(x)\}\), and local \(C^{1,1}\) regularity off \(P\). Every admissible run converges.

The full normal gradient need not have a uniform Lipschitz constant near the hyperplane. A local version uses a low trapped sublevel. For \(F(u,v)=\psi(u,v)+h(v)\), full strong convexity can be weakened to tangential strong convexity of \(\psi\), strict positive contours, compactness, and a unique optimizer with \(v_*=0\).

A separate cluster theorem permits a closed collection of \textbf{parallel genuinely nonsmooth interfaces}. It forces one affine-critical cluster point under strong convexity and a Lipschitz tangential projection. Add no false affine criticality to conclude optimality. The single-hyperplane theorem is not a literal specialization: it allows differentiable but non-\(C^{1,1}\) exceptional points.

The arbitrary-dimensional result is not a planar reduction or a block-diagonal algorithm. Its real restriction is the fixed normal direction. The explicit normal-curvature-spike examples lie outside RS. \textbf{Complete proofs and infinite-series oracle: Module U.}

\subsection{P8. Scalar convergence}\label{p8.-scalar-convergence}

In one dimension, finite convex objectives with a unique minimizer and compact initial sublevel have convergence of every infinite admissible positive-metric Armijo--weak-Wolfe run, independently of the BFGS update. This remains nonredundant because it covers functions lacking every local transported certificate. \textbf{Proof and nonnecessity example: Module G.}

\section{Noncompact homogeneous theory remains independent}\label{noncompact-homogeneous-theory-remains-independent}

For every finite convex degree-one homogeneous objective, every infinite admissible run satisfies

\[F(x_k)\to0\quad\text{or}\quad F(x_k)\to-\infty.\]

No iterate, step, matrix, or length bound is assumed, and the supporting gradient set may be infinite. Bounded-below objectives therefore have optimal value convergence; coercive ones have \(x_k\to0\).

For finite max-linear objectives, retain three refinements: a finite-value-limit run has zero in the affine hull of its recurrent gradients; recurrent gradient span of dimension at most two gives whole-iterate convergence on the finite-value branch; and purely flat-direction noncompactness admits an explicit metric-projection limit. General affine-fiber lifting gives the same type of result for reduced nonhomogeneous objectives.

The noncompact zero-value iterate branch in recurrent dimension at least three is \textbf{not completely classified}. A compact-level certificate theorem cannot replace these noncompact statements. \textbf{Full proofs: Module H.}

The two-piece finite-escape theorem from Xie--Wächter is pre-existing background and remains stronger in its finite-termination conclusion than the new many-piece infinite-run escape statement. The user's paper is supplied separately, not rebranded as a new result.

\section{Shared algebra and quantitative record}\label{shared-algebra-and-quantitative-record}

The principal-minor inequality is proved once in Module M:

\[e_j(H_{k+1})\le
\frac{\|y_k\|^2}{\beta_k^2q_k}e_{j+1}(H_k)
+\frac{\alpha_k}{\beta_k}e_j(H_k),\quad
q_k=g_k^TH_kg_k,\quad\beta_k=\frac{y_k^Ts_k}{d_k}.\]

With bounded gradients, fixed weak Wolfe and \(d_k\to0\), both the geometric and arithmetic means of \(q_k^{-1}\) diverge. Under the stronger assumption \(\sum d_k<\infty\), the final recorded strengthening is

\[\boxed{\frac1{N^2}\sum_{k<N}q_k^{-1}\to\infty.}\]

Retain the explicit finite-horizon bound as well; an asymptotic statement does not replace its finite constants. The former \(\liminf q_k=0\) and \(\Omega(N^2)\) headlines are retired as main results. They remain consequences used by older proofs.

The fixed-reference identity is

\[E_{k+1}=E_k-\frac{[g_k^T(x_k-x_*)]^2}{q_k}
+\frac{[y_k^T(x_{k+1}-x_*)]^2}{y_k^Ts_k}.\]

A cumulative correction \(o(\sum q_k^{-1})\) excludes a positive gap. In particular, by the final reciprocal result, an \(O(N^2)\) cumulative correction suffices. This last sentence is an explicit consolidation consequence of the two displayed inequalities, not a claim of an additional source theorem.

The original rate note recorded big-\(O\) corollaries. The later superquadratic reciprocal theorem strengthens them to \(o(N^{-1/2})\) for uniformly bounded radial correction, \(o(N^{-1})\) for bounded-below degree-one homogeneous objectives, and \(o(\log N/N)\) for positive-definite quadratic plus coercive gauge. These are disclosed consolidation deductions, proved in Module F while preserving the stronger explicit finite-horizon inequalities. They are not uniform rates over arbitrary certificates, sharpness statements, or line-search evaluation bounds.

\section{Geometric findings retained independently of convergence}\label{geometric-findings-retained-independently-of-convergence}

The condition

\[0\in\operatorname{aff}\partial F(p)\Longrightarrow0\in\partial F(p)\]

means that every nonoptimal nearby point admits a direction along which all subgradients predict the same strictly negative rate. Affine independence and strictly positive optimal coefficients make this true locally for finite smooth maxima, with uniform positive-band affine separation. Pointwise separation alone need not be uniform when ranks change.

For finite convex semialgebraic objectives, the affine-critical set has finitely many connected components, and the objective is constant on each. Thus only finitely many affine-critical values occur, giving no false affine criticality below a sufficiently low threshold. This is a geometric proposition, not a blanket BFGS theorem. \textbf{Proofs: Module A.}

The incoming-jump obstruction rules out all PSD certificates when the target gradient is a vanishing-cost linear combination of several eligible jumps arriving at the same target. It neither assumes regularity of the certificate nor constructs one common BFGS history. Scalar accumulating-kink examples show certificate failure despite convergence. These distinctions are retained in Module G.

\section{Verification and unresolved boundaries}\label{verification-and-unresolved-boundaries}

The fresh verification package preserves the source scripts and runs them in separate writable directories. Historical JSON and logs are never overwritten or counted as new runs. \path{verification/fresh/SUMMARY.json} records commands, script hashes, environment, exit codes, elapsed time, and output paths. The final audit report classifies each job by exact rational, symbolic, interval, mixed rational/decimal, or non-interval high-precision arithmetic.

The exact deterministic cycle, its smooth embedding, the hyperbolicity bounds, principal-minor identities, frame identities, special examples and full-objective regressions have different evidential roles. In particular, finite stable-eigenspace calculations plus the stated stable-manifold theorem supply an existence argument; no rounded initial point is certified as lying on that manifold. Geometric joining and infinite tail arguments are analytic proof steps, not tested asymptotic facts.

The remaining central gaps are: arbitrary-dimensional rotating-normal smooth-max geometry; the general local theorem for finite uniformly strongly convex \(C^2\) branches without the recorded nondegeneracy/geometric assumptions; generic or positive-measure failure; other fixed line-search implementations; arbitrary nonunique solution-set iterate convergence; noncompact high-rank homogeneous zero-value iterates; and sharp rates or universal matrix asymptotics.

The blanket question is not answered by calling all semialgebraic functions safe: global failure already occurs for affine maxima. The planar and fixed-normal positive theorems have precise additional assumptions. The latest lower-curvature counterexamples likewise do not establish local failure in a finite analytic or semialgebraic description.

\section{Supersession policy}\label{supersession-policy}

\begin{longtable}[]{@{}
  >{\raggedright\arraybackslash}p{(\columnwidth - 2\tabcolsep) * \real{0.5000}}
  >{\raggedright\arraybackslash}p{(\columnwidth - 2\tabcolsep) * \real{0.5000}}@{}}
\toprule\noalign{}
\begin{minipage}[b]{\linewidth}\raggedright
Earlier item
\end{minipage} & \begin{minipage}[b]{\linewidth}\raggedright
Operative treatment
\end{minipage} \\
\midrule\noalign{}
\endhead
\bottomrule\noalign{}
\endlastfoot
Prescribed acceptable-length counterexamples & Replaced by the final deterministic-bracketing seed and all-trial embeddings. \\
Alternative deterministic seeds & Do not mix constants. Only the final \((1,0),(-1,-1/8),(c,d)\) seed is canonical. \\
Upper-regularity local failure versus lower-curvature local failure & Retain both: differentiability and lower-curvature assumptions are incomparable. \\
Original scalar-ridge convergence headline & Corollary of the later planar theorem; the implication is proved above. Independent proof retained only for provenance. \\
All parabolic multi-branch convergence & Retain: zero-normal-slope pieces can violate the planar strict-contour hypothesis. \\
Separate \(\ell_1\) and \(u^2+|v|\) headlines & Examples of broad radial/partition families; retain special rate and matrix-limit facts as corollaries. \\
RS as the broadest sufficient condition & Replaced by the transported criterion; retain its natural equivalent and verifiable subclasses. \\
Exact versus smooth-subspace-imperfect frame transport & Retain the distinct hypotheses; no unproved strict inclusion of function classes is asserted. \\
Earlier reciprocal lower bounds & Use the superquadratic conclusion under summability, but retain the weaker-assumption theorem and finite-horizon constants. \\
Noncompact homogeneous theorems & Retain independently of compact-level results. \\
Certificate counterexamples & Keep the general obstruction and strongest nonnecessity example; do not label these as bad BFGS trajectories. \\
\end{longtable}

The scope of the positive and negative programs is now substantially clearer, but neither gives a necessary-and-sufficient classification of nonsmooth full-BFGS convergence.

% END original chapters/inventory.tex

\StartModule{M}{Shared BFGS algebra and energy budgets}{Unified complete proof of the principal-minor hierarchy, reciprocal-growth theorems, fixed and moving comparison energies, and the smooth-subspace residual budget.}

% BEGIN original chapters/foundation.tex
\section{Notation, elementary consequences, and localization}
\label{M:notation}
This module consolidates the common algebra from sources~\cite{SM,SH,SQ,SU}. No outside BFGS convergence theorem is substituted for this algebra. All results inherit the infinite-admissible-run convention in the inventory.

Put
\[
 \begin{gathered}
 B_k=H_k^{-1},\qquad q_k=g_k^TH_kg_k>0,\qquad d_k=-g_k^Ts_k=\alpha_kq_k,\\
 t_k=y_k^Ts_k,\qquad\Delta_k=F(x_k)-F(x_{k+1}),\qquad c=1-c_2>0.
 \end{gathered}
\]
Convexity, Armijo and weak Wolfe give
\begin{equation}\label{M:basic}
 c_1d_k\le\Delta_k\le d_k,\qquad t_k\ge c d_k>0.
\end{equation}
In particular $\sum d_k<\infty$ on a run with finite limiting value. The Hessian form is
\begin{equation}\label{M:hessian}
 B_{k+1}=B_k-\frac{g_kg_k^T}{q_k}+\frac{y_ky_k^T}{t_k}.
\end{equation}
Metric Cauchy--Schwarz proves
\begin{equation}\label{M:annihilation}
 M_k:=B_k-\frac{g_kg_k^T}{q_k}\succeq0,\qquad M_ks_k=0.
\end{equation}
For example $(g_k^Tz)^2\le q_k z^TB_kz$, and $B_ks_k=-\alpha_kg_k$ proves the second identity.

A finite convex function has uniformly bounded subgradients on any compact set $K$. Indeed, choose $\epsilon>0$ and bound its continuous values on $K+\epsilon\overline B$. The supporting inequality at $x$ evaluated at $x+\epsilon u$ bounds $g^Tu$ for every unit $u$, hence $\|g\|$. This does not require differentiability throughout $K$.

\begin{lemma}[A low sublevel is an invariant local basin]\label{M:localization}
Let $F$ be finite convex and have a unique minimizer $x_*$. Inside any specified structural neighborhood, choose $r>0$ with $\overline B_r(x_*)$ contained in it. Then
\[
 \omega_r=\min_{\|z\|=r}\{F(x_*+z)-F(x_*)\}>0,
\]
and $F(x_0)-F(x_*)<\omega_r$ places the entire initial sublevel in $B_r(x_*)$.
\end{lemma}
\begin{proof}
For $\|x-x_*\|\ge r$, convexity along the radial segment gives
\[
 F(x)-F(x_*)\ge\frac{\|x-x_*\|}{r}\omega_r\ge\omega_r.
\]
Continuity, compactness of the sphere, and uniqueness give $\omega_r>0$. Armijo keeps every accepted iterate in the initial sublevel. No bound on $H_0$ is needed, and rejected trials need not remain in the neighborhood. This argument also proves boundedness of every sublevel for a finite convex function with an attained unique minimizer.
\end{proof}

\begin{lemma}[Strict contours give vanishing steps on a positive limiting level]\label{M:strictcontours}
Suppose the run is contained in a compact sublevel, $F(x_k)\downarrow\ell>F_*$, and the boundary of $\{F\le\ell\}$ contains no nontrivial line segment. Then $s_k\to0$.
\end{lemma}
\begin{proof}
Otherwise, along a subsequence $x_k\to x$, $x_{k+1}\to z\ne x$ and $g_k\to g\in\partial F(x)$. Since $d_k\to0$, $F(x)=F(z)=\ell$ and $g^T(z-x)=0$. The supporting inequality and convexity imply that $F=\ell$ on $[x,z]$, contrary to the contour assumption.
\end{proof}

\section{The principal-minor inequality: complete algebra}
\label{M:principal}
For SPD $H$, let $e_j(H)$ be the elementary symmetric polynomial of degree $j$ of its eigenvalues, with $e_{n+1}=0$. Thus $e_1=\operatorname{tr}$ and $e_n=\det$.

\begin{theorem}[Principal-minor hierarchy]\label{M:hierarchy}
For an ordinary inverse-BFGS update with $s=-\alpha Hg$, $q=g^THg>0$ and $\beta=-y^THg/q>0$,
\begin{equation}\label{M:minor}
 e_j(H_+)\le \frac{\|y\|^2}{\beta^2q}e_{j+1}(H)+\frac\alpha\beta e_j(H),\qquad1\le j\le n.
\end{equation}
\end{theorem}
\begin{proof}
Set $v=Hg$, $\kappa=g^THy=-\beta q$, $P=I-vy^T/\kappa$ and $\mathcal F=PHP^T$. The update is
\[
 H_+=\mathcal F+\frac{\alpha}{\beta q}vv^T.
\]
Let $Z$ have $j$ orthonormal columns, and put
\[
 A=Z^THZ,\ a=Z^THg,\ b=Z^THy,\ r=y^THy,
 \quad u=a^TA^{-1}a,\quad w=a^TA^{-1}b,\quad z=b^TA^{-1}b.
\]
The compression is $Z^T\mathcal FZ=A-(ab^T+ba^T)/\kappa+(r/\kappa^2)aa^T$. A rank-two determinant expansion gives
\[
 \det(Z^T\mathcal FZ)=\frac{\det A}{\kappa^2}\bigl((\kappa-w)^2+u(r-z)\bigr).
\]
If $\operatorname{Gram}_H$ denotes the Gram matrix in the $H$ inner product, Schur complements give
\[
 \det\operatorname{Gram}_H(y,Z)=\det A(r-z),
\]
\[
 \det\operatorname{Gram}_H(g,y,Z)
 =\det A\bigl((q-u)(r-z)-(\kappa-w)^2\bigr).
\]
Consequently
\begin{equation}\label{M:compression}
 \det(Z^T\mathcal FZ)
 =\frac{q\det\operatorname{Gram}_H(y,Z)-\det\operatorname{Gram}_H(g,y,Z)}{\beta^2q^2}
 \le\frac{\det\operatorname{Gram}_H(y,Z)}{\beta^2q}.
\end{equation}
All Gram determinants are nonnegative. Adding the rank-one term changes the compressed determinant by exactly
\[
 \frac{\alpha}{\beta q}\det A\,u\le\frac\alpha\beta\det A,
\]
because $u\le q$. The identities remain valid for singular larger Gram matrices or singular $Z^T\mathcal FZ$; only $A$ must be invertible, which it is.

Choose an orthonormal eigenbasis of $H$, with eigenvalues $\lambda_i$, and sum over coordinate $j$-subsets $I$. Then
\[
 \sum_{|I|=j}\det\operatorname{Gram}_H(y,E_I)
 =\sum_{|J|=j+1}\left(\prod_{i\in J}\lambda_i\right)\sum_{i\in J}y_i^2
 \le\|y\|^2e_{j+1}(H).
\]
The rank-one increases sum to at most $(\alpha/\beta)e_j(H)$. This proves \eqref{M:minor}. For $j=n$, the calculation gives the exact identity $\det H_+=(\alpha/\beta)\det H$.
\end{proof}

\section{Reciprocal growth: the strongest recorded versions}
\label{M:reciprocal}
\begin{theorem}[Bounded gradients and vanishing predicted decrease]\label{M:average}
If an infinite full-BFGS run has bounded gradients, fixed weak Wolfe and $d_k\to0$, then
\[
 \left(\prod_{k<N}q_k^{-1}\right)^{1/N}\to\infty,
 \qquad\frac1N\sum_{k<N}q_k^{-1}\to\infty.
\]
\end{theorem}
\begin{proof}
Choose $G>0$ with $\|g_k\|\le G$ and write $K_0=4G^2/c^2$. Since $\beta_k\ge c$ and $\|y_k\|\le2G$, the hierarchy gives
\[
 e_j(H_{k+1})\le q_k^{-1}\{K_0 e_{j+1}(H_k)+(d_k/c)e_j(H_k)\}.
\]
Fix $\eta>0$ and a tail $k\ge J$ on which $d_k/c\le\eta$. Define
\[
 V_\eta(H)=\sum_{j=1}^n(K_0/\eta)^{j-1}e_j(H)\ge\operatorname{tr}H.
\]
Multiplying and summing yields $V_\eta(H_{k+1})\le(2\eta/q_k)V_\eta(H_k)$, while $q_k\le G^2V_\eta(H_k)$. Iterating through the penultimate step of a block and using the last endpoint gives
\[
 \prod_{k=J}^{J+m-1}q_k^{-1}\ge\frac1{G^2V_\eta(H_J)(2\eta)^{m-1}}.
\]
Taking $m$th roots and then allowing the fixed $\eta$ to be arbitrarily small proves divergence of the geometric mean; a fixed prefix has no effect. Arithmetic--geometric mean proves the other assertion.
\end{proof}

\begin{theorem}[Finite horizon and superquadratic growth under summability]\label{M:superquadratic}
Assume bounded gradients, fixed weak Wolfe, and $\sum_kd_k<\infty$. For a fixed tail $J$, let $D_J=\sum_{k\ge J}d_k>0$. For every $\eta>0$ and $m\ge2$,
\begin{equation}\label{M:finitehorizon}
 \sum_{k=J}^{J+m-1}q_k^{-1}\ge
 m[G^2V_\eta(H_J)]^{-1/m}
 \left(\eta+\frac{D_J}{c(m-1)}\right)^{-(m-1)/m}.
\end{equation}
In particular,
\begin{equation}\label{M:superquad}
 \frac1{N^2}\sum_{k<N}q_k^{-1}\longrightarrow\infty.
\end{equation}
\end{theorem}
\begin{proof}
Without any tail restriction on $d_k$, the weighted hierarchy gives
\[
 V_\eta(H_{k+1})\le\frac{\eta+d_k/c}{q_k}V_\eta(H_k).
\]
Iterate for $m-1$ steps, use the endpoint bound for $q_{J+m-1}$, and apply arithmetic--geometric mean:
\[
 \prod_{k=J}^{J+m-1}q_k
 \le G^2V_\eta(H_J)\prod_{k=J}^{J+m-2}(\eta+d_k/c)
 \le G^2V_\eta(H_J)\left(\eta+\frac{D_J}{c(m-1)}\right)^{m-1}.
\]
Arithmetic--geometric mean for the reciprocals proves \eqref{M:finitehorizon}. Choose $\eta=D_J/[c(m-1)]$. With $J$ fixed, $V_\eta(H_J)$ grows at most polynomially in $m$, of degree $n-1$. Its $m$th root tends to one. Hence
\[
 \liminf_{N\to\infty}N^{-2}\sum_{k<N}q_k^{-1}\ge \frac{c}{2D_J}.
\]
The fixed prefix does not affect the normalization. Finally let $J\to\infty$; then $D_J\to0$. This proves \eqref{M:superquad}. No uniform bound on $H_J$ as $J$ varies is used: the tail is fixed before the horizon limit.
\end{proof}
These conclusions do not say $H_k\to0$, $q_k\to0$ pointwise, or that eigenvectors converge. The finite-horizon bound and the weaker-assumption geometric-mean theorem are retained alongside the strongest summable-budget asymptotic statement.

\section{Fixed-reference and moving-frame energy identities}
\label{M:energy}
For any fixed reference point $x_*$ define $r_k=x_k-x_*$, $a_k=g_k^Tr_k$, $E_k=r_k^TB_kr_k$ and
\[
 R_k=\frac{(y_k^Tr_{k+1})^2}{t_k}.
\]
Expansion of \eqref{M:hessian} at $r_{k+1}=r_k-\alpha_kH_kg_k$ gives exactly
\begin{equation}\label{M:fixedenergy}
 E_{k+1}=E_k-\frac{a_k^2}{q_k}+R_k.
\end{equation}
The cancellation is $-2\alpha a+\alpha^2q-(a-\alpha q)^2/q=-a^2/q$. If $x_*$ is a minimizer, convexity gives $a_k\ge F(x_k)-F_*$. Therefore
\begin{equation}\label{M:energybound}
 [F(x_N)-F_*]^2\sum_{k<N}q_k^{-1}\le E_0+\sum_{k<N}R_k.
\end{equation}
Thus $\sum_{k<N}R_k=o(\sum_{k<N}q_k^{-1})$ excludes a positive gap. With the summability assumptions of Theorem~\ref{M:superquadratic}, an $O(N^2)$ cumulative correction suffices. This last observation is an explicit consolidation deduction from the recorded energy and reciprocal results; no new regularity hypothesis is inferred.

For $V_k\in\mathbb R^{n\times r}$, suppose
\[
 V_{k+1}=V_kT_k+s_kb_k^T.
\]
Set $A_k=V_k^TB_kV_k$, $a_k=V_k^Tg_k$, $c_k=V_{k+1}^Ty_k$. By \eqref{M:annihilation},
\begin{equation}\label{M:frameidentity}
 A_{k+1}=T_k^T\left(A_k-\frac{a_ka_k^T}{q_k}\right)T_k+\frac{c_kc_k^T}{t_k}.
\end{equation}
Suppose on a hypothetical nonoptimal tail $\|a_k\|\ge\epsilon>0$, $\|c_k\|^2/t_k\le C$ and the cumulative products $P_J=I$, $P_{k+1}=P_kT_k$ have bounded norms and bounded inverse norms. Congruence by $P_{k+1}^{-1}$ and a trace yield
\[
 \operatorname{tr}\widetilde A_{k+1}=\operatorname{tr}\widetilde A_k
 -\frac{\|P_k^{-T}a_k\|^2}{q_k}
 +\frac{\|P_{k+1}^{-T}c_k\|^2}{t_k},\quad
 \widetilde A_k=P_k^{-T}A_kP_k^{-1}.
\]
Positivity gives $\sum_{J\le k<N}q_k^{-1}=O(N)$, contradicting Theorem~\ref{M:average}. This is a trajectory criterion; one still has to derive its hypotheses from the objective. Neither $P_k$ nor $T_k$ is a BFGS matrix.

\section{Imperfect transport with a derived smooth-subspace budget}
\label{M:residual}
\begin{theorem}[Residuals in fixed smooth directions]\label{M:residualtheorem}
Assume bounded gradients, fixed weak Wolfe and $\sum d_k<\infty$. For a fixed subspace $S$, suppose
\[
 \|P_Sy_k\|^2\le L_St_k.
\]
A tail cannot simultaneously admit frames with
\[
 \|V_k^Tg_k\|^2\ge m>0,\qquad\|V_{k+1}^Ty_k\|^2\le C_0t_k,
\]
\[
 V_{k+1}=V_kT_k+s_kb_k^T+R_k,
\]
\[
 \|T_k\|_{\rm op}^2\le1+C_Td_k,\quad\operatorname{range}R_k\subset S,
 \quad\|R_k\|_F^2\le C_Rd_k.
\]
No full-rank frame or invertible transport is required.
\end{theorem}
\begin{proof}
Taking the trace of \eqref{M:hessian} on $S$ gives
\[
 \operatorname{tr}(P_SB_{k+1})
 =\operatorname{tr}(P_SB_k)-\frac{\|P_Sg_k\|^2}{q_k}
 +\frac{\|P_Sy_k\|^2}{t_k}
 \le\operatorname{tr}(P_SB_k)+L_S.
\]
Consequently $\operatorname{tr}(P_SB_k)\le C_S+L_Sk$. This bound is derived, not imposed.

Let $E_k=\operatorname{tr}(V_k^TB_kV_k)$ and use $M_k$ from \eqref{M:annihilation}. Then
\[
 E_{k+1}=\|M_k^{1/2}(V_kT_k+R_k)\|_F^2+\frac{\|V_{k+1}^Ty_k\|^2}{t_k}.
\]
Young's inequality with parameter $d_k$ bounds this by
\[
 (1+d_k)\|T_k\|^2\left(E_k-\frac{\|V_k^Tg_k\|^2}{q_k}\right)
 +(1+d_k^{-1})\operatorname{tr}(R_k^TM_kR_k)+C_0.
\]
The parenthesized energy is nonnegative, and
\[
 \operatorname{tr}(R_k^TM_kR_k)
 \le\|R_k\|_F^2\operatorname{tr}(P_SB_k)
 \le C_Rd_k(C_S+L_Sk).
\]
On a tail with $d_k\le1$,
\[
 E_{k+1}\le(1+C_1d_k)(E_k-m/q_k)+C_2(k+1).
\]
Normalize by $P_{k+1}=P_k(1+C_1d_k)$ with $P_J=1$. Finite $\sum d_k$ gives $1\le P_k\le\bar P<\infty$. Summing and using final-energy positivity gives
\[
 \frac m{\bar P}\sum_{k=J}^{N-1}q_k^{-1}
 \le E_J+C_2\sum_{k=J}^{N-1}(k+1)=O(N^2),
\]
contrary to Theorem~\ref{M:superquadratic}.
\end{proof}
Only residuals in the fixed subspace are controlled. Euclidean smallness alone does not control a residual in a direction where $B_k$ may grow without this estimate. No strict inclusion of all existential objective-certificate classes is claimed.

% END original chapters/foundation.tex

\StartModule{T}{Objective-level transported certificates}{Selected proof from source [ST]. Positive bands and small eligible pairs, not assumptions on BFGS matrices. Normalization is proved in Module G.}

% BEGIN original chapters/T.tex
\begingroup
\providecommand{\R}{}
\renewcommand{\R}{\mathbb R}
\providecommand{\tr}{}
\renewcommand{\tr}{\operatorname{tr}}
\providecommand{\diag}{}
\renewcommand{\diag}{\operatorname{diag}}
\providecommand{\Span}{}
\renewcommand{\Span}{\operatorname{span}}
\providecommand{\ip}{}
\renewcommand{\ip}[2]{\langle #1,#2\rangle}
\providecommand{\norm}{}
\renewcommand{\norm}[1]{\lVert #1\rVert}
\providecommand{\T}{}
\renewcommand{\T}{^{\mathsf T}}
\providecommand{\D}{}
\renewcommand{\D}{\mathcal D}
\setcounter{section}{1}
\setcounter{theorem}{0}
\setcounter{lemma}{0}
\setcounter{corollary}{0}
\setcounter{proposition}{0}
\setcounter{remark}{0}
\setcounter{equation}{2}
\section{Algorithm and objective-pair terminology}
Let $F:\R^n\to\R$ be finite and convex, let $F_*=\min F$ be attained, and let
\[
 K=\{x:F(x)\le b\},\qquad b>F_*,
\]
be compact. Ordinary full BFGS uses differentiable accepted endpoints,
\[
 g_k=\nabla F(x_k),\quad s_k=-\alpha_kH_kg_k,\quad
 x_{k+1}=x_k+s_k,\quad y_k=g_{k+1}-g_k,
\]
with $H_0\succ0$, $\alpha_k>0$, and
\[
 H_{k+1}=(I-s_ky_k\T/t_k)H_k(I-y_ks_k\T/t_k)+s_ks_k\T/t_k,
 \qquad t_k=y_k\T s_k.
\]
The fixed Armijo--weak-Wolfe conditions are
\[
 F(x_{k+1})\le F(x_k)+c_1g_k\T s_k,\qquad
 g_{k+1}\T s_k\ge c_2g_k\T s_k,\qquad 0<c_1<c_2<1.
\]
An infinite admissible run has nonzero gradients and well-defined accepted steps. These are the algorithmic quantifiers in Xie--W\"achter~\cite{XW}. No particular line-search implementation, damping, restart, or proximal operation is added.

For differentiable endpoints $x,z$, write
\[
 g=\nabla F(x),\quad y=\nabla F(z)-g,\quad s=z-x,
\]
\[
 d=-g\T s,\qquad \Delta=F(x)-F(z),\qquad t=y\T s.
\]
Call the pair \emph{eligible} if
\begin{equation}
 d>0,\qquad c_1d\le\Delta\le d,\qquad t\ge(1-c_2)d. \tag{T.2}
\end{equation}
Every accepted step is eligible. The upper bound $\Delta\le d$ follows from convexity. Eligibility does not assert compatibility of an arbitrary chain with BFGS; imposing a condition on all such pairs is an objective-only sufficient assumption.

For $s\ne0$, define the Euclidean orthogonal projector
\[
 \Pi_s=I-\frac{ss\T}{s\T s}.
\]
All matrices in the next assumption are objective-pair data, not $H_k$.


\setcounter{section}{2}
\setcounter{theorem}{0}
\setcounter{lemma}{0}
\setcounter{corollary}{0}
\setcounter{proposition}{0}
\setcounter{remark}{0}
\setcounter{equation}{3}
\section{The transported quadratic secant condition}
\begin{theorem}[One-sided transported-certificate convergence]
For each $\eta>0$ with nonempty band
\[
 K_\eta=\{x\in K:F(x)\ge F_*+\eta\},
\]
suppose there is a finite-valued symmetric positive-semidefinite field $W_\eta(x)$ on differentiable points of $K_\eta$, and constants
\[
 m_\eta>0,\quad C_\eta<\infty,\quad L_\eta\ge0,\quad \delta_\eta>0,
\]
with these properties:
\begin{align}
 g\T W_\eta(x)g&\ge m_\eta, \tag{T.V}\\
 y\T W_\eta(z)y&\le C_\eta t, \tag{T.S}\\
 \Pi_s W_\eta(z)\Pi_s&\preceq
 (1+L_\eta d)\Pi_s W_\eta(x)\Pi_s. \tag{T.T}
\end{align}
The last two inequalities need hold only for eligible pairs $x,z\in K_\eta$ with $d\le\delta_\eta$. Then every infinite admissible full-BFGS run starting in $K$ satisfies
\[
 F(x_k)\to F_*,\qquad \operatorname{dist}(x_k,\arg\min F)\to0.
\]
If the minimizer is unique, the entire iterate sequence converges to it.
\end{theorem}
Condition (V) says that the comparison directions retain a nonzero component of the gradient on a positive band. Condition (S) controls the secant correction in those directions. Condition (T) controls their expansion only in directions perpendicular to the proposed step.

There is no uniform positive-definiteness requirement on $W_\eta$. Uniform boundedness of $W_\eta$ is convenient in examples but is not used by the proof; it is enough that its values are finite. Constants and the field may depend on the positive band and the fixed line-search parameters. They may degenerate as $\eta\downarrow0$.

\subsection{Proof}
Put
\[
 B_k=H_k^{-1},\qquad q_k=g_k\T H_kg_k,\qquad
 d_k=\alpha_kq_k.
\]
The ordinary Hessian update is
\begin{equation}
 B_{k+1}=B_k-\frac{g_kg_k\T}{q_k}+\frac{y_ky_k\T}{t_k}. \tag{T.3}
\end{equation}
Its first two terms form a PSD matrix
\[
 M_k=B_k-g_kg_k\T/q_k\succeq0,
\]
by metric Cauchy--Schwarz. Also
\begin{equation}
 M_ks_k=0,\qquad M_k=\Pi_{s_k}M_k\Pi_{s_k}. \tag{T.4}
\end{equation}
This is why (T) may ignore the step direction.

Armijo implies $\sum d_k<\infty$. Gradients are bounded on $K$: a finite convex function has uniformly bounded subgradients on a compact set, as follows by comparing its values at $x$ and $x+\epsilon u$ on a slightly enlarged compact set. Thus the shared foundation in Module M applies and gives
\begin{equation}
 \frac1N\sum_{k<N}\frac1{q_k}\longrightarrow\infty. \tag{T.5}
\end{equation}
Suppose the objective gap tends to $\ell>0$. Use the field for $\eta=\ell/2$. On a tail $k\ge k_0$, all pairs lie in the band and $d_k\le\delta_\eta$. Suppress $\eta$ and set
\[
 E_k=\tr(B_kW(x_k))\ge0,\qquad a_k=1+L d_k.
\]
Using (3)--(4), the Loewner inequality (T), and (S),
\begin{align}
 E_{k+1}
 &=\tr(M_kW(x_{k+1}))+y_k\T W(x_{k+1})y_k/t_k\notag\\
 &\le a_k\tr(M_kW(x_k))+C\notag\\
 &=a_k\left(E_k-\frac{g_k\T W(x_k)g_k}{q_k}\right)+C\notag\\
 &\le a_k(E_k-m/q_k)+C. \tag{T.6}
\end{align}
Let $P_{k_0}=1$, $P_{k+1}=P_ka_k$. Since $\sum d_k<\infty$,
\[
 1\le P_k\le \exp\left(L\sum_{j\ge k_0}d_j\right)=:\overline P<\infty.
\]
Divide (6) by $P_{k+1}$ and sum:
\[
 \frac m{\overline P}\sum_{k=k_0}^{N-1}\frac1{q_k}
 \le E_{k_0}+C(N-k_0).
\]
This contradicts (5). The gap therefore tends to zero. Compactness gives convergence to the solution set, and uniqueness gives whole-sequence convergence. \qed

\subsection{Local statement and weaker band requirements}
If the assumptions hold only in a neighborhood $U$ of a unique minimizer and $K$ is a compact low sublevel, choose an initial level below
\[
 \inf_{x\in K\setminus U}F(x)>F_*.
\]
Armijo then traps all accepted iterates in $U$. This produces a local basin independent of $H_0$. Rejected trial points need not stay in $U$. For multiple minimizers, one may instead use a neighborhood of the entire compact solution set or directly assume that the initial sublevel lies in the structural region.

The proof only needs a certificate in a small one-sided level neighborhood of each possible finite limit $\ell>F_*$. Requiring a field on every whole positive band is a cleaner, stronger formulation; certificates on shells $\ell\le F\le\ell+\epsilon_\ell$ also suffice.

\subsection{Scalar form and freedom to rescale}
Taking $W(x)=v(x)v(x)\T$, a sufficient version of (T) is
\[
 v(z)=\lambda(x,z)v(x)+\beta(x,z)s,\qquad
 |\lambda(x,z)|^2\le1+L d.
\]
Conditions (V)--(S) become
\[
 |g\T v(x)|^2\ge m,\qquad |y\T v(z)|^2\le Ct.
\]
Unlike the preceding moving-frame theorem, there is no requirement that $\lambda$ be invertible or bounded away from zero. The PSD form also accommodates multiple comparison directions and varying rank, without choosing an exact frame transport.

A bounded potential may be absorbed into $W$. If (T) instead has multiplier
\[
 \exp\{L d+\Psi(x)-\Psi(z)\},
\]
where $\Psi$ is bounded on the band, replace $W$ by $e^{\Psi}W$. This gives multiplier $e^{Ld}$, which is at most $1+L'd$ for $d$ sufficiently small. The other bounds change only by band-dependent constants.


\setcounter{section}{3}
\setcounter{theorem}{1}
\setcounter{lemma}{0}
\setcounter{corollary}{0}
\setcounter{proposition}{0}
\setcounter{remark}{0}
\setcounter{equation}{6}
\section{Why this is genuinely weaker than RS}
Given (RS), take
\[
 W(x)=(x-x_*)(x-x_*)\T.
\]
Convexity gives
\[
 g\T W(x)g=(g\T(x-x_*))^2\ge(F(x)-F_*)^2\ge\eta^2.
\]
Condition (S) is exactly (RS), while
\[
 \Pi_s(z-x_*)=\Pi_s(x-x_*)
\]
makes (T) an equality with $L=0$. Thus every objective covered by RS is covered by the new condition.

The reverse implication fails. Module Q constructs fields satisfying the new condition for a multi-branch curved objective that violates even the eligible-pair radial condition on every positive neighborhood of its minimizer. This is strict inclusion of sufficient conditions, not merely a new notation for the old inequality.

Although $\Pi_s$ is written as a Euclidean projector, (T) equivalently requires
\[
 a\T W(z)a\le(1+Ld)a\T W(x)a\qquad\text{for every }a\text{ with }a\T s=0.
\]
Under an invertible affine change of variables with linear part $A$, gradients transform by $A^{-\mathsf T}$ and the certificate transforms by $W\mapsto AWA\T$. The three conditions are therefore affine-covariant; no preferred coordinate system is essential.

The extra transport clause cannot simply be omitted: changing a comparison vector creates terms involving $B_k$, which may be unbounded. A small Euclidean change is not by itself a small change in the $B_k$ energy. Projection and the multiplicative budget are precisely what make the new hypothesis sufficient without assumptions on $B_k$.


\endgroup

% END original chapters/T.tex

\StartModule{R}{Radial smoothness and finite partitions}{Selected complete arguments from source [SR]. RS is a special case, not the most general sufficient condition. Boundary examples disprove a criterion, not convergence.}

% BEGIN original chapters/R.tex
\begingroup
\providecommand{\R}{}
\renewcommand{\R}{\mathbb R}
\providecommand{\ip}{}
\renewcommand{\ip}[2]{\langle #1,#2\rangle}
\providecommand{\DF}{}
\renewcommand{\DF}{D_F}
\providecommand{\st}{}
\renewcommand{\st}{x_*}
\providecommand{\fs}{}
\renewcommand{\fs}{F_*}
\makeatletter\let\c@lemma\c@theorem\let\c@corollary\c@theorem\let\c@proposition\c@theorem\let\c@remark\c@theorem\makeatother
\renewcommand{\thelemma}{\thetheorem}\renewcommand{\thecorollary}{\thetheorem}\renewcommand{\theproposition}{\thetheorem}\renewcommand{\theremark}{\thetheorem}
\renewcommand{\thetheorem}{R.\arabic{section}.\arabic{theorem}}
\setcounter{section}{0}
\setcounter{theorem}{0}
\setcounter{equation}{0}
\section{Setup: the existing condition}
Let $F:\R^n\to\R$ be finite and convex, with a minimizer $\st$ and minimum value $\fs$. Let
\[
 K=\{x:F(x)\le b\},\qquad b>\fs,
\]
be compact. Write $r_x=x-\st$. Fix bounds $R>0$ and $G>0$ such that
\[
 \|r_x\|\le R,\qquad \|g\|\le G\quad(x\in K,\ g\in\partial F(x)).
\]
Such a subgradient bound exists because a finite convex function is locally Lipschitz on the whole space.

For differentiable $x,z$, put
\begin{align*}
 g&=\nabla F(x),& h&=\nabla F(z),&y&=h-g,&s&=z-x,\\
 T(x,z)&=y^Ts,&N(x,z)&=y^Tr_z.
\end{align*}
The supporting-plane error, or Bregman expression, is
\[
 \DF(z,x)=F(z)-F(x)-g^T(z-x)\ge0.
\]
Unlike the strictly convex smooth setting of many standard presentations \cite{BB}, only differentiability of the second endpoint is needed to define this expression here. We have
\[
 T=\DF(z,x)+\DF(x,z),\qquad 0\le T\le T_{\max}:=4GR.
\]
The preceding report assumes the following positive-band condition:
\begin{equation}
 \forall\eta>0\ \exists C_\eta<\infty:\quad
 N(x,z)^2\le C_\eta T(x,z)
 \quad\text{if }\fs+\eta\le F(z)\le F(x),\ x,z\in K,\ T>0.
 \tag{R.RS}\label{R:eq:RS}
\end{equation}
Endpoints in this condition are differentiable. Constants may diverge as $\eta\downarrow0$. The hypothesis is imposed only on nonempty bands.

The algorithm remains ordinary full BFGS with arbitrary $H_0\succ0$, differentiable accepted iterates, and positive steps satisfying fixed Armijo--weak-Wolfe parameters $0<c_1<c_2<1$, as in \cite{XW}. ``Every run'' below always means every infinite admissible run; a zero gradient is successful termination. General all-ray existence of differentiable acceptable points is not asserted.


\setcounter{section}{1}
\setcounter{theorem}{0}
\setcounter{equation}{1}
\section{Equivalent formulation: smooth only toward the minimizer}
\begin{theorem}[Inward radial smoothness is equivalent to (RS)]\label{R:thm:equiv}
On the compact sublevel $K$, condition (RS) is equivalent, up to constants, to the following hypothesis. For every $\eta>0$, there exist $M_\eta<\infty$ and $\theta_\eta\in(0,1/2]$ such that, at every differentiable $x\in K$ with $F(x)\ge\fs+\eta$,
\begin{equation}
 0\le F(x-t r_x)-F(x)+t\nabla F(x)^Tr_x
 \le \frac{M_\eta}{2}t^2,
 \qquad 0\le t\le\theta_\eta.
 \tag{R.IR}\label{R:eq:IR}
\end{equation}
No differentiability of $x-t r_x$ is required. The sampled point is a convex combination of $x$ and $\st$, so it remains in $K$.
\end{theorem}
The label ``inward radial smoothness'' is descriptive terminology in this note, not a claim of a standardized definition or of novelty for the terminology. The parameter $t$ is a fraction of the distance to the minimizer. A factor $\|r_x\|^2$ can alternatively be placed in the upper bound: on a positive objective band, the two conventions are equivalent because $\eta/G\le\|r_x\|\le R$.

\subsection{Proof that (IR) implies (RS)}
Fix a positive band and write $M,\theta$ for its constants. Consider endpoints $x,z$ on that band. If $N\ge0$, evaluate the supporting plane at $x$ at the inward test point $z-t r_z$. Convexity and (IR) at $z$ give
\[
 tN\le \DF(z,x)+\frac M2t^2,\qquad 0\le t\le\theta.
\]
If $N<0$, use the inward test point $x-t r_x$ and the supporting plane at $z$. Since
\[
 (g-h)^Tr_x=T-N\ge|N|,
\]
we instead obtain
\[
 t(T-N)\le\DF(x,z)+\frac M2t^2.
\]
In either case there is a number $A\in[|N|,2GR]$ and a supporting error $D\in\{\DF(z,x),\DF(x,z)\}$ such that
\[
 tA\le D+\frac M2t^2\quad(0\le t\le\theta).
\]
When $M>0$ and $A\le M\theta$, take $t=A/M$ to get $A^2\le2MD$. Otherwise take $t=\theta$; if $A>M\theta$, then $D\ge A\theta/2$, so
\[
 A^2\le \frac{4GR}{\theta}D.
\]
For $M=0$, choosing $t=\theta$ gives an even stronger bound. Thus
\begin{equation}
 N^2\le C T,\qquad
 C=\max\left\{2M,\frac{4GR}{\theta}\right\}.
 \label{R:eq:IRtoRS}
\end{equation}
This proof also shows why controlling only the direction \emph{toward} the minimizer is enough: the sign of $N$ determines which endpoint supplies the inward probe. No outward point is needed.

\subsection{Proof that (RS) implies (IR)}
Fix $\eta>0$ and choose
\[
 \theta_\eta=\min\left\{\frac12,\frac{\eta}{4GR}\right\}.
\]
For $x$ on the $\eta$-band and $z=x-t r_x$, the Lipschitz bound on $K$ gives
\[
 F(z)\ge F(x)-GtR\ge\fs+3\eta/4.
\]
First suppose both endpoints are differentiable. Put
\[
 a=(\nabla F(x)-\nabla F(z))^Tr_x\ge0.
\]
Then
\[
 T(x,z)=ta,\qquad N(x,z)=-(1-t)a.
\]
Using (RS) on the $\eta/2$-band, if $a>0$,
\[
 (1-t)^2a^2\le C_{\eta/2}ta,
 \qquad
 a\le\frac{C_{\eta/2}t}{(1-t)^2}.
\]
Convexity supplies
\[
 0\le\DF(z,x)\le ta\le4C_{\eta/2}t^2.
\]
If $a=0$, the same conclusion follows from $\DF(z,x)\le ta=0$. Thus (IR) holds with $M_\eta=8C_{\eta/2}$.

For completeness, the inner endpoint may be nondifferentiable. For fixed $t<1$, approximate $x$ by points $x_j$ in the interior of $K$ such that both $x_j$ and $x_j-t(x_j-\st)$ are differentiable. This is possible because the differentiability set has full measure and the affine dilation is invertible. The positive-band inequalities persist, using the spare $\eta/4$ margin above. Convexity and differentiability at $x$ imply $\nabla F(x_j)\to\nabla F(x)$. Passing to the limit in the function-value inequality proves (IR) without differentiability of the inner endpoint. Boundary points of $K$ are approximated from its interior. The case $t=0$ is trivial. \qed

\subsection{Interpretation and a derivative-free version}
For a unit direction $d$, inspect only the scalar function
\[
 \varphi_d(\rho)=F(\st+\rho d),\qquad \rho>0.
\]
On compact positive objective bands, the equivalent conditions control the change of its scalar derivative uniformly over directions. Angular derivatives may jump. In a smooth region, the infinitesimal quantity being bounded is
\[
 r_x^T\nabla^2F(x)r_x,
\]
not the largest Hessian eigenvalue. For a nonsmooth function, checking this expression almost everywhere is \emph{not enough}: a radial gradient jump contributes singular curvature and must also be excluded.

On slightly nested positive bands where both dilations remain in $K$, an equivalent function-value formulation is a uniform bound
\begin{equation}
 F(\st+(1+t)r)+F(\st+(1-t)r)-2F(\st+r)\le M_\eta t^2.
 \label{R:eq:central}
\end{equation}
Indeed, this bound directly bounds each nonnegative tangent-plane remainder at a differentiable middle point, yielding (IR). Conversely, apply (RS) to the outer/middle and middle/inner radial pairs. If $a_+$ and $a_-$ are their scalar derivative differences in direction $r$, then
\[
 a_+\le C t,\qquad a_-\le Ct/(1-t)^2.
\]
The central difference is at most $t(a_++a_-)\le5Ct^2$ for $t\le1/2$. Approximation by differentiable triples extends the function-value inequality to all centers. This is a local equivalence with constants allowed to change; an outward point need not remain in the original top boundary of $K$.


\setcounter{section}{2}
\setcounter{theorem}{0}
\setcounter{equation}{4}
\section{A second interpretation: what tangent planes predict at the optimum}
Define, at a differentiable point,
\[
 b_F(x)=\fs-\bigl[F(x)+\nabla F(x)^T(\st-x)\bigr]
       =\nabla F(x)^Tr_x-(F(x)-\fs)\ge0.
\]
This is the amount by which the supporting plane at $x$ underestimates the minimum when extrapolated to $\st$. It is also the supporting error $\DF(\st,x)$.

A direct calculation gives
\begin{equation}
 N(x,z)=\DF(z,x)+b_F(z)-b_F(x).
 \label{R:eq:intercept}
\end{equation}
Since $0\le\DF(z,x)\le T\le T_{\max}$, the radial condition is equivalent on the same bounded domains to
\begin{equation}
 |b_F(z)-b_F(x)|^2\le C' T(x,z).
 \tag{R.SP}\label{R:eq:SP}
\end{equation}
For example, (RS) with $C$ implies (SP) with $2C+2T_{\max}$, and conversely. Thus the full gradient need not be Lipschitz; only the \emph{height at the optimizer} of its supporting plane needs this curvature-relative regularity. The square root of $T$ is not being asserted to be a metric.

For a degree-one homogeneous function minimized at zero, $b_F\equiv0$. For
\[
 F(u,v)=u^2+|v|,
\]
we have $b_F(u,v)=u^2$: crossing $v=0$ flips a gradient component but does not produce a jump in this scalar. For a nonnegative degree-$p$ homogeneous function, $b_F=(p-1)F$. These identities explain much of the breadth of the previous sufficient classes.

In the four-piece counterexample, the optimizer is $(-1,5)$, the minimum is $-1$, and the first three branches have values $-1,-1,-29/4$ there. Their supporting errors at the optimizer are therefore
\[
 0,\quad0,\quad25/4.
\]
The bad trajectory repeatedly switches between these incompatible scalar predictions while $T$ tends to zero. This is exactly what (SP) prohibits. The supporting-plane calculations here do not change or independently re-prove the earlier infinite BFGS certificate.


\setcounter{section}{4}
\setcounter{theorem}{0}
\setcounter{equation}{8}
\section{A broader verifiable family: polyhedrally piecewise smooth convex functions}
\begin{theorem}[Finite polyhedral partitions]\label{R:thm:plq}
Let $F$ be finite and convex, with a unique minimizer $\st$. Suppose a neighborhood of $\st$ is covered by finitely many closed full-dimensional polyhedra $P_i$, with compatible function values on overlaps, and
\[
 F(x)=f_i(x)\quad(x\in P_i)
\]
there. Assume each $f_i$ has a $C^{1,1}$ extension to an open neighborhood of the relevant part of its cell. Then every sufficiently low-starting infinite admissible full-BFGS run converges to $\st$, for arbitrary SPD initialization.
\end{theorem}
\begin{proof}
Discard all cells not containing $\st$ by shrinking to a ball whose radius is below their positive distances from $\st$. There are only finitely many cells. If a remaining cell contains $x$, it also contains the segment from $x$ to $\st$ by convexity. At a differentiable point of $F$, its gradient agrees with the smooth gradient of each adjacent full-dimensional cell: approach from the interior of that cell and use differentiability and convexity of $F$.

Let $L$ bound the Lipschitz constants of the finitely many smooth extensions on a smaller ball of radius $R$. For $x$ in one such cell, the usual Taylor estimate along the inward segment gives
\[
 0\le F(x-t r_x)-F(x)+t\nabla F(x)^Tr_x
 \le\frac L2 t^2\|r_x\|^2\le\frac{LR^2}{2}t^2.
\]
The lower bound follows from convexity of $F$; individual smooth extensions need not be convex outside their cells. Thus uniform (IR), hence (RS), holds locally. A low compact sublevel is trapped there, and the preceding convergence theorem applies. No cell or active-manifold identification is assumed of the iterates.
\end{proof}
In particular this covers every finite convex \emph{piecewise linear-quadratic} objective, and a convex smooth loss plus a finite convex piecewise linear-quadratic regularizer. This is broader than merely having affine pieces plus one common smooth summand: different cells can have different quadratic Hessians and different gradient jumps. Convexity of the full objective remains essential. A finite maximum of arbitrary convex quadratics need not have a polyhedral partition and is not automatically included.


\setcounter{section}{5}
\setcounter{theorem}{0}
\setcounter{equation}{8}
\section{Common families through the one-dimensional profile}
All conclusions in this section require the indicated minimizer and compact-basin hypotheses. Local statements are not global arbitrary-initialization statements.

\subsection{Norms, gauges, and mixed powers}
For a finite sublinear $\gamma$ centered at zero,
\[
 \gamma((1-t)z)=(1-t)\gamma(z),\qquad 0\le t<1.
\]
There is exactly zero radial Taylor error, regardless of angular nonsmoothness. This includes polyhedral norms, Euclidean norms, the spectral norm and nuclear norm when centered at zero, and asymmetric finite gauges.

For nonnegative $\gamma$ and $p\ge1$,
\[
 \gamma((1-t)z)^p=(1-t)^p\gamma(z)^p.
\]
On $0\le t\le1/2$, the second derivative in $t$ is bounded uniformly on any bounded set of $z$. This includes $1<p<2$, without demanding a Lipschitz full gradient near zero. Therefore finite sums of differently powered, linearly composed gauges,
\[
 F(\st+z)=\psi(z)+\sum_j\gamma_j(A_jz)^{p_j},\qquad p_j\ge1,
\]
are covered when $\psi$ is convex $C^{1,1}$ locally and the full objective has the required minimizer and compact sublevel. The individual terms need not be coercive.

\subsection{Polyhedral losses and penalties at nonzero solutions}
Locally around any optimizer, a nonzero scalar affine residual has fixed sign and is smooth; a zero residual becomes a centered absolute value or hinge. Hence smooth losses plus $\ell_1$, fused/generalized $\ell_1$, finite hinge and quantile losses are covered locally under uniqueness. The minimizer is not required to be zero. Finite polyhedral regularizers, and their piecewise linear-quadratic extensions from Theorem~\ref{R:thm:plq}, are also covered. Extended-valued constraints remain outside the finite-function formulation.

\subsection{Affine Euclidean group residuals}
For terms $\|A_jx-b_j\|_2$, a nonzero optimal residual is locally smooth, while a zero optimal residual becomes $\|A_jz\|_2$, which is homogeneous. Thus sums of such terms, including overlapping Euclidean group penalties and sums of distances, added to a suitable convex smooth objective, satisfy the local condition when the minimizer is unique. This does not imply that every spectral penalty is smooth near every nonzero residual.

\subsection{Finite maxima: a simple branch test}
If locally $F=\max_i f_i$ with finitely many convex $C^2$ branches, it suffices that
\[
 (\nabla f_i(w)-\nabla f_j(w))^T(w-\st)=0
 \quad(i,j\in I(w))
\]
at every active tie near $\st$. Along a fixed ray, active branches then agree in the scalar first derivative. That derivative is a continuous selection of finitely many Lipschitz branch derivatives, and has their common Lipschitz bound on compact intervals. This proves (IR) directly, shortening the earlier segment-based proof.

The necessity for the old radial condition is also clear on genuine active ties: if the active radial derivatives differ by $J>0$, the one-dimensional central second difference at the tie is $J|t|+O(t^2)$, not $O(t^2)$. Thus radial compatibility describes the exact no-jump part of the criterion for this finite smooth-branch class. It is not a necessary condition for BFGS convergence itself.


\setcounter{section}{6}
\setcounter{theorem}{0}
\setcounter{equation}{8}
\section{Important boundaries: curved ridges and low-rank spectral penalties}
\subsection{The parabolic ridge}
For
\[
 F(u,v)=\max\{u^2+v,2u^2-v\},\qquad \st=0,
\]
the switching curve is $v=u^2/2$. At $w=(u,u^2/2)$, $u\ne0$, the active gradient difference is $(-2u,2)$ and its radial inner product is $-u^2\ne0$. A ray crosses a genuine one-dimensional kink there. For the descending pair
\[
 x_\epsilon=(u,u^2/2-\epsilon),\qquad
 z_\epsilon=(u,u^2/2+\epsilon/2),
\]
the ratio is exactly
\[
 \frac{(y^Tz_\epsilon)^2}{y^T(z_\epsilon-x_\epsilon)}
 =\frac{(u^2-\epsilon)^2}{3\epsilon}\to\infty.
\]
This fails (RS) on a positive objective band in every neighborhood of zero. It is \emph{not} a nonconvergent BFGS trajectory and does not establish failure of BFGS.

\subsection{A common spectral regularizer not automatically covered}
Work in symmetric $2\times2$ matrices with the Frobenius inner product. Let
\[
 F(X)=\frac12\|X-C\|_F^2+\|X\|_*,\qquad
 C=\operatorname{diag}(2,0),\quad X_*=\operatorname{diag}(1,0).
\]
This is strongly convex and uniquely minimized at $X_*$. Indeed $X_*-C=\operatorname{diag}(-1,0)$ is canceled by a nuclear-norm subgradient $\operatorname{diag}(1,0)$ there.

Take orthonormal vectors
\[
 v=(\cos\theta,\sin\theta),\qquad u=(-\sin\theta,\cos\theta),
\]
and let $W=vv^T$, $U=uu^T$, $a=\sin^2\theta>0$. For
\[
 X_\epsilon=W-\epsilon U,\qquad Z_\epsilon=W+\epsilon U,
\]
both nuclear norms are differentiable, since the matrices have nonzero eigenvalues. The pair is descending: $F(Z_\epsilon)-F(X_\epsilon)=-4a\epsilon<0$. Its gradient difference is $y=2(1+\epsilon)U$, so
\[
 T=4\epsilon(1+\epsilon),\qquad
 N=2(1+\epsilon)(\epsilon-a),\qquad
 \frac{N^2}{T}=\frac{(1+\epsilon)(\epsilon-a)^2}{\epsilon}\to\infty.
\]
Also $F(W)-F(X_*)=2a>0$. Choosing small $\theta$ puts such points arbitrarily close to $X_*$. Thus this familiar smooth-plus-nuclear-norm problem need not satisfy the original local radial criterion at a nonzero rank-deficient minimizer. The nuclear norm centered at zero \emph{is} covered. No failure of BFGS on this spectral problem is being claimed.


\endgroup

% END original chapters/R.tex

\StartModule{F}{Structural corollaries and rates}{Consolidated proofs from [SM], [SP], [SE] and [SQ]. Includes centered composites, transformed gauges, explicit basins, lifting, conservative rates and line-search qualifications.}

% BEGIN original chapters/families.tex
\section{Centered smooth-plus-sublinear objectives}
This module collects the nonredundant structural consequences from [M], [P], [E] and [Q]. The finite-polyhedral-partition and powered-gauge radial arguments themselves are reproduced in Module R. The discussion here retains the global, nonunique, explicit-basin and quantitative clauses that are not supplied by a bare local uniqueness theorem.

\begin{theorem}[Centered smooth-plus-sublinear convergence]\label{F:centered}
Let $F(x)=\psi(x)+\gamma(x-x_*)$, where $\psi$ is convex and locally $C^{1,1}$ on an open convex neighborhood of a compact initial sublevel, and $\gamma$ is finite convex and positively homogeneous of degree one. Suppose $x_*$ minimizes $F$. Every infinite admissible full-BFGS run has $F(x_k)\to F(x_*)$. If $x_*$ is unique, then $x_k\to x_*$. A representation only near a unique minimizer gives a low-starting local theorem by Lemma~\ref{M:localization}.
\end{theorem}
\begin{proof}
Translate $x_*=0$ and let $R$ bound the initial sublevel. At differentiable endpoints, $\gamma$ is differentiable because $F$ and $\psi$ are. Write $y=y_\psi+y_\gamma$, and set
\[
 a=y_\psi^Ts,\quad b=y_\gamma^Ts,\quad
 \xi=y_\psi^Tx',\quad\zeta=y_\gamma^Tx'.
\]
Then $a,b\ge0$ and $t=a+b$. Convex local $L$-smoothness gives $\|y_\psi\|^2\le La$, hence $\xi^2\le LR^2a$. This cocoercivity statement is valid on an open convex neighborhood; the source cites Wachsmuth--Wachsmuth for the local version.

If $h\in\partial\gamma(x)$, homogeneity and the supporting inequality give $h^Tx=\gamma(x)$ and $h^Tz\le\gamma(z)$. Hence $0\le\zeta\le b$. All subgradients of $\gamma$ are bounded by some $G_\gamma$, so $\zeta\le2G_\gamma R$. Weighted Cauchy--Schwarz gives
\[
 \frac{(y^Tx')^2}{t}
 \le\frac{\xi^2}{a}+\frac{\zeta^2}{b}
 \le LR^2+2G_\gamma R.
\]
Terms with zero denominator are zero since their numerators vanish. Apply the fixed-reference energy bound and the reciprocal-average theorem in Module M. Compactness and uniqueness give the iterate conclusion.
\end{proof}
The original summands need not separately be minimized at the center. Optimality implies $-\nabla\psi(x_*)\in\partial\gamma(0)$. After translating, subtract the smooth tangent plane from $\psi$ and add it to $\gamma$; both adjusted summands are then nonnegative and minimized at zero. This is a permissible rewriting of the same objective, not invariance of BFGS under a change of objective.

\section{Smooth increasing transforms of coercive gauges}
\begin{theorem}[Transformed-gauge convergence]\label{F:transform}
Let $\gamma$ be finite convex, degree-one homogeneous and coercive. Let $\Phi:[0,\infty)\to\mathbb R$ be continuous, convex, strictly increasing, with $\Phi(0)=0$, $\Phi\in C^2((0,\infty))$ and $\Phi'>0$ there. Then every infinite admissible full-BFGS run on $F(x)=\Phi(\gamma(x-x_*))$ converges to $x_*$.
\end{theorem}
\begin{proof}
Coercivity of $\gamma$ and unboundedness of $\Phi$ give compact sublevels. Translate $x_*=0$. Suppose $F_k\downarrow L>0$. Euler's identity and the chain rule give
\[
 a_k=g_k^Tx_k=\gamma(x_k)\Phi'(\gamma(x_k))=\Theta(F_k),
\quad\Theta(t)=\Phi^{-1}(t)\Phi'(\Phi^{-1}(t)).
\]
On $[L,F_0]$, the derivative
\[
 \Theta'(t)=1+\frac{r\Phi''(r)}{\Phi'(r)},\qquad r=\Phi^{-1}(t),
\]
is bounded. Thus $|a_{k+1}-a_k|\le C\Delta_k$. Since
\[
 y_k^Tx_{k+1}=d_k+a_{k+1}-a_k,
\]
\eqref{M:basic} implies
\[
 R_k\le\frac{(1+C)^2}{(1-c_2)c_1}\Delta_k.
\]
The corrections are summable, contradicting a positive gap in \eqref{M:energybound}. Hence $F_k\to0$, and coercivity and strict increase give $x_k\to0$.
\end{proof}
This proof never claims invariance of BFGS under a nonlinear scalar transform. It includes cases whose scalar second derivative is unbounded at zero. Finite sums of differently powered gauges are instead covered by the radial argument in Module R, including $1<p<2$.

\section{Finite max-affine terms and a certified basin}
Let $F=\psi+p$, with $\psi$ convex locally $C^{1,1}$ and $p(x)=\max_i\ell_i(x)$, $\ell_i(x)=a_i^Tx+b_i$, a finite real-valued maximum. At a unique minimizer $x_*$, put $I_*=\{i:\ell_i(x_*)=p(x_*)\}$. Inactive pieces have positive gaps, so on a sufficiently small ball
\begin{equation}\label{F:active}
 p(x_*+z)=p(x_*)+\max_{i\in I_*}a_i^Tz.
\end{equation}
Indeed, with any fixed $i_0\in I_*$, the difference between its value and an inactive piece $j$ is
\[
 \delta_j+(a_{i_0}-a_j)^Tz,\qquad \delta_j=p(x_*)-\ell_j(x_*)>0,
\]
which stays positive if $\|z\|\le\delta_j/(2\|a_{i_0}-a_j\|)$; a parallel inactive piece stays dominated everywhere. This exact local identity and Theorem~\ref{F:centered} prove local convergence without strict complementarity or Hessian nondegeneracy.

For an explicit sufficient level, define the closed active region
\[
 P_j=\{x:\ell_j(x)=p(x)\}
 =\{x:(a_j-a_i)^Tx\ge b_i-b_j\ \forall i\}
\]
and
\begin{equation}\label{F:basin}
 \tau_* =\min_{j\notin I_*}\inf_{x\in P_j}F(x),
\end{equation}
with empty infima or minima equal to $+\infty$. Then $\tau_*>F_*$. To prove it, note that unique-minimizer convexity gives compact sublevels; a finite infimum on a nonempty closed $P_j$ is attained. It cannot equal $F_*$ because $x_*\notin P_j$. Finitely many inactive pieces preserve the strict gap.

Every admissible run with $F(x_0)<\tau_*$ converges. On its full compact initial sublevel, only indices in $I_*$ can be active, so \eqref{F:active} and the centered correction estimate apply throughout. A tail that first enters this sublevel enjoys the same conclusion for its current arbitrary SPD matrix. The threshold is sufficient, not a characterization of every successful initialization. Each inner infimum is a convex program over the displayed linear inequalities.

\section{Residual penalties and affine-fiber lifting}
For a scalar affine residual nonzero at the solution, its absolute-value or hinge term is locally affine. A zero optimal residual produces a centered sublinear term. The same decomposition applies to an affine-composed Euclidean norm: a nonzero optimal vector residual is locally smooth, and a zero residual becomes a centered norm. Finite sums with a convex locally $C^{1,1}$ smooth part reduce to Theorem~\ref{F:centered}. For a more general gauge, local $C^{1,1}$ smoothness at that particular nonzero residual must be checked; a spectral penalty may fail it.

\begin{theorem}[Exact linear-image lifting]\label{F:lifting}
Let $A\in\mathbb R^{r\times n}$ have full row rank and $F(x)=\bar F(Ax)$. Suppose the projected admissible run converges to a unique optimizer $z_*$ of $\bar F$, under any of the retained sufficient hypotheses. Then
\[
 x_k\to x_0+H_0A^T(AH_0A^T)^{-1}(z_*-Ax_0).
\]
\end{theorem}
\begin{proof}
With $z_k=Ax_k$ and $\bar H_k=AH_kA^T$, projection of the inverse-BFGS update gives exactly full BFGS on $\bar F$, with the same step lengths and line-search comparisons. Let $S=\operatorname{range}A^T$. Since gradients and secant differences lie in $S$, the update gives $H_{k+1}S\subset H_kS$; both spaces have dimension $r$, so equality holds. All displacements lie in $H_0S$. The restriction of $A$ to that space has right inverse
\[
 L_0=H_0A^T(AH_0A^T)^{-1}.
\]
Thus $x_k-x_0=L_0(z_k-z_0)$, proving the result. This is the $H_0^{-1}$-metric projection onto the affine fiber, not a claim about arbitrary nonunique minimizer sets.
\end{proof}

\section{Finite-horizon rates and strengthened asymptotic corollaries}
Let $S_N=\sum_{k<N}q_k^{-1}$. Retain the explicit lower bound \eqref{M:finitehorizon}, which is meaningful at a fixed horizon. All rates below concern fixed dimension, problem, initialization and line-search constants, and count accepted iterations only. The source rate note recorded big-$O$ bounds. The little-$o$ versions below are explicit consolidation deductions from the later superquadratic theorem; they do not assert a sharp rate.

\paragraph{Uniformly bounded reference correction.}
If $R_k\le\bar R$, then \eqref{M:energybound} gives
\[
 F(x_N)-F_*\le\sqrt{\frac{E_0+\bar RN}{S_N}}=o(N^{-1/2}),
\]
because $S_N/N^2\to\infty$. The same estimate holds for a linear cumulative correction budget. This is not a rate for arbitrary band-dependent certificates.

\paragraph{Bounded-below degree-one objectives.}
Module H proves
\[
 \sum_{k<N}\frac{F(x_k)^2}{q_k}\le\Phi_0,
\quad \Phi_0=x_0^TH_0^{-1}x_0+\frac{1-c_1}{c_1}F(x_0).
\]
Hence $F(x_N)\le\sqrt{\Phi_0/S_N}=o(N^{-1})$. Coercivity gives $F(x)\ge\mu\|x\|$ and the same order for the iterate norm. The little-$o$ conclusion uses the superquadratic theorem; the finite-horizon bound retains its explicit constants. No rate sharpness is asserted.

\paragraph{Quadratic plus coercive gauge.}
For $F(u,v)=u^TQu/2+\gamma(v)$, $Q\succ0$, let $\mu\|v\|\le\gamma(v)\le\Lambda\|v\|$. Split $t=a+b$ with
\[
 a=\delta u^TQ\delta u,\quad b=(h'-h)^T\delta v,\quad
 \xi=(Q\delta u)^Tu',\quad\zeta=(h'-h)^Tv',
\]
where $h\in\partial\gamma(v)$ and $h'\in\partial\gamma(v')$. Then
\[
 \xi^2/a\le u'^TQu',\qquad0\le\zeta\le b,
 \qquad\zeta\le(1+\Lambda/\mu)\gamma(v').
\]
Thus
\[
 R_k\le C_\gamma F(x_{k+1}),\qquad
 C_\gamma=\max\{2,1+\Lambda/\mu\}.
\]
Choose $a_0>0,N_0$ with $S_N\ge a_0N^2$ for $N\ge N_0$, and put
\[
 A_N=E_0+C_\gamma\sum_{j=1}^NF(x_j).
\]
Then $F(x_N)\le\sqrt{A_N}/(\sqrt{a_0}N)$, and
\[
 \sqrt{A_N}-\sqrt{A_{N-1}}
 =\frac{C_\gamma F(x_N)}{\sqrt{A_N}+\sqrt{A_{N-1}}}
 \le\frac{C_\gamma}{\sqrt{a_0}N}.
\]
Summing first recovers the source bounds $\sqrt{A_N}=O(\log N)$ and $F(x_N)=O(\log N/N)$. For the strengthened conclusion, keep the exact denominator:
\[
 \sqrt{A_N}-\sqrt{A_{N-1}}\le \frac{C_\gamma}{\sqrt{S_N}}=o(N^{-1}).
\]
A sum of $o(1/N)$ increments is $o(\log N)$: for any fixed $\epsilon>0$, bound the sufficiently late increments by $\epsilon/N$, then divide their sum by $\log N$ and let $\epsilon\downarrow0$. Thus $\sqrt{A_N}=o(\log N)$. Dividing by $\sqrt{S_N}=N\sqrt{S_N/N^2}$ yields
\[
 \boxed{F(x_N)=o(\log N/N).}
\]
These strengthened asymptotic corollaries follow from already recorded estimates; they are not uniform complexity bounds over all instances and do not establish linear or superlinear convergence.

\section{Matrix limits and line-search existence are different questions}
The convergence theorems do not assert a universal inverse-Hessian limit. For $u^2+|v|$, choose $A>0$, $v_0>0$ and $0<\rho<(1-c_1)/(1+c_1)$. With
\[
 x_0=(0,v_0),\quad
 H_0=\operatorname{diag}\left(A,\frac{1+\rho}{2\rho}v_0\right),\quad\alpha_k=2\rho,
\]
ordinary BFGS gives
\[
 x_k=(0,(-\rho)^kv_0),\quad
 H_k=\operatorname{diag}\left(A,\frac{1+\rho}{2\rho}\rho^kv_0\right).
\]
The Armijo ratio is $(1-\rho)/(1+\rho)>c_1$ and the weak-Wolfe ratio is $-1$. The unexcited tangential entry tends to its arbitrary initial value $A$, not necessarily the inverse true tangential curvature.

For a convex bounded-below line restriction with negative initial derivative $a$, let $\tau$ be the first time its right derivative reaches $c_2a$. Before $\tau$, its slopes are at most $c_2a$, giving strict Armijo because $c_2a<c_1a$. Nearby times to the right satisfy weak Wolfe. If full-space nondifferentiability occurs at only finitely many ray times, a differentiable acceptable point can be selected. This applies to finite affine absolute-value/hinge sums and suitable Euclidean residual sums. It is an existence argument, not a complexity guarantee for every implementation.

For arbitrary max-functions there is a genuine exception. The coercive function
\[
 h(x,y,z)=\max\{-x,y+z,-y+z,x/4,-z\}
\]
has unique minimizer zero. At $(-2,0,1)$ with $H_0=I$, the direction is $(1,0,0)$ and the ray function is $\max\{2-\alpha,1,(\alpha-2)/4\}$. For $(c_1,c_2)=(1/4,1/2)$, times $0<\alpha<1$ fail weak Wolfe; $1\le\alpha\le6$ are full-space nondifferentiable; and $\alpha>6$ fails Armijo, which would require $\alpha\le5$. Scaling the point and matrix places the same finite obstruction arbitrarily close to the minimizer. It is not an infinite admissible run and does not contradict the positive theorems.

% END original chapters/families.tex

\StartModule{H}{Homogeneous objectives without compactness}{Selected proof from source [SH]. Retains the value dichotomy, affine-hull obstruction, low-effective-rank iterates and affine-fiber cases. Superseded matrix arguments are replaced by Module M.}

% BEGIN original chapters/H.tex
\begingroup

\setcounter{section}{0}
\setcounter{theorem}{0}
\setcounter{lemma}{0}
\setcounter{corollary}{0}
\setcounter{proposition}{0}
\setcounter{remark}{0}
\setcounter{equation}{0}
\section{Scope and main
conclusions}\label{H:scope-and-main-conclusions}

This note continues the investigation of ordinary, full-memory BFGS
under the algorithmic formulation in Xie--W\"achter~\cite{XW}. The accepted
steps satisfy fixed Armijo and \textbf{weak} Wolfe inequalities.
Accepted iterates are differentiable; the inverse matrices are positive
definite; no damping, restarts, smoothing, spectral safeguards, or
particular line-search implementation is imposed.

The principal family is

\[
f(x)=\max_{1\le i\le m}a_i^\top x,\qquad m<\infty.
\]

The objective-value theorem below holds more generally for every finite
convex, positively homogeneous function on the whole space. ``Finite''
in this broader statement means real-valued everywhere, not a finite
number of pieces.

\textbf{Research status.} The arguments below are new derivations in
this investigation. They have been checked algebraically and accompanied
by exact-rational finite-step tests. They have not undergone independent
mathematical peer review, and no exhaustive priority claim is made. The
infinite-sequence statements rest on the proofs, not on the finite
tests.

\subsection{Theorem H.A: an unconditional objective-value
dichotomy}\label{H:theorem-a-an-unconditional-objective-value-dichotomy}

Let \(f:\mathbb R^n\to\mathbb R\) be finite, convex, and positively
homogeneous of degree one. Every infinite admissible full-BFGS run
satisfies

\[
\boxed{f(x_k)\longrightarrow0\quad\text{or}\quad f(x_k)\longrightarrow-\infty.}
\]

There is \textbf{no assumption} that the iterates, displacements,
accepted step lengths, or matrices are bounded. There is also no
assumption of a finite gradient alphabet for this theorem.

Consequently, if \(f\) is bounded below, then \(f(x_k)\to0=\min f\). If
\(f\) is coercive, then \(x_k\to0\), its unique minimizer. If an
infinite run ever attains \(f(x_k)\le0\), its values must tend to
\(-\infty\); the finite-value branch has \(f(x_k)>0\) at every
iteration.

\subsection{Theorem H.B: a many-piece extension of the common-tilt
geometry}\label{H:theorem-b-a-many-piece-extension-of-the-common-tilt-geometry}

For a finite max-linear objective, let \(G\) be the set of gradients
occurring infinitely often. If the objective values have a finite limit,
then

\[
\boxed{0\in\operatorname{aff}G.}
\]

This no longer needs a bounded-displacement hypothesis. In particular,

\[
0\notin\operatorname{aff}\{a_1,\ldots,a_m\}
\quad\Longrightarrow\quad
f(x_k)\longrightarrow-\infty
\]

for every infinite admissible run. A useful arbitrary-many-piece family
is

\[
\boxed{f(u,t)=t+\max_{1\le i\le m}b_i^\top u.}
\]

All its gradients have last component one, so every infinite admissible
run drives the values to \(-\infty\). This is an infinite-run result,
not a proof of finite termination for every many-piece objective.

\subsection{Theorem H.C: complete iterate dichotomy in the planar
unbounded-below
case}\label{H:theorem-c-complete-iterate-dichotomy-in-the-planar-unbounded-below-case}

For a finite max-linear objective on \(\mathbb R^2\) with
\(\inf f=-\infty\), every infinite admissible run satisfies exactly one
of

\[
\boxed{x_k\longrightarrow0,\quad f(x_k)\downarrow0,}
\]

or

\[
\boxed{f(x_k)\longrightarrow-\infty,\quad\|x_k\|\longrightarrow\infty.}
\]

In the first alternative the limit is nonstationary. The previously
constructed exact three-piece example attains that alternative. In the
second, convergence of the \emph{directions} \(x_k/\|x_k\|\) is not
asserted.

More generally, for any finite max-linear objective, if the span of the
recurrent gradients has dimension at most two, every finite-value-limit
run has convergent iterates. Section H.7 proves this broader statement,
including the noncoercive bounded-below case.

\subsection{What this does not yet
establish}\label{H:what-this-does-not-yet-establish}

The objective-value classification is now unconditional in every
dimension. The remaining unresolved iterate question in this
investigation is the zero-value branch for noncompact objectives when
the recurrent gradient span has dimension at least three. We have
neither proved that all such iterate sequences converge nor constructed
a nonconvergent zero-value-limit trajectory.

This distinction is essential: the results close the previous
unbounded-displacement gap for \textbf{function values}, but are not
being presented as a full all-dimensional classification of the
\textbf{iterates}.


\setcounter{section}{1}
\setcounter{theorem}{0}
\setcounter{lemma}{0}
\setcounter{corollary}{0}
\setcounter{proposition}{0}
\setcounter{remark}{0}
\setcounter{equation}{0}
\section{Notation and BFGS
identities}\label{H:notation-and-bfgs-identities}

Write

\[
\begin{aligned}
g_k&=\nabla f(x_k),&p_k&=-H_kg_k,&s_k&=\alpha_kp_k,\\
x_{k+1}&=x_k+s_k,&y_k&=g_{k+1}-g_k,&B_k&=H_k^{-1},\\
q_k&=g_k^\top H_kg_k>0,&d_k&=-g_k^\top s_k=\alpha_kq_k>0,\\
t_k&=y_k^\top s_k,&\Delta_k&=f(x_k)-f(x_{k+1}).
\end{aligned}
\]

The line-search conditions are

\[
f(x_{k+1})\le f(x_k)+c_1g_k^\top s_k,
\qquad
g_{k+1}^\top s_k\ge c_2g_k^\top s_k,
\qquad0<c_1<c_2<1.
\]

In particular, with \(c=1-c_2>0\),

\[
t_k\ge c\,d_k>0,
\qquad
\beta_k:=\frac{t_k}{d_k}
=-\frac{y_k^\top H_kg_k}{q_k}\ge c.
\]

The full inverse-BFGS update is

\[
H_{k+1}=
\left(I-\frac{s_ky_k^\top}{t_k}\right)
H_k
\left(I-\frac{y_ks_k^\top}{t_k}\right)
+\frac{s_ks_k^\top}{t_k}.
\]

Its Hessian form, after substituting \(s_k=-\alpha_kH_kg_k\), is

\[
\boxed{B_{k+1}=B_k-\frac{g_kg_k^\top}{q_k}
+\frac{y_ky_k^\top}{t_k}.}
\tag{H.2.1}
\]

A zero gradient corresponds to a stationary termination rather than an
infinite strict-descent run. Throughout, ``infinite admissible''
therefore includes nonzero accepted gradients and well-defined
positive-curvature updates.

\Needspace{19\baselineskip}


\setcounter{section}{4}
\setcounter{theorem}{0}
\setcounter{lemma}{0}
\setcounter{corollary}{0}
\setcounter{proposition}{0}
\setcounter{remark}{0}
\setcounter{equation}{0}
\section{Homogeneous potential and proof of Theorem
A}\label{H:homogeneous-potential-and-proof-of-theorem-a}

\subsection{5.1 Bounded subgradients and Euler's
identity}\label{H:bounded-subgradients-and-eulers-identity}

For finite convex positively homogeneous \(f\), every subgradient
\(g\in\partial f(x)\) satisfies

\[
g^\top x=f(x),\qquad g^\top z\le f(z)\quad\text{for every }z.
\]

The first equality follows by applying the subgradient inequality to
\(0\) and \(2x\); the second follows by substituting that equality back
into the inequality. In particular, all subgradients belong to
\(\partial f(0)\). Since \(f\) is finite and continuous on the unit
sphere,

\[
\|g\|\le\max_{\|z\|=1}f(z)<\infty.
\]

Thus the bounded-gradient assumption in Theorem~\ref{M:average} is automatic, even for
an infinite-piece sublinear function.

\subsection{5.2 The exact energy
identity}\label{H:the-exact-energy-identity}

Set \(f_k=f(x_k)\) and

\[
h_k=y_k^\top x_{k+1}=d_k-\Delta_k,\qquad
\ell_k=-y_k^\top x_k.
\]

Convexity and homogeneity give \(h_k,\ell_k\ge0\) and
\(t_k=h_k+\ell_k\). Armijo and convexity give

\[
c_1d_k\le\Delta_k\le d_k.
\]

Expanding (2.1) at \(x_{k+1}=x_k-\alpha_kH_kg_k\) yields

\[
\boxed{
x_{k+1}^\top B_{k+1}x_{k+1}
=x_k^\top B_kx_k-\frac{f_k^2}{q_k}+\frac{h_k^2}{t_k}.
}
\tag{H.5.1}
\]

The explicit terms involving \(\alpha_k\) cancel because
\(g_k^\top x_k=f_k\):

\[
-2\alpha_k f_k+\alpha_k^2q_k
-\frac{(f_k-\alpha_kq_k)^2}{q_k}
=-\frac{f_k^2}{q_k}.
\]

Also,

\[
\frac{h_k^2}{t_k}\le h_k=d_k-\Delta_k
\le\frac{1-c_1}{c_1}\Delta_k.
\]

Consequently, with \(\gamma=(1-c_1)/c_1\),

\[
\Phi_k=x_k^\top B_kx_k+\gamma f_k,
\qquad
\boxed{\Phi_{k+1}\le\Phi_k-\frac{f_k^2}{q_k}.}
\tag{H.5.2}
\]

This potential inequality was established in the earlier note. The new
ingredient is Theorem~\ref{M:average}, which replaces its former bounded-step
reciprocal-sum argument.

\subsection{5.3 Every finite value limit is
zero}\label{H:every-finite-value-limit-is-zero}

Armijo gives strictly decreasing values. Either they tend to \(-\infty\)
or they have a finite limit \(L\).

In the latter case,

\[
\sum_k d_k\le\frac{f_0-L}{c_1}<\infty,
\]

so \(d_k\to0\) and Theorem~\ref{M:average} gives \(\liminf q_k=0\). On the other hand,
\(\Phi_k\ge\gamma L\), and telescoping (5.2) gives

\[
\sum_k\frac{f_k^2}{q_k}<\infty.
\]

If \(L\ne0\), this would imply \(\sum_k1/q_k<\infty\) and therefore
\(q_k\to\infty\), a contradiction.

There is also a direct pointwise argument. The potential bounds

\[
x_k^\top B_kx_k\le\Phi_0-\gamma L=:M<\infty.
\]

Metric Cauchy--Schwarz and Euler's identity imply

\[
f_k^2=(g_k^\top x_k)^2
\le(g_k^\top H_kg_k)(x_k^\top B_kx_k)
\le Mq_k.
\]

A finite nonzero limit would thus force \(q_k\) to be bounded away from
zero, again contradicting Theorem~\ref{M:average}. Hence \(L=0\), proving Theorem H.A.
\(\square\)

\subsection{5.4 Bounded-below and coercive
consequences}\label{H:bounded-below-and-coercive-consequences}

A positively homogeneous function bounded below cannot take a negative
value: otherwise \(f(tx)=tf(x)\to-\infty\). Thus its minimum value is
zero, attained at the origin, and Theorem H.A gives unconditional
objective convergence.

If its sublevels are bounded, then \(f(x)>0\) for \(x\ne0\). Compactness
of the unit sphere gives \(f(x)\ge\mu\|x\|\) for some \(\mu>0\), so
\(f_k\to0\) forces \(x_k\to0\).

These consequences now hold for all finite convex sublinear objectives,
not only finite polyhedral maxima. No rate or finite-path-length
estimate is asserted.


\setcounter{section}{5}
\setcounter{theorem}{0}
\setcounter{lemma}{0}
\setcounter{corollary}{0}
\setcounter{proposition}{0}
\setcounter{remark}{0}
\setcounter{equation}{0}
\section{Affine-hull obstruction and a many-piece escape
theorem}\label{H:affine-hull-obstruction-and-a-many-piece-escape-theorem}

\textbf{Lemma H.3.} Suppose a full-BFGS run satisfies Armijo--weak-Wolfe,
has a finite gradient alphabet, and has objective values bounded below
along the run. Let \(G\) be the recurrent gradients. Then
\(0\in\operatorname{aff}G\).

\textbf{Proof.} Finite objective decrease gives \(d_k\to0\). Bounded
gradients and Theorem~\ref{M:average} give \(\liminf q_k=0\), hence
\(\sum_k1/q_k=\infty\).

Discard a finite prefix so that all remaining gradients belong to \(G\).
If \(0\notin\operatorname{aff}G\), a vector \(v\) exists with

\[
v^\top g=1\quad(g\in G).
\]

In particular \(v^\top y_k=0\). Equation (2.1) then implies

\[
v^\top B_{k+1}v=v^\top B_kv-\frac1{q_k}.
\]

Positive definiteness forces the reciprocal sum to be finite, a
contradiction. \(\square\)

This lemma does not need homogeneity, or even convexity, beyond what may
be needed to obtain the specified run. It is a necessary condition for
any finite-gradient, finite-value-limit trajectory under these BFGS and
line-search assumptions.

For max-linear objectives, if \(0\notin\operatorname{aff}\{a_i\}\), no
recurrent subset can have an affine hull containing zero. A finite value
limit is impossible; monotonicity therefore forces \(f_k\to-\infty\).
This proves Theorem H.B.

\subsection{Connection to the two-piece
paper}\label{H:connection-to-the-two-piece-paper}

The model \(|u|+t\) has gradients \((1,1)\) and \((-1,1)\), both with
last component one. The arbitrary-many-piece model

\[
f(u,t)=t+\max_i b_i^\top u
\]

has precisely the same common-component property. The new theorem proves
that an infinite accepted-iterate sequence cannot stagnate on any finite
value for this entire family.

Xie--W\"achter~\cite{XW} proves the stronger finite-termination conclusion
for its two-piece model: a search direction exposing unboundedness is
reached after finitely many iterations. The argument here does
\textbf{not} establish finite termination for the general many-piece
family. It does show that any infinite admissible continuation must
decrease the values without bound.

If \(r=\dim\operatorname{span}G\), a finite-value-limit run needs at
least \(r+1\) distinct recurrent gradients. Indeed,
\(0\in\operatorname{aff}G\) implies that the gradient differences span
the same \(r\)-dimensional space; \(|G|-1\) differences must span it.
The earlier three-gradient, rank-two stagnation example attains this
lower bound.


\setcounter{section}{6}
\setcounter{theorem}{0}
\setcounter{lemma}{0}
\setcounter{corollary}{0}
\setcounter{proposition}{0}
\setcounter{remark}{0}
\setcounter{equation}{0}
\section{Iterate convergence in effective dimension at most
two}\label{H:iterate-convergence-in-effective-dimension-at-most-two}

\subsection{7.1 Exact invariant subspace and
lifting}\label{H:exact-invariant-subspace-and-lifting}

Suppose that, after some index \(K\), all gradients lie in a subspace
\(S\). Let \(U\) have orthonormal columns spanning \(S\), and define

\[
\widehat x_k=U^\top x_k,\qquad
\widehat H_k=U^\top H_kU,\qquad
\widehat g_k=U^\top g_k.
\]

Projection of the inverse-BFGS formula gives exactly the same
inverse-BFGS update for the hatted variables and projected steps. The
scalars \(q_k,d_k,t_k\) and the line-search inequalities are unchanged.

Moreover, all subsequent steps belong to the \textbf{fixed} subspace
\(H_KS\). To see this, initially \(H_KS\) is the image of \(S\); if
\(H_kS\subset H_KS\), then \(s_k\in H_KS\), and the inverse update maps
\(S\) back into \(H_KS\), since \(g_k,y_k\in S\). Induction proves the
claim.

The restriction of \(U^\top\) to \(H_KS\) is invertible because
\(U^\top H_KU\succ0\). Therefore

\[
\boxed{
x_k=x_K+L_K(\widehat x_k-\widehat x_K),\qquad
L_K=H_KU(U^\top H_KU)^{-1}.
}
\tag{H.7.1}
\]

In particular, convergence of the projected iterates implies convergence
of the original iterates. If \(\widehat x_k\to0\), the limit is

\[
\boxed{x_*=x_K-H_KU(U^\top H_KU)^{-1}U^\top x_K.}
\tag{H.7.2}
\]

The projected Hessian is \((U^\top H_kU)^{-1}\), not \(U^\top B_kU\).
This distinction is necessary for a valid reduction.

\subsection{7.2 A two-gradient endpoint
observation}\label{H:a-two-gradient-endpoint-observation}

Assume \(f_k\to0\). Then \(d_k\to0\), and every endpoint satisfies

\[
g_{k+1}^\top x_{k+1}=f_{k+1}\longrightarrow0,
\qquad
g_k^\top x_{k+1}=f_k-d_k\longrightarrow0.
\tag{H.7.3}
\]

In two dimensions, if \(g_k\) and \(g_{k+1}\) are linearly independent,
these two equations determine \(x_{k+1}\). There are only finitely many
distinct pairs, so their inverse norms are uniformly bounded. Thus the
endpoints of all independent-gradient transitions converge to zero, even
if the full sequence was not known to be bounded.

An unchanged nonzero gradient cannot satisfy weak Wolfe on a strict
descent step. Distinct gradients that are positive multiples of one
another also cannot both be active at positive-value iterates: among
positive multiples of a slope, only the largest multiple can attain a
positive maximum. Consequently, any dependent transition on the
positive-value branch is between opposite slopes.

\subsection{7.3 The unbounded-below planar
case}\label{H:the-unbounded-below-planar-case}

If the global objective is unbounded below, its slope set cannot contain
an opposite pair \(a,-\kappa a\) with \(\kappa>0\): their maximum alone
is nonnegative everywhere. Every accepted gradient transition is
therefore linearly independent in the plane.

Equation (7.3) now gives \(x_k\to0\) directly. Combined with Theorem H.A,
this proves Theorem H.C. In the other branch, global Lipschitz continuity
gives \(|f(x)|\le M\|x\|\), so \(f_k\to-\infty\) implies
\(\|x_k\|\to\infty\).

For an unbounded-below max-linear objective there is no stationary
point: a stationary point of a convex function would be a global
minimizer, and no minimizer exists. Thus the zero limit is genuinely
nonstationary.

\subsection{7.4 The general planar finite-value
branch}\label{H:the-general-planar-finite-value-branch}

For completeness, opposing-gradient transitions can also be handled,
giving convergence for \textbf{all} planar finite max-linear objectives,
including noncompact bounded-below ones.

Call an endpoint of an independent-gradient transition an \emph{anchor}.
By (7.3), anchors converge to zero.

Consider consecutive anchors \(x_a,x_b\). If there is no intermediate
iterate, there is nothing to prove. Otherwise, all gradients from
\(g_a\) through \(g_{b-1}\) lie on a single gradient line; the
intermediate dependent transitions alternate between opposite slopes. By
the invariant-subspace argument, all steps from \(x_a\) up to and
including the step arriving at \(x_b\) lie on a fixed primal line.
Therefore all points in this block are collinear with its endpoints.

Write its initial gradient as \(g\), and the next opposite gradient as
\(-\kappa g\), with \(\kappa>0\). The first dependent transition has

\[
d_a=f_a+f_{a+1}/\kappa.
\]

Armijo implies

\[
f_{a+1}\le\frac{1-c_1}{1+c_1/\kappa}f_a
\le(1-c_1)f_a,
\]

so monotonicity gives \(f_b\le(1-c_1)f_a\). Every accepted point in the
block obeys

\[
-f_a/\kappa\le g^\top x\le f_a,
\]

because \(g\) and \(-\kappa g\) are both slopes of the objective and
\(f(x)\le f_a\). The endpoint difference satisfies

\[
g^\top(x_b-x_a)\le-c_1 f_a.
\]

Writing \(x=x_a+\theta(x_b-x_a)\) consequently gives

\[
0\le\theta\le\frac{1+1/\kappa}{c_1}.
\]

Only finitely many slope pairs occur, so this bound is uniform across
blocks. Since both anchor endpoints tend to zero, every intermediate
point does too. Arbitrarily long blocks cause no difficulty.

If there are only finitely many anchors, the tail gradients lie on a
fixed line, and all tail iterates lie on one fixed primal line. The
slope coordinate tends to zero because \(f_k\to0\) and only finitely
many nonzero slope multipliers occur. Positive definiteness makes that
slope coordinate nonconstant along the primal line, so the entire tail
converges to its unique intersection with the zero-slope hyperplane.

This proves planar iterate convergence in every finite-value-limit case.
Reducing to the span of the recurrent gradients and using (7.1) proves
the same conclusion whenever that span has dimension at most two in any
ambient dimension. The limit may be nonzero in the original coordinates.


\setcounter{section}{7}
\setcounter{theorem}{0}
\setcounter{lemma}{0}
\setcounter{corollary}{0}
\setcounter{proposition}{0}
\setcounter{remark}{0}
\setcounter{equation}{0}
\section{A noncompact convergence theorem in arbitrary
dimension}\label{H:a-noncompact-convergence-theorem-in-arbitrary-dimension}

One source of noncompactness is a subspace of directions on which the
function is constant. When this is the only source of noncompactness,
the case is fully tractable in all dimensions.

Let

\[
S=\operatorname{span}\{a_1,\ldots,a_m\},\qquad
P=\operatorname{conv}\{a_1,\ldots,a_m\},
\]

and suppose \(0\in\operatorname{relint}P\), with relative interior taken
in \(S\). Then \(f\) depends only on the projection to \(S\) and is
coercive there:

\[
f(x)\ge\mu\|U^\top x\|\quad\text{for some }\mu>0.
\]

Theorem H.A implies \(U^\top x_k\to0\). The invariant subspace formula
with \(K=0\) gives

\[
\boxed{
x_k\longrightarrow
x_0-H_0U(U^\top H_0U)^{-1}U^\top x_0.
}
\tag{H.8.1}
\]

The limit is a minimizer, and in fact the projection of \(x_0\) onto
\(S^\perp\) in the \(H_0^{-1}\) metric. It is independent of the choices
of admissible accepted lengths. This includes all finite polyhedral
seminorms and noncoercive gauges whose minimizer set is a linear
subspace.

The same argument applies after a finite prefix if
\(0\in\operatorname{relint}\operatorname{conv}G\) for the recurrent
gradients. The projected tail objective is then coercive even if some
original, nonrecurrent slopes have disappeared. Formula (7.2) supplies
the limit.

For \(f\equiv0\), the algorithm is already stationary; there is no
nonzero-gradient infinite run. The projection formula above concerns a
nontrivial slope span.


\endgroup

% END original chapters/H.tex

\StartModule{Q}{Arbitrarily many parabolic branches}{Selected complete application from [ST]. Flat contour segments are permitted, so this family is not subsumed by Module C.}

% BEGIN original chapters/Q.tex
\begingroup
\providecommand{\R}{}
\renewcommand{\R}{\mathbb R}
\providecommand{\tr}{}
\renewcommand{\tr}{\operatorname{tr}}
\providecommand{\diag}{}
\renewcommand{\diag}{\operatorname{diag}}
\providecommand{\Span}{}
\renewcommand{\Span}{\operatorname{span}}
\providecommand{\ip}{}
\renewcommand{\ip}[2]{\langle #1,#2\rangle}
\providecommand{\norm}{}
\renewcommand{\norm}[1]{\lVert #1\rVert}
\providecommand{\T}{}
\renewcommand{\T}{^{\mathsf T}}
\providecommand{\D}{}
\renewcommand{\D}{\mathcal D}
\setcounter{section}{4}
\setcounter{theorem}{1}
\setcounter{lemma}{0}
\setcounter{corollary}{0}
\setcounter{proposition}{0}
\setcounter{remark}{0}
\setcounter{equation}{6}
\section{A new application: arbitrarily many curved quadratic branches}
\begin{theorem}[Finite multi-branch parabolic maxima]
Let $m<\infty$ and
\[
 F(u,v)=\max_{i\le m}\{a_i u^2+b_i v\},\qquad
 a_i>0,\quad\min_i b_i<0<\max_i b_i.
\]
Every infinite admissible unmodified full-BFGS run, for arbitrary $H_0\succ0$ and fixed $0<c_1<c_2<1$, satisfies
\[
 (u_k,v_k)\to(0,0),\qquad F(u_k,v_k)\to0.
\]
The number of essential branches is arbitrary. In particular, branches with $b_i=0$, giving a region flat in the normal coordinate, are allowed.
\end{theorem}
This is a concrete application of the objective-pair theorem, not the definition of its scope. The abstract condition is dimension-independent; this explicitly verified new class has two variables.

\subsection{Compactness and the normalized upper envelope}
Set $a_{\min}=\min_i a_i>0$. A convex combination of one positive and one negative normal slope has slope zero, so
\[
 F(u,v)\ge a_{\min}u^2.
\]
The positive extreme slope bounds $v$ above on a sublevel, and the negative extreme bounds it below. Thus every sublevel is compact. At $u=0$, $F(0,v)>0$ for $v\ne0$, proving uniqueness of the minimum zero.

Discard globally redundant pieces, merge equal normal slopes by retaining the largest $a_i$, and remove pieces active only on a tie. Order the remaining pieces by strictly increasing $b_i$. They form the upper envelope
\[
 \varphi(\xi)=\max_i(a_i+b_i\xi),\qquad \xi=v/u^2\quad(u\ne0).
\]
The neighboring breakpoints $\xi_i$ increase strictly. Put
\[
 A_i=a_i+b_i\xi_i=a_{i+1}+b_{i+1}\xi_i>0.
\]
Choose constants $k_i$ recursively:
\begin{equation}
 k_{i+1}-k_i=-\frac{a_{i+1}-a_i}{2A_i}. \tag{Q.7}
\end{equation}
At a differentiable point with active branch $i$, define
\begin{equation}
 r(u,v)=(u,2v)\T,\qquad
 \theta(u,v)=\exp\left\{\frac{a_i u^2}{2F(u,v)}+k_i\right\},
 \qquad W=\theta^2rr\T. \tag{Q.8}
\end{equation}
The constants make $\theta$ agree across every switching interface. At a tie $v=\xi_i u^2$ the exponent difference is
\[
 \frac{a_{i+1}-a_i}{2A_i}+k_{i+1}-k_i=0.
\]
The field is well-defined on $F>0$, including by continuous extension across ties. It is only used at differentiable endpoints in the theorem. Since $0\le u^2/F\le1/a_{\min}$, there are global constants
\[
 0<\theta_{\min}\le\theta\le\theta_{\max}<\infty.
\]
There is no coordinate change in the BFGS algorithm. The normalized upper envelope is used solely to construct an analytical certificate.

\subsection{Visibility and secant correction}
Fix a compact positive band $\eta\le F\le B$. For old branch $i$ at $x=(u,v)$ and new branch $j$ at $z=(u',v')$, set
\[
 h=u'-u,\quad \ell=v'-v,\quad s=(h,\ell),
\]
\[
 e=F(z)-(a_i u'^2+b_i v')\ge0,\qquad
 \bar e=F(x)-(a_j u^2+b_jv)\ge0.
\]
Exact identities are
\begin{align}
 d-\Delta&=a_i h^2+e, \tag{Q.9}\\
 t&=(a_i+a_j)h^2+e+\bar e, \tag{Q.10}\\
 g\T r(x)&=2F(x), \tag{Q.11}\\
 y\T r(z)&=2a_i h u'+2e. \tag{Q.12}
\end{align}
For eligible pairs, $h^2\le d/a_{\min}$ and $e\le d$. Weighted Cauchy--Schwarz gives
\[
 \frac{(y\T r(z))^2}{t}
 \le \frac{4a_i^2}{a_i+a_j}u'^2+4e.
\]
The right side is uniformly bounded on the compact band for small $d$. Multiplication by $\theta(z)^2$ proves (S). Equation (11) proves (V) with $m=4\theta_{\min}^2\eta^2$.

\subsection{The step quotient cannot degenerate on an eligible small-decrease pair}
Define
\[
 D=\det(s,r(x))=2vh-u\ell.
\]
We next prove, uniformly on the positive band and for sufficiently small eligible $d$,
\begin{equation}
 d\le C|D|,\qquad |h|\le C|D|. \tag{Q.13}
\end{equation}
This is the essential denominator estimate. It does not assume $s\to0$ in both coordinates or that successive active pieces are adjacent.

Let $N=y\T r(z)$. Since $r(z)-r(x)=\diag(1,2)s$,
\[
 y\T r(x)=-2t+N+y_u h.
\]
By (12), bounded endpoints and gradients, and $e\le d$,
\[
 y\T r(x)=-2t+\varepsilon,\qquad |\varepsilon|\le C(|h|+d).
\]
The two-dimensional determinant identity yields
\[
 \det(g,y)D=(g\T s)(y\T r(x))-(g\T r(x))(y\T s)
 =-d\,y\T r(x)-2F(x)t.
\]
Consequently
\[
 (2F(x)-2d)t\le C|D|+Cd(|h|+d).
\]
Use $F(x)\ge\eta$, $t\ge(1-c_2)d$, and $|h|\le\sqrt{d/a_{\min}}$. Taking $d$ sufficiently small absorbs the final term into the left side, proving $d\le C|D|$. The exact identity
\begin{equation}
 b_iD=2F(x)h+ud \tag{Q.14}
\end{equation}
then proves $|h|\le C|D|$, including when $b_i=0$. In particular $D\ne0$.

\subsection{The rescaling cancels first-order transport}
Since $D\ne0$, define
\[
 \lambda=1+\frac{h\ell}{D}.
\]
Then $r(z)-\lambda r(x)$ is parallel to $s$. A direct identity using $d=-2a_iuh-b_i\ell$ is
\begin{equation}
 \lambda-1+\frac{a_i u}{F(x)}h
 =-\frac{vhd}{F(x)D}. \tag{Q.15}
\end{equation}
By (13), the right side is $O(d)$ uniformly on the band. Since $h^2=O(d)$,
\[
 \lambda=1+O(\sqrt d)>0,\qquad
 \log\lambda=-\frac{a_i u}{F(x)}h+O(d).
\]
For the old branch's exponent $L_i(x)=a_i u^2/(2F(x))+k_i$, smoothness in the two scalar arguments $(u,F)$ on the positive band gives
\[
 L_i(z)-L_i(x)=\frac{a_i u}{F(x)}h+O(h^2+\Delta)
 =\frac{a_i u}{F(x)}h+O(d).
\]
Thus $\log\lambda+L_i(z)-L_i(x)=O(d)$.

We still need to account for a change from branch $i$ to branch $j$. The finite-envelope geometry implies
\begin{equation}
 |L_j(z)-L_i(z)|\le C\frac{e}{F(z)}\le C_\eta d. \tag{Q.16}
\end{equation}
A complete justification is as follows. Normalize a point by
\[
 t_0=u'^2/F(z),\qquad v_0=v'/F(z).
\]
It lies on the compact polygonal arc
\[
 S=\{(t_0,v_0):t_0\ge0,\ \max_i(a_it_0+b_iv_0)=1\}.
\]
Compactness follows from $t_0\le1/a_{\min}$ and the two nonzero extreme normal slopes. On the segment where $j$ is active, both
\[
 e/F(z)=1-a_it_0-b_iv_0,\qquad
 L_j-L_i=(a_j-a_i)t_0/2+k_j-k_i
\]
are affine functions. The first is nonnegative. Whenever it vanishes, the second vanishes by (7). On an adjacent segment the two functions are proportional; on a nonadjacent segment the first has a positive minimum. There are finitely many segments and pairs, proving (16). Removal of zero-length envelope pieces avoids spurious triple junctions; it does not change the objective or its differentiable gradients.

It follows that the compensated transport multiplier
\[
 \widetilde\lambda=\theta(z)\lambda/\theta(x)
\]
satisfies
\begin{equation}
 |\log\widetilde\lambda|\le C_\eta d. \tag{Q.17}
\end{equation}
Because
\[
 \Pi_s(\theta(z)r(z))=\widetilde\lambda\Pi_s(\theta(x)r(x)),
\]
we obtain
\[
 \Pi_sW(z)\Pi_s\preceq e^{2C_\eta d}\Pi_sW(x)\Pi_s
 \preceq(1+L_\eta d)\Pi_sW(x)\Pi_s
\]
on a sufficiently small-decrease tail. This is (T). The general theorem proves convergence from every admissible initialization. \qed


\setcounter{section}{5}
\setcounter{theorem}{2}
\setcounter{lemma}{0}
\setcounter{corollary}{0}
\setcounter{proposition}{0}
\setcounter{remark}{0}
\setcounter{equation}{13}
\section{Explicit three-branch example and strict nonnecessity of RS}
Consider
\begin{equation}
 F(u,v)=\max\{u^2+v,\ 3u^2,\ 2u^2-v\}. \tag{Q.18}
\end{equation}
The switching parabolas are $v=-u^2$ and $v=2u^2$. Between them the objective is flat in $v$, but the origin is its unique minimizer. The theorem applies globally.

An especially simple scaling for (8) is
\[
 \theta(u,v)=
 \begin{cases}
 \exp\{u^2/F-1/3\},&v<-u^2,\\
 1,&-u^2<v<2u^2,\\
 \exp\{u^2/(2F)-1/6\},&v>2u^2.
 \end{cases}
\]
It is continuous across both interfaces, and $e^{-1/3}\le\theta\le1$. Take $W=\theta^2(u,2v)(u,2v)\T$.

Fix $u\ne0$, $A>0$, and $0<\epsilon<6u^2$. Let
\[
 x_\epsilon=(u,2u^2+\epsilon),\qquad
 z_\epsilon=(u,2u^2-\epsilon/2),
\]
\[
 H_\epsilon=
 \begin{pmatrix}
 A&-2uA\\-2uA&4u^2A+3\epsilon/2
 \end{pmatrix}\succ0.
\]
The gradients are $g=(2u,1)\T$ and $g'=(6u,0)\T$. Directly,
\[
 -H_\epsilon g=z_\epsilon-x_\epsilon,
 \quad d=t=3\epsilon/2,\quad \Delta=\epsilon.
\]
Thus $\alpha=1$ is a valid BFGS step, with strict Armijo whenever $c_1<2/3$ and strict weak Wolfe for every allowed $c_2>0$.

Here $y=(4u,-1)\T$, so
\[
 \frac{(y\T z_\epsilon)^2}{t}
 =\frac{(2u^2+\epsilon/2)^2}{3\epsilon/2}\longrightarrow\infty.
\]
The values tend to the positive level $3u^2$ while $d\to0$. Such pairs exist in every neighborhood of zero. The original RS condition, even restricted to eligible near-stalling pairs for ordinary line-search constants, fails.

In contrast, $\theta(z_\epsilon)=1$ and
\[
 y\T r(z_\epsilon)=\epsilon,\qquad
 \frac{y\T W(z_\epsilon)y}{t}=\frac{2\epsilon}{3}\to0.
\]
The unscaled quotient multiplier is $\lambda=1$, and the compensated multiplier satisfies exactly
\[
 \log\widetilde\lambda
 =\frac{\epsilon}{6(3u^2+\epsilon)}=O(d)
\]
on each fixed positive band. This displays both parts of the mechanism: removing the radial mismatch makes the secant correction small, while a mild objective-defined rescaling controls accumulated transport.


\endgroup

% END original chapters/Q.tex

\StartModule{G}{What certificates mean and cannot imply}{Selected geometric proofs from [SG]. Certificate nonexistence is distinct from a bad BFGS trajectory. The scalar convergence theorem shows strict nonnecessity.}

% BEGIN original chapters/G.tex
\begingroup
\providecommand{\R}{}
\renewcommand{\R}{\mathbb R}
\providecommand{\PSD}{}
\renewcommand{\PSD}{\succeq0}
\providecommand{\range}{}
\renewcommand{\range}{\operatorname{range}}
\providecommand{\aff}{}
\renewcommand{\aff}{\operatorname{aff}}
\providecommand{\conv}{}
\renewcommand{\conv}{\operatorname{conv}}
\providecommand{\norm}{}
\renewcommand{\norm}[1]{\left\|#1\right\|}
\providecommand{\ip}{}
\renewcommand{\ip}[2]{\langle #1,#2\rangle}
\makeatletter\let\c@lemma\c@theorem\let\c@corollary\c@theorem\let\c@proposition\c@theorem\let\c@remark\c@theorem\makeatother
\renewcommand{\thelemma}{\thetheorem}\renewcommand{\thecorollary}{\thetheorem}\renewcommand{\theproposition}{\thetheorem}\renewcommand{\theremark}{\thetheorem}
\renewcommand{\thetheorem}{G.\arabic{section}.\arabic{theorem}}
\setcounter{section}{0}
\setcounter{theorem}{0}
\setcounter{equation}{0}
\section{The precise condition under study}
Let $F:\R^n\to\R$ be finite and convex, with attained minimum $F_*$. Fix a compact sublevel $K=\{F\le b\}$, a positive band
\[
 K_\eta=\{x\in K:F(x)\ge F_*+\eta\},\qquad \eta>0,
\]
and fixed $0<c_1<c_2<1$. Endpoints below are differentiable. Write
\[
 g=\nabla F(x),\quad g_z=\nabla F(z),\quad y=g_z-g,\quad s=z-x,
\]
\[
 d=-g^Ts,\quad \Delta=F(x)-F(z),\quad t=y^Ts.
\]
A pair is \emph{eligible} when
\begin{equation}
 d>0,\qquad c_1d\le\Delta\le d,\qquad t\ge(1-c_2)d. \tag{G.E}
\end{equation}
These are the necessary endpoint inequalities of the Armijo--weak-Wolfe formulation in Xie--W\"achter~\cite{XW}. In particular $t>0$ and $s\ne0$.

On each positive band, the condition asks for a finite-valued PSD field $W(x)$ and constants $m>0$, $C<\infty$, $L\ge0$, $\delta>0$ such that
\begin{align}
 g_x^TW(x)g_x&\ge m, \tag{G.V}\\
 y^TW(z)y&\le Ct, \tag{G.S}\\
 \Pi_sW(z)\Pi_s&\preceq(1+Ld)\Pi_sW(x)\Pi_s,
 \qquad \Pi_s=I-ss^T/(s^Ts). \tag{G.T}
\end{align}
The last two apply to all eligible pairs in the band with $d\le\delta$. The field may be singular, discontinuous, and unbounded as a function on the band; only its value at each point is finite. All constants may depend on the band and the fixed line-search parameters. No claim below silently imposes continuity or local boundedness on $W$.

The preceding theorem says that this condition implies objective convergence for every admissible infinite full-BFGS run in $K$, and iterate convergence under uniqueness. A local theorem requires the condition only near the minimizer, then uses a sufficiently low trapped sublevel. Failure on a large sublevel and failure in every neighborhood of a minimizer are different statements.


\setcounter{section}{1}
\setcounter{theorem}{0}
\setcounter{equation}{1}
\section{The expansion allowance can be removed completely}
\begin{proposition}[Exact nonexpansion normalization]
The existential condition (V)--(T) is equivalent to one in which $L=0$:
\begin{equation}
 \Pi_s\widehat W(z)\Pi_s\preceq\Pi_s\widehat W(x)\Pi_s. \tag{G.T0}
\end{equation}
The other constants may change, but no domain or eligible-pair hypothesis is added.
\end{proposition}
\begin{proof}
Given a field satisfying (T), put $\kappa=L/c_1$ and
\[
 \widehat W(x)=e^{\kappa(F(x)-F_*)}W(x).
\]
On the compact value band this multiplier is bounded above and away from zero, preserving finite values, (V), and (S), with changed constants. Since $\Delta\ge c_1d$,
\[
 e^{-\kappa\Delta}(1+Ld)\le e^{-Ld}(1+Ld)\le1.
\]
Multiplying (T) by $e^{\kappa(F(z)-F_*)}$ gives (T0). The converse is the case $L=0$.
\end{proof}
This is an equivalence, not a further relaxation. The accumulated-expansion budget was useful in the proof, but is not essential to the existential definition: an objective-value rescaling absorbs it.

\subsection{An ellipsoid interpretation}
Associate with $W(x)$ the possibly degenerate centered ellipsoid
\[
 \mathcal E_x=W(x)^{1/2}\{v:\norm v\le1\}.
\]
Its support function is $h_{\mathcal E_x}(a)=\sqrt{a^TW(x)a}$. Thus the normalized condition means:
\begin{enumerate}[label=(\roman*),itemsep=3pt]
\item $\mathcal E_x$ contains a direction with descent magnitude at least $\sqrt m$;
\item every incoming gradient change has magnitude at most $\sqrt{Ct}$ when tested on $\mathcal E_z$;
\item $\Pi_s\mathcal E_z\subseteq\Pi_s\mathcal E_x$ for every eligible small-decrease pair.
\end{enumerate}
Equivalently, after quotienting out the executed step direction, the comparison ellipsoid can only shrink. The ellipsoids are analytical comparison directions, not BFGS steps or approximations of $H_k$.

A necessary rank consequence is
\begin{equation}
 \range W(z)\subseteq\range W(x)+\operatorname{span}\{s\}. \label{G:eq:range}
\end{equation}
Indeed, Loewner ordering implies inclusion of the ranges of the projected PSD matrices. The range of $\Pi_sW\Pi_s$ is $\Pi_s\range W$, because it equals the range of $(\Pi_sW^{1/2})(\Pi_sW^{1/2})^T$. Lifting the projected inclusion proves \eqref{G:eq:range}. For a rank-one field, a new comparison direction must therefore lie in the plane spanned by the old comparison direction and the step.


\setcounter{section}{2}
\setcounter{theorem}{0}
\setcounter{equation}{3}
\section{Removing the matrix from pointwise feasibility}
Fix a band and a predicted-decrease threshold $\delta$. For each differentiable endpoint $z$, let $\mathcal I(z)$ be all differentiable predecessors $x$ for which $(x,z)$ is eligible, belongs to the band, and has $d\le\delta$. Define
\[
 \mathcal D(z)=\left\{v:\left|[\nabla F(z)-\nabla F(x)]^Tv\right|
       \le\sqrt{t(x,z)}\ \text{for every }x\in\mathcal I(z)\right\},
\]
\begin{equation}
 \nu(z)=\sup_{v\in\mathcal D(z)}|\nabla F(z)^Tv|\in[0,\infty]. \label{G:eq:nu}
\end{equation}
An empty predecessor set gives $\mathcal D(z)=\R^n$. The set $\mathcal D(z)$ is the polar of the closed absolutely convex hull of the normalized incoming jumps $y/\sqrt t$.

\begin{proposition}[Exact characterization of (V) and (S), without transport]
There is a finite-valued PSD field satisfying (V) and (S) with uniform positive $m$ and finite $C$ if and only if
\[
 \inf_z\nu(z)>0.
\]
This statement concerns only (V) and (S). It does not assert that the pointwise selections satisfy (T).
\end{proposition}
\begin{proof}
If $C>0$ and $W$ satisfies (V)--(S), set $h=\nabla F(z)$ and
\[
 v=\frac{W(z)h}{\sqrt C\sqrt{h^TW(z)h}}.
\]
PSD Cauchy--Schwarz gives $|y^Tv|\le\sqrt t$ for every incoming pair and
$|h^Tv|\ge\sqrt{m/C}$. If $C=0$, the same argument without $\sqrt C$ gives a nonzero visible direction annihilating every incoming jump, which can be scaled arbitrarily.

Conversely, if $\inf\nu\ge\nu_0>0$, choose, for each $z$, a finite vector $v_z\in\mathcal D(z)$ with $|\nabla F(z)^Tv_z|\ge\nu_0/2$. Then $W(z)=v_zv_z^T$ satisfies (S) with $C=1$ and (V) with $m=\nu_0^2/4$. No continuity is needed for these pointwise selections.
\end{proof}
Thus higher rank adds no pointwise power to (V)--(S); its potential advantage is compatibility of comparison geometries across endpoints through (T). The scalar $\nu(z)$ has a direct interpretation: how much descent remains visible after all low-curvature incoming jumps have been suppressed?

\subsection{A nonexistence test independent of the certificate}
\begin{proposition}[Incoming-jump obstruction]\label{G:prop:star}
Suppose $z_\epsilon$ and finitely many predecessors $x_{j,\epsilon}$ give eligible pairs in a fixed positive band with $d_{j,\epsilon}\to0$. Write $y_{j,\epsilon}=\nabla F(z_\epsilon)-\nabla F(x_{j,\epsilon})$ and $t_{j,\epsilon}=y_{j,\epsilon}^T(z_\epsilon-x_{j,\epsilon})$. If
\[
 \nabla F(z_\epsilon)=\sum_{j=1}^r\lambda_{j,\epsilon}y_{j,\epsilon},
 \qquad
 \sum_{j=1}^r|\lambda_{j,\epsilon}|\sqrt{t_{j,\epsilon}}\longrightarrow0,
\]
then no field satisfies (V) and (S), and hence no transported certificate exists on that band.
\end{proposition}
\begin{proof}
For any PSD field satisfying (S), the triangle inequality for the seminorm induced by $W(z_\epsilon)$ gives
\[
 \sqrt{\nabla F(z_\epsilon)^TW(z_\epsilon)\nabla F(z_\epsilon)}
 \le\sqrt C\sum_j|\lambda_{j,\epsilon}|\sqrt{t_{j,\epsilon}}\to0,
\]
contradicting (V). Neither transport nor boundedness or continuity of $W$ was used.
\end{proof}
In particular, the obstruction holds if $n$ incoming jumps converge to a basis of $\R^n$, their curvature products tend to zero, and endpoint gradients remain bounded. The coefficients expressing the gradient in that basis then remain bounded.

The predecessors need not belong to one BFGS history. This is not a loophole: the condition quantifies over all such pairs. Moreover, every eligible pair is individually realizable as a first BFGS step. If $d=-g^Ts>0$, then
\[
 H=\frac{ss^T}{d}+I-\frac{gg^T}{g^Tg}\succ0,\qquad Hg=-s.
\]
Hence $\alpha=1$ realizes it. Positive definiteness follows because a vector in both nullspaces would be parallel to $g$ and orthogonal to $s$, and therefore zero.


\setcounter{section}{3}
\setcounter{theorem}{0}
\setcounter{equation}{4}
\section{Connection with tangent and normal geometry}
At a nonsmooth point $p$, let $g_1,\dots,g_r$ be limiting active gradients. Define
\[
 N=\operatorname{span}\{g_i-g_j\},\qquad U=N^\perp.
\]
The projected gradients $P_Ug_i$ agree. They are the common part of the first-order models that is not affected by gradient jumps. In suitable partly smooth settings $N$ is the normal space and $U$ the tangent space; this interpretation and its relation to the affine span of the subdifferential are standard in partial smoothness~\cite{LL}.

For a certificate with a finite continuous extension to $p$, accessible incoming pairs with $t\to0$ force
\[
 W(p)(g_i-g_j)=0.
\]
Thus its range is tangential to the accessible jump directions. Visibility then requires a nonzero common tangential gradient. Algebraically,
\[
 P_Ug_i=0\quad\Longleftrightarrow\quad0\in\aff\{g_1,\dots,g_r\}.
\]
The condition $0\in\aff G$ is weaker than $0\in\conv G$: a point can have no tangent descent of this type yet remain nonstationary for the convex objective.

This limiting picture is an interpretation under the stated continuity and accessibility qualifications, \emph{not} a proof for arbitrary fields. Proposition~\ref{G:prop:star} supplies the rigorous obstruction without those regularity assumptions. Eligibility matters: some local derivative jumps cannot be crossed by arbitrarily small steps satisfying the fixed Wolfe condition.

For $u^2+|v|$ at $(u,0)$ with $u\ne0$, the jump space is vertical and a nonzero horizontal derivative remains. At a generic two-dimensional vertex with independent accessible jumps in both coordinates, there may be no remaining comparison direction. The common-target obstruction above makes that interpretation rigorous without continuity of a certificate.


\setcounter{section}{5}
\setcounter{theorem}{0}
\setcounter{equation}{5}
\section{In one dimension the condition can be fully eliminated}
\begin{proposition}[Exact scalar characterization]
On a compact positive objective band in $\R$, the transported certificate exists if and only if, for some small-decrease threshold, eligible pairs satisfy
\begin{equation}
 \sup_{(x,z)}\frac{|F'(z)-F'(x)|}{|z-x|\,|F'(z)|^2}<\infty. \label{G:eq:scalar}
\end{equation}
Equivalently, the derivative is uniformly Lipschitz on these eligible pairs:
\[
 |F'(z)-F'(x)|\le A_\eta|z-x|.
\]
This is not a Lipschitz assertion on all pairs or across the minimizer.
\end{proposition}
\begin{proof}
In one dimension $\Pi_s=0$, so transport imposes no restriction. From (V),
$W(z)\ge m/|F'(z)|^2$. From (S), using $t=|y||s|$,
$W(z)|y|\le C|s|$. These imply \eqref{G:eq:scalar}. Conversely, choose
$W(z)=1/|F'(z)|^2$; visibility equals one and \eqref{G:eq:scalar} gives (S).

On a compact positive band, convexity gives $|F'(z)|\ge\eta/R$, where $R$ bounds distance to a fixed minimizer, and local boundedness gives $|F'(z)|\le G$. These bounds make \eqref{G:eq:scalar} equivalent to the displayed Lipschitz inequality.
\end{proof}
This reveals a genuine restriction: an eligible arbitrarily small crossing of a nonoptimal derivative jump rules out every certificate.

\subsection{Convergence itself does not need this scalar restriction}
\begin{proposition}[Scalar convergence, independently of BFGS updates]\label{G:prop:scalarconv}
Let $f:\R\to\R$ be finite and convex with unique minimizer and compact initial sublevel. Every infinite differentiable-iterate Armijo--weak-Wolfe run using any positive scalar metric converges to that minimizer.
\end{proposition}
\begin{proof}
Suppose the monotone objective gap tends to $L>0$. Convexity gives $|g_k|\ge L/R$ on the compact sublevel. In one dimension a positive-metric descent step obeys $d_k=|g_k||s_k|$. Armijo makes $\sum d_k<\infty$, hence $\sum|s_k|<\infty$. Thus the iterates converge to a nonminimizer. On a tail they lie on one side of the unique minimizer, so all gradients have the same sign. Weak Wolfe then gives $|g_{k+1}|\le c_2|g_k|$, contradicting the positive lower bound. Objective convergence and compactness with uniqueness give iterate convergence.
\end{proof}
Already $f(x)=\frac12(x+2)^2+|x|$ gives a finite strongly convex example with global convergence but no global certificate on a sublevel containing zero. For $x=\epsilon$, $z=-\epsilon/2$,
\[
 g_x=3+\epsilon,\quad g_z=1-\epsilon/2,\quad y=-2-3\epsilon/2,
\]
\[
 d=\frac92\epsilon+\frac32\epsilon^2,\quad
 \Delta=\frac72\epsilon+\frac38\epsilon^2,\quad
 t=3\epsilon+\frac94\epsilon^2.
\]
These pairs are strictly eligible for $(1/4,1/2)$ for small $\epsilon$, and their scalar quotient in \eqref{G:eq:scalar} diverges.


\setcounter{section}{6}
\setcounter{theorem}{0}
\setcounter{equation}{6}
\section{Local nonexistence despite global convergence}
The difference between a local certificate and actual convergence remains even after shrinking every neighborhood of the optimizer.

Define $a_j=4^{-j}$ and
\begin{equation}
 h(x)=\frac12x^2+3\sum_{j=1}^\infty a_j(x-a_j)_+. \label{G:eq:infinite}
\end{equation}
The series converges uniformly on compact sets: its absolute value is bounded by a constant times $\sum a_j$. Every summand is convex, and $h(x)\ge x^2/2$ with $h(0)=0$. Thus $h$ is finite, $1$-strongly convex, coercive, and uniquely minimized at zero.

At a knot $a=a_j$, the one-sided derivatives are $2a$ and $5a$. These formulas come from exact geometric sums, not a truncation. For example, when $a/4<x<a$,
\[
 h(x)=\frac12x^2+ax-\frac15a^2,
\]
while just above $a$,
\[
 h(x)=\frac12x^2+4ax-\frac{16}{5}a^2.
\]
Also $h(a)=13a^2/10$.

Fix any knot, then let $x=a+\epsilon$, $z=a-\epsilon$, with $\epsilon\downarrow0$. Exactly,
\[
 g_x=5a+\epsilon,\quad g_z=2a-\epsilon,\quad y=-3a-2\epsilon,
\]
\begin{align*}
 d&=10a\epsilon+2\epsilon^2,&
 \Delta&=7a\epsilon,&
 t&=6a\epsilon+4\epsilon^2.
\end{align*}
The pairs are strictly eligible for $(c_1,c_2)=(1/4,1/2)$ for small $\epsilon$. Yet
\[
 \frac{|y|}{|s|\,g_z^2}
 =\frac{3a+2\epsilon}{2\epsilon(2a-\epsilon)^2}\longrightarrow\infty.
\]
No certificate exists on a positive band containing this knot. Every neighborhood of zero contains such a knot. Fixing a knot first, then taking $\epsilon\to0$, is essential: deterioration of band constants as $\eta\downarrow0$ cannot repair the divergence on the fixed knot's band.

Nevertheless Proposition~\ref{G:prop:scalarconv} proves convergence of every admissible run on $h$. The minimizer is differentiable for this scalar example, but the objective is nonsmooth in every neighborhood of it. Thus the transported condition is not necessary even for local convergence. A further source example is retained in the provenance archive; it is not needed for the nonnecessity conclusion here.


\setcounter{section}{8}
\setcounter{theorem}{0}
\setcounter{equation}{8}
\section{Why numerical nonexistence is intrinsically delicate}
\begin{proposition}[Every finite sample is feasible if constants are free]
Take any finite set of differentiable nonstationary endpoints and any finite set of eligible pairs between them. There are finite constants $m>0$, $C<\infty$ and a constant field $W(x)=I$ satisfying (V), (S), and (T0) on that finite sample.
\end{proposition}
\begin{proof}
Choose $m=\min_x\norm{\nabla F(x)}^2>0$ and
\[
 C=\max_{(x,z)}\frac{\norm{\nabla F(z)-\nabla F(x)}^2}{t(x,z)}<\infty.
\]
There are finitely many pairs and all their denominators are positive. Transport is equality for a constant field.
\end{proof}
Consequently a finite-sample feasibility problem with unrestricted constants cannot prove nonexistence. One needs a known bound on the constants, extra restrictions on the field, or an analytic limiting family that forces the required constants to infinity. The examples above supply the latter. Conversely, passing many finite tests is not an existence theorem for a uniform field on the continuum.


\endgroup

% END original chapters/G.tex

\StartModule{A}{Affine criticality and local geometry}{Source-derived geometric deductions from [SA]. Equal-rate descent, rank stability and finite semialgebraic affine-critical values; no blanket dynamical theorem is inferred.}

% BEGIN original chapters/affine_geometry.tex
\section{Equal-rate descent and affine criticality}
This module preserves the geometric propositions of the affine-criticality note~\cite{SA}. They do not depend on the BFGS principal-minor lemma. The originally proposed dynamical step is now proved only in the completed settings of Modules C and U; no blanket theorem for all finite smooth maxima is asserted.

For a finite convex $F$, put
\[
 V(p)=\operatorname{span}(\partial F(p)-\partial F(p)),\quad U(p)=V(p)^\perp,
 \quad a(p)=\operatorname{proj}_{\operatorname{aff}\partial F(p)}0.
\]
The affine hull equals $a(p)+V(p)$ and $a(p)\perp V(p)$. Hence every subgradient satisfies
\[
 P_{U(p)}g=a(p),\qquad g^Ta(p)=\|a(p)\|^2.
\]
The condition
\begin{equation}\label{A:AC}
 0\in\operatorname{aff}\partial F(p)\Longrightarrow0\in\partial F(p)
\end{equation}
is therefore equivalent, at every nonoptimal point under consideration, to the existence of a direction $v$ with $g^Tv=-1$ for every $g\in\partial F(p)$. One may take $v=-a(p)/\|a(p)\|^2$. Conversely such an equal-rate direction rules out zero in the affine hull. This is stronger than ordinary strict descent, which requires only negative, not equal, inner products.

On a compact admissible run, all cluster points share the limiting value $\ell$. If at least one cluster point is affine critical and \eqref{A:AC} holds on the basin, that point is optimal, so $\ell=F_*$. Compactness gives convergence to the solution set, and uniqueness gives convergence of the entire sequence. This reduction does not assume an iterate limit. It would be incorrect to replace affine criticality of a cluster point by vanishing gradients at differentiable iterates: for example, the gradients of $|x|$ have unit magnitude away from its minimizer.

\section{Why pointwise affine separation need not be uniform}
Let $\gamma(p)=\operatorname{dist}(0,\operatorname{aff}\partial F(p))$. A useful quantitative condition is a positive minimum of $\gamma$ on every compact positive objective band. Pointwise nonvanishing plus compactness does not establish it when ranks change.

Consider
\[
 F(u,v)=u+u^2+v^2+|uv|
 =\max\{u+u^2+v^2+uv,\ u+u^2+v^2-uv\}.
\]
Both branch Hessians have eigenvalues $1$ and $3$, so $F$ is strongly convex, with unique minimizer $(-1/2,0)$ and value $-1/4$. At $(0,t)$, $t\ne0$, the active gradients are $(1+t,2t)$ and $(1-t,2t)$, giving $\gamma(0,t)=2|t|$. At zero the gradient is $(1,0)$, so $\gamma(0,0)=1$. On the other switching axis $\gamma(u,0)=|1+2u|$, vanishing only at the optimizer. Smooth points are affine critical only when stationary. Thus \eqref{A:AC} holds globally, but $\gamma(0,t)\to0$ while the objective gap tends to $1/4$.

This is a geometric rank-loss example, not a BFGS counterexample. A sufficiently small neighborhood of its actual optimizer excludes the rank-loss point.

\section{Nondegenerate smooth maxima give local affine-critical correctness}
Let $F=\max_{i\le m}f_i$ with convex $C^1$ branches and a minimizer $x_*$. Restrict to $I_*$, the active branches there; inactive branches have positive gaps and disappear in a sufficiently small neighborhood. Suppose the active gradients are affinely independent and admit strictly positive coefficients
\[
 \lambda_i^*>0,\quad\sum_{i\in I_*}\lambda_i^*=1,
 \quad\sum_{i\in I_*}\lambda_i^*\nabla f_i(x_*)=0.
\]
Then \eqref{A:AC} holds near $x_*$.

\begin{proof}
Let $A(x)$ have columns $(\nabla f_i(x),1)$ for $i\in I_*$. Full column rank at $x_*$ persists nearby. At an affine-critical point $p$, extend the affine coefficients of its active gradients by zero on missing indices. They solve
\[
 A(p)\lambda=(0,\ldots,0,1)^T.
\]
Whenever this system is solvable, its unique solution equals
\[
 \lambda(p)=(A(p)^TA(p))^{-1}A(p)^T(0,\ldots,0,1)^T.
\]
This continuous expression equals the strictly positive vector $\lambda^*$ at $x_*$. It remains positive nearby. No missing index can therefore have zero coefficient, and the affine combination is a convex combination. Thus $0\in\partial F(p)$.
\end{proof}
This is the source note's barycentric-stability argument, closely related to the geometric nondegeneracy results it cites. It is not itself a BFGS convergence theorem.

There is also local lower semicontinuity of $\gamma$. For each nonempty $J\subset I_*$, the distance to the affine hull of the gradients in $J$ is continuous, because the relevant difference Gram matrix remains nonsingular. For $p_j\to p$, pass to a subsequence of constant active set $J\subset I(p)$. Then
\[
 \lim\gamma(p_j)=\operatorname{dist}(0,\operatorname{aff}\{\nabla f_i(p):i\in J\})\ge\gamma(p).
\]
On a compact positive band, local \eqref{A:AC} therefore gives a strictly positive minimum. This establishes visibility, not the full transport needed in arbitrary dimension.

\section{Semialgebraic objectives have finitely many affine-critical values}
\begin{proposition}[A low-level affine-critical gap]\label{A:semialgebraic}
If $F:\mathbb R^n\to\mathbb R$ is finite convex and semialgebraic and attains $F_*$, then some $\epsilon>0$ satisfies
\[
 F(p)<F_*+\epsilon,\quad0\in\operatorname{aff}\partial F(p)
 \quad\Longrightarrow\quad0\in\partial F(p).
\]
\end{proposition}
\begin{proof}
The subgradient graph is semialgebraic by its quantified supporting-plane description and quantifier elimination. Using at most $n+1$ subgradients in an affine representation shows that the affine-critical set $C_{\rm aff}$ is semialgebraic. It has finitely many semialgebraically path-connected components.

On a piecewise $C^1$ semialgebraic path $p(t)$ in that set, the composition $F\circ p$ is continuous and piecewise $C^1$. At a time where both derivatives exist, the two-sided convex subgradient inequality implies
\[
 (F\circ p)'(t)=g^Tp'(t)\qquad\text{for every }g\in\partial F(p(t)).
\]
Taking an affine combination of these subgradients equal to zero proves $(F\circ p)'(t)=0$. Thus $F$ is constant on each path component. There are finitely many affine-critical values. Since $F_*$ is one of them, choose $\epsilon$ below the least strictly larger value, or any positive number if no larger value occurs. Convexity turns equality with $F_*$ into stationarity.
\end{proof}
The standard semialgebraic path and projection facts are the external geometric ingredients cited in~\cite{SA}, including Bolte--Daniilidis--Lewis--Shiota~\cite{BDLS}. This proposition concerns affine criticality rather than merely Clarke criticality. It does not imply a blanket local BFGS theorem without a dynamical result forcing such a cluster point. Modules C and U provide completed restricted cases.

\section{What remains a conjecture}
The recorded general-dimensional conjecture concerns finite convex $C^{2,1}$ maxima with unique minimizer, compact low sublevels, affinely independent active gradients, positive optimal weights, and positive weighted Hessian on
\[
 U_* =\{v:(\nabla f_i(x_*)-\nabla f_1(x_*))^Tv=0\ \forall i\in I_*\}.
\]
The proposed conclusion is local full-BFGS convergence for arbitrary SPD initialization. The planar theorem proves a strict-contour version; the fixed-normal theorem proves a different arbitrary-dimensional class. Neither establishes this full conjecture. Straightening an active manifold by a nonlinear coordinate change does not preserve ordinary BFGS, and cannot supply the missing implication.

% END original chapters/affine_geometry.tex

\StartModule{C}{Forced affine-critical cluster points in the plane}{Complete retained proof from [SC], excluding its repeated matrix appendix. Strict positive contours and regular switching are substantive assumptions.}

% BEGIN original chapters/C.tex
\begingroup
\providecommand{\R}{}
\renewcommand{\R}{\mathbb R}
\providecommand{\aff}{}
\renewcommand{\aff}{\operatorname{aff}}
\providecommand{\conv}{}
\renewcommand{\conv}{\operatorname{conv}}
\providecommand{\tr}{}
\renewcommand{\tr}{\operatorname{tr}}
\providecommand{\dist}{}
\renewcommand{\dist}{\operatorname{dist}}
\providecommand{\Ca}{}
\renewcommand{\Ca}{\mathcal C_{\rm aff}}
\providecommand{\ip}{}
\renewcommand{\ip}[2]{\langle #1,#2\rangle}
\providecommand{\norm}{}
\renewcommand{\norm}[1]{\lVert #1\rVert}
\makeatletter\let\c@lemma\c@theorem\let\c@corollary\c@theorem\let\c@proposition\c@theorem\let\c@remark\c@theorem\makeatother
\renewcommand{\thelemma}{\thetheorem}\renewcommand{\thecorollary}{\thetheorem}\renewcommand{\theproposition}{\thetheorem}\renewcommand{\theremark}{\thetheorem}
\setcounter{section}{0}
\setcounter{theorem}{0}
\setcounter{equation}{0}
\section{The new statements}
Let $F:\R^2\to\R$ be finite and convex, attain $F_*$, and have a compact sublevel
\[
K=\{x:F(x)\le b\},\qquad b>F_*.
\]
A point is \emph{affine critical} when
\[
 p\in\Ca:=\{p:0\in\aff\partial F(p)\}.
\]
Ordinary stationarity, $0\in\partial F(p)$, implies affine criticality, but not conversely.

We use two objective-level hypotheses.

\textbf{Strictly convex positive sublevels.} For every $F_*<t\le b$, the convex body $K_t=\{F\le t\}$ is strictly convex: its boundary contains no nontrivial line segment. This is a condition on the contours, not a uniform lower Hessian bound. Strong convexity of $F$ suffices.

\textbf{Regular switching away from affine criticality.} At every $p\in K\setminus\Ca$, either $F$ is convex $C^2$ on a neighborhood of $p$, or on such a neighborhood
\[
 F=\max\{f_1,f_2\},
\]
where $f_1,f_2$ are convex $C^2$, both are active at $p$, and $\nabla f_1(p),\nabla f_2(p)$ are linearly independent. In particular, two different smooth branches meet transversally both to each other and to the positive level. No assumptions are made on the local structure at an affine-critical point.

\begin{theorem}[Forced affine-critical cluster point]\label{C:thm:cluster}
Under these hypotheses, every infinite admissible full-BFGS run starting in $K$ has at least one affine-critical cluster point. No iterate limit, bounded condition number, limiting eigenspace, matrix accuracy at initialization, or finite path length is assumed.
\end{theorem}

Here admissibility means $H_0\succ0$, differentiable accepted iterates with nonzero gradients, $s_k=-\alpha_kH_kg_k$, $\alpha_k>0$, and fixed Armijo--weak-Wolfe conditions
\begin{equation}\label{C:eq:AW}
 F(x_{k+1})\le F(x_k)+c_1g_k^Ts_k,\qquad
 g_{k+1}^Ts_k\ge c_2g_k^Ts_k,\qquad 0<c_1<c_2<1.
\end{equation}
The usual full inverse update is
\begin{equation}\label{C:eq:bfgs}
 H_{k+1}=\left(I-\frac{s_ky_k^T}{t_k}\right)H_k
 \left(I-\frac{y_ks_k^T}{t_k}\right)+\frac{s_ks_k^T}{t_k},\quad
 y_k=g_{k+1}-g_k,\quad t_k=y_k^Ts_k>0.
\end{equation}
Existence of differentiable acceptable lengths on every exceptional search ray is not part of the general theorem.

\begin{corollary}[No false affine criticality gives convergence]\label{C:cor:AC}
If, in addition, $\Ca\cap K\subseteq\arg\min F$, then every such run satisfies $F(x_k)\to F_*$ and $\dist(x_k,\arg\min F)\to0$. A unique minimizer gives convergence of the entire iterate sequence. A local version holds when the hypotheses hold only on a sufficiently low compact sublevel inside a structural neighborhood of a unique minimizer.
\end{corollary}

\begin{corollary}[Nondegenerate planar smooth maxima]\label{C:cor:max}
Suppose a finite convex $F:\R^2\to\R$ has a unique minimizer $x_*$ and, near $x_*$,
\[
 F=\max_{i\le m}f_i,
\]
with convex $C^2$ branches. Suppose sufficiently low positive sublevels are strictly convex. Let $I_*$ be the active set at $x_*$, and assume its gradients are affinely independent and admit strictly positive optimal weights:
\begin{equation}\label{C:eq:nondeg}
 \lambda_i^*>0,\quad \sum_{i\in I_*}\lambda_i^*=1,\quad
 \sum_{i\in I_*}\lambda_i^*\nabla f_i(x_*)=0.
\end{equation}
Then some $\varepsilon>0$ satisfies
\[
 F(x_0)<F_*+\varepsilon\quad\Longrightarrow\quad x_k\to x_*
\]
for every infinite admissible full-BFGS run, with arbitrary $H_0\succ0$ and any fixed valid line-search constants. The basin can be chosen independently of these algorithmic choices. A simple sufficient replacement for the strictly convex-sublevel hypothesis is a common positive lower Hessian bound on the finitely many locally relevant branches.
\end{corollary}

This proves a planar, strictly-convex-contour version of the smooth-max conjecture recorded in the previous note~\cite{SA}. It does not prove its general-dimensional or merely weighted-tangential-curvature version.


\setcounter{section}{1}
\setcounter{theorem}{3}
\setcounter{equation}{3}
\section{Why one affine-critical cluster point is enough}
Armijo makes the objective values decrease to a limit $\ell\ge F_*$. Every cluster point has value $\ell$. If any cluster point is stationary, then $\ell=F_*$. Compactness gives distance to the solution set tending to zero; uniqueness gives whole-sequence convergence.

The geometric condition
\begin{equation}\label{C:eq:AC}
 0\in\aff\partial F(p)\ \Longrightarrow\ 0\in\partial F(p)
\end{equation}
therefore reduces the problem to Theorem~\ref{C:thm:cluster}. This reduction was proved in~\cite{SA}; the missing algorithmic assertion is the subject of this note.

For completeness, if $a(p)$ is the minimum-norm point of $\aff\partial F(p)$, then
\[
 g^Ta(p)=\norm{a(p)}^2\qquad(g\in\partial F(p)).
\]
At a non-affine-critical point, $-a(p)/\norm{a(p)}^2$ decreases every supporting model at rate one. Thus \eqref{C:eq:AC} says that every nonoptimal point retains equal-rate descent across its active models. It does not say that the gradients at differentiable accepted points tend to zero.


\setcounter{section}{2}
\setcounter{theorem}{3}
\setcounter{equation}{4}
\section{Vanishing steps without an assumed iterate limit}
Put
\[
 q_k=g_k^TH_kg_k,\quad d_k=-g_k^Ts_k=\alpha_kq_k,\quad
 \Delta_k=F(x_k)-F(x_{k+1}),\quad c=1-c_2>0.
\]
Convexity and \eqref{C:eq:AW} give
\begin{equation}\label{C:eq:eligible}
 d_k>0,\qquad c_1d_k\le\Delta_k\le d_k,\qquad t_k\ge c d_k.
\end{equation}
In particular $\sum_k d_k<\infty$. A finite convex function has bounded subgradients on a compact set: compare supporting planes at $x\in K$ with values on a fixed enlarged compact neighborhood. Hence the accepted gradients are bounded, and Module~\ref{M:chapter} applies.

\begin{lemma}[Strict contours force vanishing steps]\label{C:lem:steps}
If $F(x_k)\downarrow\ell>F_*$ and $K_\ell$ is strictly convex, then $s_k\to0$.
\end{lemma}
\begin{proof}
Otherwise a subsequence has $x_k\to x$, $x_{k+1}\to z\ne x$, and $g_k\to g\in\partial F(x)$ by compactness and closedness of the convex subgradient graph. Since $d_k\to0$,
\[
 F(x)=F(z)=\ell,\qquad g^T(z-x)=0.
\]
The supporting inequality at $x$ and convexity make $F$ identically $\ell$ on $[x,z]$. This is a nontrivial boundary segment of $K_\ell$, a contradiction.
\end{proof}

We will argue by contradiction that the compact cluster set $\Gamma$ contains no affine-critical point. Then $\ell>F_*$, every point of $\Gamma$ satisfies the regular-switching hypothesis, and Lemma~\ref{C:lem:steps} applies. This does not assume that $\Gamma$ is a singleton, or even use its connectedness.


\setcounter{section}{3}
\setcounter{theorem}{4}
\setcounter{equation}{5}
\section{Joining comparison directions around a regular positive level}
This is the new geometric construction. Fix a compact subset $C\subset\{F=\ell\}$ containing no affine-critical point and satisfying the regularity hypothesis. Eventually we use $C=\Gamma$.

\subsection{Finitely many relevant corners and their level velocities}
Near a nonsmooth point of $C$, $F=\max\{f_1,f_2\}$ and the matrix with rows $\nabla f_1^T,\nabla f_2^T$ is nonsingular. The equations
\[
 f_1(p)=f_2(p)=t
\]
therefore define a $C^2$ point curve $p_j(t)$ near $t=\ell$. On the level $\ell$, the switching point is isolated. Smooth points of $C$ have smooth neighborhoods. A compactness argument gives finitely many relevant curves $p_j(t)$ and a neighborhood of $C$ whose nonsmooth points on nearby levels lie on these curves. Explicitly, an infinite sequence of distinct level-$\ell$ corners in a compact regular neighborhood would accumulate at a smooth chart, or at a chart with only one solution of $f_1=f_2=\ell$; either is impossible. Shrinking the neighborhood and the level interval removes all other switching curves.

Define
\begin{equation}\label{C:eq:vel}
 v_j(t)=p_j'(t),\qquad
 \nabla f_{j,1}(p_j(t))^Tv_j(t)
 =\nabla f_{j,2}(p_j(t))^Tv_j(t)=1.
\end{equation}
This vector is tangent to the switching curve and detects a fixed nonzero increase in the objective. It is $C^1$ in $t$.

\subsection{Why strictly convex contours permit global joining}
At $p_j(t)$, the nonnegative combinations of its two active gradients generate a closed arc of unit supporting normals of $K_t$. This arc has width less than $\pi$, and every unit normal $n$ in it satisfies $n^Tv_j(t)>0$.

The normal arcs of distinct corners at $t=\ell$ are disjoint. Indeed, a shared normal would support $K_\ell$ at two different points, making their connecting segment lie on its boundary. Strict convexity prohibits this. The finitely many compact arcs consequently have positive angular separation. After slightly enlarging them and shrinking the interval of $t$, choose \emph{fixed} smooth angular cutoffs $\chi_j$ with disjoint supports, equal to one on neighborhoods of the entire moving normal arcs. The supports can be chosen where $n^Tv_j(t)>0$ uniformly.

Let
\[
 n(\theta)=(\cos\theta,\sin\theta),\qquad
 T(\theta)=(-\sin\theta,\cos\theta),
\]
and put
\begin{equation}\label{C:eq:h}
 h(t,\theta)=1+\sum_j\chi_j(\theta)\,[n(\theta)^Tv_j(t)-1].
\end{equation}
The supports are disjoint, so this is a convex combination of positive numbers at each angle. Hence $h$ is bounded above and bounded away from zero. It is $C^1$ in $t$ and smooth in $\theta$, with the mixed derivatives used below bounded. No convex-support-function property of $h$ is asserted or needed.

At a smooth point near $C$, write $n(\theta(x))=\nabla F(x)/\norm{\nabla F(x)}$ and define
\begin{equation}\label{C:eq:w}
 w(x)=h(F(x),\theta(x))n(\theta(x))
       +h_\theta(F(x),\theta(x))T(\theta(x)).
\end{equation}
These expressions are periodic in $\theta$, so no choice of a global angle coordinate is needed.

Where $\chi_j=1$, $h=n^Tv_j$ and $h_\theta=T^Tv_j$, so
\begin{equation}\label{C:eq:nearcorner}
 w(x)=v_j(F(x)).
\end{equation}
This holds on a whole neighborhood of the corresponding switching curve, on both sides. Define $w$ by \eqref{C:eq:nearcorner} on the curve itself. The field thus extends continuously and locally Lipschitzly across the gradient jump. Away from the curves it is $C^1$. On a smaller compact neighborhood of $C$, all close-pair Lipschitz and piecewise derivative bounds are uniform.

At smooth points,
\begin{equation}\label{C:eq:visibility}
 \nabla F(x)^Tw(x)=\norm{\nabla F(x)}h(F(x),\theta(x))>0.
\end{equation}
At a switching point, every subgradient has inner product one with $w$. Compactness on a positive level neighborhood gives a uniform positive lower bound in \eqref{C:eq:visibility} and uniform bounds on $w$.

\subsection{The planar property that makes transport inexpensive}
At fixed $t$, differentiation gives the exact identity
\begin{equation}\label{C:eq:angleidentity}
 \partial_\theta(hn+h_\theta T)=(h+h_{\theta\theta})T.
\end{equation}
Therefore, at a smooth point, $Dw(x)$ maps the one-dimensional tangent line of the level curve into itself. It may vary arbitrarily in the normal direction. In the orthonormal basis $(T,n)$, it has the form
\begin{equation}\label{C:eq:triangular}
 Dw=\begin{pmatrix} a&b\\0&e\end{pmatrix}.
\end{equation}
The coefficients are bounded in the compact neighborhood. The sign of $h+h_{\theta\theta}$ is immaterial.


\setcounter{section}{4}
\setcounter{theorem}{4}
\setcounter{equation}{12}
\section{Three endpoint estimates}
An eligible pair means differentiable $x,z$ satisfying \eqref{C:eq:eligible}, with $s=z-x$, $g=\nabla F(x)$, $y=\nabla F(z)-g$, and $t=y^Ts$. The following estimates hold for sufficiently close eligible pairs in a buffered neighborhood of $C$. All intermediate segment points lie in the larger structural neighborhood.

\begin{lemma}[Uniform transversality]\label{C:lem:trans}
There is $\kappa>0$ such that
\[
 |\det(s,w(x+\theta s))|\ge\kappa\norm{s}
 \quad(0\le\theta\le1)
\]
for all sufficiently close eligible pairs.
\end{lemma}
\begin{proof}
Otherwise take pairs coalescing to $p$ in the compact neighborhood, unit steps $s/\norm{s}\to e$, and segment points where the normalized determinant tends to zero. Continuity of $w$ implies $e$ is parallel to $w(p)$. Subgradient limits $a,b\in\partial F(p)$ have the same strictly positive inner product with $w(p)$: this is automatic at a smooth point and follows from \eqref{C:eq:vel} at a corner. Descent selects $e=-w(p)/\norm{w(p)}$, so $d/\norm{s}\to-a^Te>0$. But
\[
 t/\norm{s}\to(b-a)^Te=0,
\]
contradicting $t\ge(1-c_2)d$. Compactness makes the bound uniform, including all segment points.
\end{proof}

\begin{lemma}[Transport is charged to objective decrease]\label{C:lem:transport}
For the same pairs, there are scalars $\lambda>0,\beta$ satisfying
\begin{equation}\label{C:eq:transport}
 w(z)=\lambda w(x)+\beta s,\qquad
 |\log\lambda|\le C d.
\end{equation}
In particular $\lambda^2\le1+C'd$ for sufficiently small $d$.
\end{lemma}
\begin{proof}
Let $q(\theta)=x+\theta s$, $\varphi(\theta)=F(q(\theta))$, and $D(\theta)=\det(s,w(q(\theta)))$. Lemma~\ref{C:lem:trans} makes $D$ nonzero with constant sign. Both $w\circ q$ and $\varphi$ are absolutely continuous.

On a smooth part of the segment use \eqref{C:eq:triangular}. With $s=s_TT+s_n n$,
\[
 \det(s,Dw\,s)=s_n[(e-a)s_T-bs_n].
\]
The uniform derivative bound, transversality, and a positive lower bound on $\norm{\nabla F}$ yield
\[
 |(\log|D|)'|\le C|\nabla F(q)^Ts|=C|\varphi'|.
\]
Inside a corner neighborhood, \eqref{C:eq:nearcorner} gives directly
\[
 (w\circ q)'=v_j'(\varphi)\varphi'\quad\hbox{a.e.},
\]
which gives the same bound even on the nondifferentiable subset. This argument does not assume finitely many intersections of a straight segment with a $C^2$ switching curve.

Convexity of $\varphi$, $\varphi'(0)=-d$, and $\varphi(1)\le\varphi(0)$ imply
\begin{equation}\label{C:eq:variation}
 \int_0^1|\varphi'(\theta)|\,d\theta\le2d.
\end{equation}
Indeed the negative variation is at most $d$, because every later derivative is at least $-d$; the positive variation is no greater than the negative variation. Integration proves $|\log(D(1)/D(0))|\le2Cd$. Since $s,w(x)$ form a basis, the coefficient of $w(x)$ in $w(z)$ is exactly $\lambda=D(1)/D(0)>0$. This proves the claim.
\end{proof}

\begin{lemma}[Bounded secant correction]\label{C:lem:secant}
There is a uniform $C_0<\infty$ such that
\[
 \frac{(y^Tw(z))^2}{t}\le C_0
\]
for all sufficiently close eligible pairs in the same neighborhood.
\end{lemma}
\begin{proof}
Use a finite buffered chart cover. For close pairs in a common smooth chart, or with the same active branch in a two-branch chart, convex $C^2$ smoothness gives $\norm{y}^2\le L t$ by Hessian integration. Bounded $w$ proves the claim.

For different active branches $i,j$ in a two-branch chart, the segment has a tie point $p$. At that point the jump is orthogonal to $w(p)$. The branch gradients and $w$ are Lipschitz, so
\begin{equation}\label{C:eq:numerator}
 |y^Tw(z)|\le C_1\norm{s}.
\end{equation}
For the denominator put $a=\nabla f_i(x)$, $b=\nabla f_j(x)$, and $r=\nabla f_j(z)-\nabla f_j(x)$, with $\norm{r}\le L\norm{s}$. Then
\[
 a^Ts=-d,\qquad b^Ts=t-d-r^Ts.
\]
The two gradient rows are uniformly independent on the chart. Thus
\[
 \norm{s}\le A(2d+t+L\norm{s}^2)
 \le A\left(1+\frac{2}{1-c_2}\right)t+AL\norm{s}^2.
\]
For small $s$, absorb the final term to obtain $t\ge c_0\norm{s}$. Combining with \eqref{C:eq:numerator} bounds the correction by a constant times $\norm{s}$, hence uniformly. This is a crossing-specific linear curvature estimate, not a uniform positive Hessian assumption.
\end{proof}


\setcounter{section}{5}
\setcounter{theorem}{7}
\setcounter{equation}{15}
\section{The BFGS contradiction}
\begin{proof}[Proof of Theorem~\ref{C:thm:cluster}]
Suppose the compact cluster set $\Gamma$ misses $\Ca$. Then the limiting value $\ell$ is strictly above $F_*$; otherwise all cluster points are minimizers. Lemma~\ref{C:lem:steps} gives $s_k\to0$. The distance of $x_k$ to $\Gamma$ tends to zero, so the entire tail lies in the buffered regular neighborhood constructed above with $C=\Gamma$.

Let $B_k=H_k^{-1}$ and
\[
 E_k=w(x_k)^TB_kw(x_k),\qquad
 M_k=B_k-\frac{g_kg_k^T}{q_k}\succeq0,\qquad M_ks_k=0.
\]
The Hessian form of \eqref{C:eq:bfgs}, the representation in Lemma~\ref{C:lem:transport}, and $M_ks_k=0$ give the exact identity
\begin{equation}\label{C:eq:energy}
 E_{k+1}=\lambda_k^2\left(E_k-\frac{(g_k^Tw(x_k))^2}{q_k}\right)
       +\frac{(y_k^Tw(x_{k+1}))^2}{t_k}.
\end{equation}
The parenthesized expression is nonnegative. Uniform visibility and the last two lemmas imply, on a tail,
\[
 E_{k+1}\le (1+C d_k)(E_k-m/q_k)+C_0,
 \qquad m>0.
\]
Set $P_K=1$ and $P_{k+1}=P_k(1+Cd_k)$. Since $\sum d_k<\infty$, $1\le P_k\le\bar P<\infty$. Divide by $P_{k+1}$ and telescope. Positivity gives
\[
 \frac{m}{\bar P}\sum_{k=K}^{N-1}\frac1{q_k}
 \le E_K+C_0(N-K).
\]
This contradicts Module~\ref{M:chapter}, which forces the average of $1/q_k$ to diverge. Thus $\Gamma\cap\Ca\ne\varnothing$.
\end{proof}

The argument constructs a field near the \emph{whole hypothetical cluster set}. It does not first select one cluster point and assume all later iterates stay near it. Nor is the field an extra algorithmic input: it is an analytical object derived from the objective, independent of every $H_k$.


\setcounter{section}{6}
\setcounter{theorem}{7}
\setcounter{equation}{16}
\section{Local consequences and natural assumptions}
\begin{proof}[Proof of Corollary~\ref{C:cor:AC}]
All cluster points have the limiting objective value. An affine-critical cluster point is optimal by the extra hypothesis; hence that value is $F_*$. Compactness gives distance to the solution set tending to zero. Under uniqueness every convergent subsequence has the same limit, proving whole-sequence convergence. For the local statement, compactness and uniqueness put a sufficiently low sublevel inside any given neighborhood of $x_*$; Armijo then traps the accepted iterates there.
\end{proof}

\begin{proof}[Proof of Corollary~\ref{C:cor:max}]
Remove branches inactive at $x_*$ by shrinking the neighborhood. Form the augmented matrix $A(x)$ whose columns are $(\nabla f_i(x),1)$ for $i\in I_*$. Affine independence makes $A(x_*)$ full column rank; that rank persists. Whenever $A(p)\lambda=(0,0,1)^T$ has a solution, it is unique and equals
\[
 (A(p)^TA(p))^{-1}A(p)^T(0,0,1)^T.
\]
This is continuous and agrees with the positive vector $\lambda^*$ at $x_*$. Shrink the neighborhood so every component remains positive. An affine representation of zero using the gradients active at $p$, extended by zero coefficients on inactive indices, must equal this vector. It therefore has no inactive indices and is a convex representation of zero. This proves local \eqref{C:eq:AC}. This geometric argument is the one recorded in~\cite{SA}, with related nondegeneracy arguments in~\cite{BG}.

In the plane at most three gradients are affinely independent. At a non-affine-critical nearby point, three cannot all be active, since their affine hull is $\R^2$. A singleton gives a smooth chart. For two active gradients, the difference is nonzero; if the gradients were linearly dependent, their affine line would contain zero, again contradicting non-affine-criticality. The other branches have strict gaps at that point, so the required local two-branch representation follows.

Choose a low sublevel contained in this neighborhood. Its strictly convex positive sublevels and the regular-switching property allow Theorem~\ref{C:thm:cluster}, then Corollary~\ref{C:cor:AC}. The localization uses only the objective, not the line-search constants or the initial matrix. Uniformly strongly convex branches on a convex local neighborhood make their maximum strongly convex there, supplying the sufficient strict-contour condition.
\end{proof}

\subsection{A tame-geometric alternative}
The previous note proved that a finite convex semialgebraic function has only finitely many affine-critical values~\cite{SA}. For clarity, its subgradient graph and affine-critical set are semialgebraic. The latter has finitely many piecewise-smooth path components. On such a path the two-sided subgradient inequality makes the derivative of $F$ equal to $g^T\dot x$ for \emph{every} subgradient $g$. An affine combination giving zero forces that derivative to vanish. Thus $F$ is constant on each component; choose a low level excluding the finitely many larger values. This is a geometric consequence consistent with the stratification framework of~\cite{BDLS}, not a BFGS theorem from that source.

Combining that fact with Theorem~\ref{C:thm:cluster} gives a second local sufficient class: \emph{planar finite convex semialgebraic objectives with strictly convex low sublevels and the stated regular two-branch structure at non-affine-critical points}. No positive-multiplier condition at the minimizer is needed if local \eqref{C:eq:AC} is obtained this way. This does not cover arbitrary semialgebraic functions: degenerate noncritical ties and flat contours remain outside this theorem.


\setcounter{section}{7}
\setcounter{theorem}{7}
\setcounter{equation}{16}
\section{An explicit new quadratic example and basin}
Consider
\begin{align*}
 f_0(u,v)&=u+u^2+\tfrac12v^2,\\
 f_1(u,v)&=v+\tfrac12u^2+\tfrac32v^2,\\
 f_2(u,v)&=-u-v+\tfrac32u^2+v^2,\qquad F=\max_i f_i.
\end{align*}
All branch Hessians lie between $I$ and $3I$. Also $(f_0+f_1+f_2)/3=u^2+v^2$, so $F\ge u^2+v^2$ and zero is its unique minimizer. At zero the active gradients are $(1,0),(0,1),(-1,-1)$, with positive weights $1/3$.

The affine coefficient calculation is exact. With $D=3(1+v-uv)$, the solution of $\sum\lambda_i\nabla f_i=0$, $\sum\lambda_i=1$ is
\[
 \lambda_0=\frac{1-4u+3v-7uv}{D},\quad
 \lambda_1=\frac{1+2u-3v-uv}{D},\quad
 \lambda_2=\frac{1+2u+3v+5uv}{D}.
\]
On $|u|,|v|\le1/10$, $D\ge267/100$ and the numerators are at least $23/100$, $49/100$, and $45/100$. Hence affine independence and local \eqref{C:eq:AC} hold on this box. Since $F\ge\norm{x}^2$,
\begin{equation}\label{C:eq:explicitbasin}
 \boxed{F(x_0)<1/100\quad\Longrightarrow\quad x_k\to0}
\end{equation}
for every infinite admissible full-BFGS run with arbitrary SPD initialization.

\subsection{This is outside the earlier radial criterion}
For small $u>0$, the active tie $f_0=f_1$ has
\[
 v=v(u)=\frac{-1+\sqrt{1+4u+2u^2}}2
 =u-\tfrac12u^2+O(u^3),
\]
with $f_2$ inactive. The gradient jump at the tie is
\[
 \nabla f_0-\nabla f_1=(1+u,-1-2v),\qquad
 (\nabla f_0-\nabla f_1)^T(u,v)=\tfrac12u^2-v^2\ne0.
\]
Take $x=(u,v+\epsilon)$ on branch 1 and $z=(u,v-\epsilon)$ on branch 0. Then $s=(0,-2\epsilon)$ and
\[
 d=2\epsilon(1+3v+3\epsilon),\quad
 t=2\epsilon(1+2v+4\epsilon),\quad
 \Delta=\epsilon(1+4v)+\epsilon^2.
\]
For each sufficiently small fixed $u>0$ and $\epsilon\downarrow0$, these are eligible for $c_1=1/4,c_2=1/2$. The radial numerator tends to the nonzero quantity $\tfrac12u^2-v^2$ while $t\to0$. Thus RS fails on a positive band in every neighborhood of zero, even though \eqref{C:eq:explicitbasin} holds. Curved interfaces between arbitrary quadratic branches are not being replaced by finite polyhedral cells or by a nonlinear change of BFGS coordinates.


\setcounter{section}{8}
\setcounter{theorem}{7}
\setcounter{equation}{17}
\section{What can already be proved in any dimension}
There is a useful partial statement before transport. Suppose a compact regular region has a finite convex $C^2$ maximum representation. Assume that near each non-affine-critical point $p$, the gradient-difference matrix of the fixed family $I(p)$ has constant rank. Affine independence is an easy sufficient assumption.

The minimum-norm point $a_J(x)$ of the affine hull of these gradients depends smoothly on $x$ in a smaller chart, and remains nonzero. Hence
\[
 v_J(x)=a_J(x)/\norm{a_J(x)}^2,
 \qquad \nabla f_i(x)^Tv_J(x)=1\quad(i\in J).
\]
A finite partition of unity on buffered charts yields a bounded smooth field $v$ with $g^Tv=1$ at every active gradient. For sufficiently close endpoints, each chart contributing at $z$ also contains $x$. If $i$ is active at $x$, both $\nabla F(z)^Tv(z)$ and $\nabla f_i(z)^Tv(z)$ equal one. Therefore
\[
 y^Tv(z)=[\nabla f_i(z)-\nabla f_i(x)]^Tv(z).
\]
The local one-sided convex smooth estimate and the finite-max supporting gap give
\[
 \norm{\nabla f_i(z)-\nabla f_i(x)}^2
 \le2L D_{f_i}(z,x)\le2L y^Ts.
\]
To justify the local one-sided estimate, subtract the tangent plane at $x$ and apply the descent lemma at $z$ to the step $-(\nabla f_i(z)-\nabla f_i(x))/L$; for sufficiently close buffered endpoints this step stays in the larger convex chart. Thus visibility and the secant correction can be constructed without a BFGS matrix assumption.

\textbf{What does not follow:} smoothness of this field does not provide the projected transport needed by the BFGS energy. The planar proof succeeds by making its derivative along the level tangent parallel to that tangent. In dimension three and higher, mapping a tangent space into itself is not enough: it may rotate or shear the executed step within the tangent space, while $M_ks_k=0$ removes only the component parallel to that one step.

A concrete warning uses the entirely smooth quadratic $F(x)=x^TQx/2$, $Q=\operatorname{diag}(1,2,3)$, $x=(1,1,1)$, and the canonical common-rate field $v=\nabla F/\norm{\nabla F}^2$. Take
\[
 s=(\epsilon,0,-(\epsilon+\epsilon^2)/3),\quad y=Qs.
\]
Then
\[
 d=\epsilon^2,\quad t=\frac{\epsilon^2(\epsilon^2+2\epsilon+4)}3,\quad
 \det[g,s,y]=\frac43\epsilon^2(1+\epsilon)\ne0.
\]
For small positive $\epsilon$ the pair is Armijo--weak-Wolfe eligible with $c_1=1/4,c_2=1/2$, but $v(x+s)\notin\operatorname{span}\{v(x),s\}$. Its rank-one projected transport therefore fails for any finite multiplier. Another field, the radial field, works for this quadratic. This is not nonexistence of all certificates or a BFGS counterexample; it demonstrates the gap in an argument based only on a smoothly chosen equal-rate direction.


\endgroup

% END original chapters/C.tex

\StartModule{U}{Arbitrary dimension with fixed jump normals}{Selected objective proofs from [SU], using the new residual budget in Module M. Parallel interfaces and a single exceptional hyperplane are distinct statements.}

% BEGIN original chapters/U.tex
\begingroup

\setcounter{section}{0}
\setcounter{theorem}{0}
\setcounter{lemma}{0}
\setcounter{corollary}{0}
\setcounter{proposition}{0}
\setcounter{remark}{0}
\setcounter{equation}{0}
\section{A concrete dimension-independent local
theorem}\label{U:a-concrete-dimension-independent-local-theorem}

Let \(a\in\mathbb R^n\) be a unit vector, \(b\in\mathbb R\), and write
\[
 P=\{x:a^Tx=b\},\qquad S=a^\perp,\qquad P_S=I-aa^T.
\] The normal direction is fixed; the dimension \(n-1\) of the tangent
space is unrestricted.

\phantomsection\label{U:one}\textbf{Theorem U.1 (One exceptional hyperplane).} Let
\(F:\mathbb R^n\to\mathbb R\) be finite and convex, with unique
minimizer \(x_*\in P\). Let \(K=\{F\le F(x_0)\}\) be compact. On an open
convex neighborhood of \(K\), assume:

\begin{enumerate}
\def\labelenumi{\arabic{enumi}.}
\tightlist
\item
  \(F\) is \(\mu\)-strongly convex for some \(\mu>0\).
\item
  All subgradients have the same tangential projection,
  \[P_Sg=A(x)\qquad(g\in\partial F(x)),\] where \(A\) is Lipschitz
  continuous on that neighborhood.
\item
  \(F\) is locally \(C^{1,1}\) away from \(P\).
\end{enumerate}

Then every infinite admissible full-BFGS run starting at \(x_0\)
satisfies \[
 \boxed{F(x_k)\to F_*,\qquad x_k\to x_*.}
\] The hyperplane itself may include differentiable points; it is
allowed to contain failures of gradient-Lipschitz regularity as well as
actual kinks. No uniform bound on the full gradient-Lipschitz constants
as the hyperplane is approached is required.

If these hypotheses hold only in a neighborhood of a unique minimizer,
some \(\varepsilon>0\) gives the same conclusion for every start with
\(F(x_0)<F_*+\varepsilon\). The same basin works for every SPD
initialization and every fixed valid line-search pair. Compact sublevel
trapping supplies this local statement: choose a low sublevel contained
in the structural neighborhood, then use Armijo monotonicity.

The assumptions concern the objective. In particular, the claim that a
part of \(B_k=H_k^{-1}\) has a controlled growth rate will be
\textbf{derived}, not assumed.

\subsection{1.1 A simple family and a weaker-curvature
version}\label{U:a-simple-family-and-a-weaker-curvature-version}

An immediate family is \[
 \boxed{F(u,v)=\psi(u,v)+h(v),\qquad u\in\mathbb R^{n-1},\quad v\in\mathbb R,}
\] where \(\psi\) is convex with locally Lipschitz gradient, \(h\) is
finite convex and locally \(C^{1,1}\) on \(\mathbb R\setminus\{0\}\),
\(F\) is strongly convex on the relevant neighborhood, and the minimizer
has \(v_*=0\). Here \(A=\nabla_u\psi\), irrespective of the singular
normal derivative of \(h\). The objective and the initial matrix may be
fully coupled.

Full strong convexity can be weakened for this composite family. It is
enough that \[
 \psi(u,v)-\frac{\mu}{2}\|u\|^2\quad\hbox{is convex},\qquad \mu>0,
\] that the relevant positive sublevels of \(F\) are strictly convex,
and that the minimizer is unique with \(v_*=0\). A compact initial
sublevel is still required. The proof of this extension is given in
Section U.6. It needs positive curvature only in the smooth tangential
directions, not a positive lower bound in the normal direction.

This is not a claim that all convex nonsmooth functions admit such a
splitting. Nor does a common hyperplane through the optimizer describe
arbitrary curved switching geometry.


\setcounter{section}{1}
\setcounter{theorem}{0}
\setcounter{lemma}{0}
\setcounter{corollary}{0}
\setcounter{proposition}{0}
\setcounter{remark}{0}
\setcounter{equation}{0}
\section{A forced affine-critical cluster-point
theorem}\label{U:a-forced-affine-critical-cluster-point-theorem}

The method also fits the affine-criticality program directly.

\phantomsection\label{U:parallel}\textbf{Theorem U.2 (Parallel singular interfaces).} Let \(F\) be finite
convex with attained minimum and compact sublevel \(K\). Assume \(F\) is
\(\mu\)-strongly convex on an open convex neighborhood of \(K\). Let
\(a\) be a unit vector and \(E\subset\mathbb R\) closed. On that
neighborhood suppose:

\begin{itemize}
\tightlist
\item
  the singular set is \(\Sigma=\{x:a^Tx\in E\}\), every point of
  \(\Sigma\) is genuinely nondifferentiable, and \(F\) is locally
  \(C^{1,1}\) off \(\Sigma\);
\item
  \(P_S\partial F(x)=\{A(x)\}\) for a Lipschitz field \(A\), where
  \(S=a^\perp\).
\end{itemize}

Then the cluster set of every infinite admissible full-BFGS run in \(K\)
contains a point \(p\) with \[
 \boxed{0\in\operatorname{aff}\partial F(p).}
\] There may be several or countably many parallel interfaces, subject
to the closed-set and regularity assumptions above. No finite
affine-critical set is assumed.

At a point of \(\Sigma\), the subdifferential is a nontrivial compact
segment parallel to \(a\), because every subgradient has the same \(S\)
projection. Hence \[
 \operatorname{aff}\partial F(p)=A(p)+\mathbb R a,
 \qquad
 0\in\operatorname{aff}\partial F(p)\iff A(p)=0.
\] At a smooth point, affine criticality is ordinary stationarity.

Consequently, adding \[
 0\in\operatorname{aff}\partial F(p)\implies0\in\partial F(p)
\] on \(K\), or on a sufficiently low trapped neighborhood, gives
convergence to the minimum value and, under uniqueness, convergence of
the whole sequence.

Theorem U.1 is slightly different from a literal specialization of Theorem U.2: it allows the exceptional hyperplane to contain smooth but not
locally \(C^{1,1}\) points. Having \(x_*\in P\) makes every zero of
\(A\) on \(P\) optimal, so genuine nondifferentiability at every
exceptional point is unnecessary there.


\setcounter{section}{4}
\setcounter{theorem}{0}
\setcounter{lemma}{0}
\setcounter{corollary}{0}
\setcounter{proposition}{0}
\setcounter{remark}{0}
\setcounter{equation}{0}
\section{Proof of the parallel-interface and single-interface
results}\label{U:proof-of-the-parallel-interface-and-single-interface-results}

We give a common argument that covers Theorems 1 and 2.

\subsection{5.1 Preliminary bounds and a contradiction
set}\label{U:preliminary-bounds-and-a-contradiction-set}

Strong convexity and convexity give \[
 d_k-(F_k-F_{k+1})\ge\frac\mu2\|s_k\|^2,
 \qquad t_k\ge\mu\|s_k\|^2.\tag{U.10}
\] Thus \(\sum\|s_k\|^2<\infty\) and \(s_k\to0\). Let \(\Gamma\) be the
compact cluster set. Every point of \(\Gamma\) has the common limit
value \(\ell\).

For Theorem U.2, assume \(\Gamma\) has no affine-critical point. Then
\(\ell>F_*\), so it contains no stationary point. For Theorem U.1, assume
directly \(\ell>F_*\).

Write \(r(x)=a^Tx\) in Theorem U.2, and \(r(x)=a^Tx-b\) in Theorem U.1. Let
\(E\) be the exceptional scalar set; in Theorem U.1 use \(E=\{0\}\).

The common tangential projection and its Lipschitz constant \(L_A\)
imply \[
 \|P_S y_k\|^2=\|A(x_{k+1})-A(x_k)\|^2
 \le L_A^2\|s_k\|^2\le\frac{L_A^2}{\mu}t_k.\tag{U.11}
\] Hence Theorem~\ref{M:residualtheorem}'s smooth-subspace estimate holds.

\subsection{5.2 Why tangential descent cannot disappear on the
exceptional
set}\label{U:why-tangential-descent-cannot-disappear-on-the-exceptional-set}

For Theorem U.2, at an exceptional point the nontrivial subgradient
segment has affine hull \(A(p)+\mathbb R a\). Thus \(A(p)=0\) would make
it affine critical, contrary to the assumption on \(\Gamma\).

For Theorem U.1, if \(p\in P\) and \(g\in\partial F(p)\), strong
monotonicity with \(0\in\partial F(x_*)\) gives \[
 A(p)^T(p-x_*)=g^T(p-x_*)\ge\mu\|p-x_*\|^2.
\] Thus \(A(p)=0\) on \(P\) is possible only at \(x_*\). Such a point
cannot belong to a positive limiting level.

In either case the compact set \[
 Z=\{p\in\Gamma:A(p)=0\}
\] has scalar image \(r(Z)\) disjoint from the closed set \(E\). If
\(Z\) is nonempty, these sets have positive separation
\(\delta=\operatorname{dist}(r(Z),E)>0\). If \(E\) is empty there is no
singular region and we may take the cutoff below identically one; if
\(Z\) is empty it may be identically zero.

Choose a scalar Lipschitz function \(\tau:\mathbb R\to[0,1]\) which is
zero where \(\operatorname{dist}(r,E)\le\delta/4\), and one where
\(\operatorname{dist}(r,E)\ge\delta/2\). Such a cutoff is obtained by
composing a piecewise-linear cutoff with the distance function. In
particular \(\tau=1\) on \(r(Z)\). The cutoff is an analytical object
chosen from \(F\) and the hypothetical cluster set, not an algorithmic
input.

\subsection{5.3 A frame on the entire cluster
set}\label{U:a-frame-on-the-entire-cluster-set}

Let the columns of \(U\) be an orthonormal basis of \(S\), and set \[
 \boxed{V(x)=[\,U,\ \tau(r(x))a\,].}\tag{U.12}
\] Its tangential columns are always present. Its normal column is
switched off near singular interfaces and restored away from them.

At differentiable points, \[
 \|V(x)^T\nabla F(x)\|^2
 =\|A(x)\|^2+\tau(r(x))^2|a^T\nabla F(x)|^2.
\] On \(\Gamma\setminus Z\) the first term is positive. On \(Z\), the
cutoff equals one and the full gradient is nonzero. The displayed
expression extends continuously across singular interfaces because the
cutoff vanishes on a whole strip around them; gradients are bounded and
\(A\) is continuous. Compactness therefore provides \(m>0\) and a
buffered neighborhood of all of \(\Gamma\) in which visibility holds.

The entire late trajectory lies in this neighborhood because its
distance to its compact cluster set tends to zero. There is no selection
of one cluster point followed by an assumption that the sequence
converges to it.

\subsection{5.4 The new secant term is
bounded}\label{U:the-new-secant-term-is-bounded}

The tangential part follows from (11). If \(\tau(r(z))=0\), there is no
normal contribution.

If \(\tau(r(z))>0\), the endpoint is a fixed positive distance from the
exceptional scalar set. For sufficiently close endpoints, the entire
segment lies in a slightly smaller smooth strip. On a compact buffered
subset of that strip, a finite convex-chart cover gives a common local
gradient-Lipschitz constant. Convex cocoercivity on each sufficiently
small convex chart yields \[
 \|\nabla F(z)-\nabla F(x)\|^2\le L_\delta\,y^Ts.
\] Therefore \[
 \|V(z)^Ty\|^2\le C_0t\tag{U.13}
\] for every sufficiently late accepted pair. This uses local smoothness
away from the exceptional set; \(L_\delta\) need not stay bounded as
\(\delta\downarrow0\).

\subsection{5.5 The only transport error is tangential and
inexpensive}\label{U:the-only-transport-error-is-tangential-and-inexpensive}

For an endpoint pair put \[
 \beta=\begin{cases}
 [\tau(r(z))-\tau(r(x))]/(a^Ts),&a^Ts\ne0,\\
 0,&a^Ts=0.
 \end{cases}
\] Then \(|\beta|\le\operatorname{Lip}(\tau)\). With \(e_n\) the last
frame-coordinate vector, \[
 \boxed{
 V(z)=V(x)+s\,\beta e_n^T-\beta P_Ss\,e_n^T.
 }\tag{U.14}
\] Thus (6) holds with \(T=I\) and \[
 R=-\beta P_Ss\,e_n^T\in S,
 \qquad
 \|R\|_F^2\le\operatorname{Lip}(\tau)^2\|P_Ss\|^2
 \le\frac{2\operatorname{Lip}(\tau)^2}{\mu}d.\tag{U.15}
\] All conditions of Theorem~\ref{M:residualtheorem} hold. The contradiction proves Theorem U.2,
and proves \(\ell=F_*\) for Theorem U.1. Compactness and uniqueness give
whole-sequence convergence. \(\square\)


\setcounter{section}{5}
\setcounter{theorem}{0}
\setcounter{lemma}{0}
\setcounter{corollary}{0}
\setcounter{proposition}{0}
\setcounter{remark}{0}
\setcounter{equation}{0}
\section{Only tangential strong convexity is needed for a composite
subclass}\label{U:only-tangential-strong-convexity-is-needed-for-a-composite-subclass}

Consider \(F(u,v)=\psi(u,v)+h(v)\) under the weaker hypotheses in
Section U.1.1. Set \(S=\mathbb R^{n-1}\times\{0\}\). Since
\(\psi-\mu\|u\|^2/2\) and \(h\) are convex, \[
 d_k-(F_k-F_{k+1})\ge\frac\mu2\|P_Ss_k\|^2.\tag{U.16}
\] This is the bound needed for the transport residual, even without a
positive normal curvature bound.

At differentiable accepted points, \[
 P_Sy=P_S[\nabla\psi(z)-\nabla\psi(x)].
\] Convex \(L\)-smoothness of \(\psi\) on the compact neighborhood, and
monotonicity of the scalar subgradient of \(h\), give \[
 \|P_Sy\|^2
 \le L[\nabla\psi(z)-\nabla\psi(x)]^Ts
 \le Lt.\tag{U.17}
\] This replaces (11). On the exceptional plane \(v=0\), the restriction
of \(F\) is a differentiable strongly convex function of \(u\), uniquely
minimized at \(u_*\). Thus the needed nonvanishing of \(A=\nabla_u\psi\)
on a nonoptimal cluster set still holds.

Finally, strict convexity of the positive sublevels supplies \(s_k\to0\)
in place of full strong convexity. The proof is dimension independent:
if a subsequence had \(x_k\to x\), \(x_{k+1}\to z\ne x\), \(g_k\to g\),
and \(F_k\to\ell>F_*\), then \(d_k\to0\) gives \(g^T(z-x)=0\). Convexity
and the supporting inequality make \(F\) constant at \(\ell\) on
\([x,z]\), contrary to the strictly convex boundary of \(\{F\le\ell\}\).

Equations (16)--(17) and the same cutoff frame complete the proof. A
quadratic-plus-absolute-value objective is included, but those familiar
cases were already proved in earlier notes. The extension below
illustrates the additional normal irregularity allowed here.


\setcounter{section}{6}
\setcounter{theorem}{0}
\setcounter{lemma}{0}
\setcounter{corollary}{0}
\setcounter{proposition}{0}
\setcounter{remark}{0}
\setcounter{equation}{0}
\section{An explicit all-dimensional example outside
RS}\label{U:an-explicit-all-dimensional-example-outside-rs}

Let \(a_j=2^{-j}\) for \(j\ge2\), and let \[
 \eta(t)=\begin{cases}(1-t^2)^2,&|t|\le1,\\0,&|t|>1.\end{cases}
\] This nonnegative bump is \(C^1\), has \(\eta(0)=1\), and integral
\(16/15\). Its supports at the following scales are pairwise disjoint.
Define \[
 k(v)=\sum_{j=2}^{\infty}a_j^{-3}
 \eta\!\left(\frac{v-a_j}{a_j^5}\right),\qquad v>0,
\] \[
 r(v)=\begin{cases}
 \displaystyle\int_0^v(v-t)k(t)\,dt,&v>0,\\
 0,&v\le0.
 \end{cases}\tag{U.18}
\] Each bump has mass \((16/15)a_j^2\). The total mass is finite. The
sum is locally finite away from zero, so \(r\) is convex, \(C^1\)
globally, and \(C^3\) on \((0,\infty)\) and \((-\infty,0)\). In
particular it is locally \(C^{1,1}\) away from zero. Direct bounds give
\[
 0\le r'(v)\le6v^2,\qquad 0\le r(v)\le6v^3
\] near zero. However, \[
 r''(a_j)=a_j^{-3}\to\infty.\tag{U.19}
\] (The second derivative at zero itself exists and equals zero, but is
not continuous there.)

For any \(n\ge2\), consider \[
 \boxed{
 F(u,v)=\frac12\|u\|^2+\frac12v^2+|v|+r(v),
 \qquad u\in\mathbb R^{n-1}.
 }\tag{U.20}
\] It is 1-strongly convex and coercive, with unique minimizer zero. Its
only nondifferentiable hyperplane is \(v=0\); its tangential gradient is
\(u\), which is globally Lipschitz. Theorem U.1 therefore proves global
convergence of every admissible full-BFGS run, for arbitrary coupled SPD
initialization.

Yet \textbf{RS fails in every neighborhood of zero}. Fix any
sufficiently small nonzero \(u_0\), and take \(x_j=(u_0,a_j)\). Their
objective values approach the positive level \(\|u_0\|^2/2\). At these
smooth points, \[
 x_j^T\nabla^2F(x_j)x_j
 =\|u_0\|^2+a_j^2+a_j^{-1}\longrightarrow\infty.\tag{U.21}
\] A uniform inward Taylor estimate on that fixed positive band would,
by taking its second-order limit at each \(x_j\), bound the left side of
(21). This is impossible. The recorded equivalence between inward radial
smoothness and RS therefore rules out RS on that band. Choosing \(u_0\)
arbitrarily small gives failure in every neighborhood of the optimizer.

This is a \textbf{function-class} gain beyond RS, but not a claim that
the function lacks every prior transported certificate. It also expands
the allowed normal curvature rather than settling the original arbitrary
curved finite-\(C^2\)-max conjecture. Where each piece has a \(C^{1,1}\)
extension through the flat interface, the older
finite-polyhedral-partition theorem already gave local convergence.

\subsection{7.1 The infinite series is evaluated exactly in the
verifier}\label{U:the-infinite-series-is-evaluated-exactly-in-the-verifier}

There is no need to truncate (18). All bumps below a query can be summed
using geometric series, with at most one partially included bump. Define
\[
 P(t)=t-\frac23t^3+\frac15t^5,\qquad
 Q(t)=\frac12t^2-\frac12t^4+\frac16t^6.
\] For a query inside the bump centered at \(a=a_j\), with
\(t=(v-a)/a^5\), the partial mass is \(a^2[P(t)+8/15]\), and its first
moment is \[
 a^3[P(t)+8/15]+a^7[Q(t)-1/6].
\] All smaller bumps have total mass \((16/45)a^2\) and first moment
\((16/105)a^3\). This gives exact rational function and derivative
values for rational queries. At a peak, \[
 r(a)=\frac{64}{315}a^3+\frac16a^7,\qquad
 r'(a)=\frac89a^2,\qquad r''(a)=a^{-3}.
\] The numerical full-BFGS tests use this complete infinite-series
oracle and a dense positive-definite quadratic in place of the diagonal
smooth term in (20). They do not assume block-diagonal matrices or
perform a reduced two-dimensional algorithm.

\subsection{7.2 The objective need not split as a smooth function plus
one scalar
penalty}\label{U:the-objective-need-not-split-as-a-smooth-function-plus-one-scalar-penalty}

For a fixed vector \(c\in\mathbb R^{n-1}\) with \(\|c\|<1\), another
example is \[
 F(u,v)=\frac12(\|u\|^2+v^2)+|(1+c^Tu)v|+r(v).
\] The two quadratic functions before adding \(r\) have Hessians \[
 \begin{pmatrix}I&c\\c^T&1\end{pmatrix},\qquad
 \begin{pmatrix}I&-c\\-c^T&1\end{pmatrix},
\] both bounded below by \((1-\|c\|)I\). Their maximum plus convex
\(r(v)\) is strongly convex. Since the absolute-value term and \(r\) are
nonnegative, zero is the unique minimizer.

On a neighborhood with \(|c^Tu|<1/2\), the only exceptional hyperplane
is \(v=0\), and the tangential gradient is \[
 A(u,v)=u+c|v|,
\] which is Lipschitz across that hyperplane. Theorem U.1 gives local
full-BFGS convergence. The jump magnitude changes with \(u\), so the
theorem is not relying on an additive smooth-plus-univariate
decomposition. The same spike calculation makes RS fail locally for this
example as well.


\endgroup

% END original chapters/U.tex

\StartModule{D}{Exact deterministic local failure}{Canonical source [SD]. Algebraic root, full matrix closure, all accepted and rejected trials, and a globally $C^{2,1}$ convex-branch embedding. No positive lower curvature is asserted.}

% BEGIN original chapters/D.tex
\begingroup
\providecommand{\R}{}
\renewcommand{\R}{\mathbb R}
\providecommand{\B}{}
\renewcommand{\B}{\mathcal B}
\providecommand{\I}{}
\renewcommand{\I}{I}
\setcounter{section}{0}
\setcounter{theorem}{0}
\setcounter{lemma}{0}
\setcounter{corollary}{0}
\setcounter{proposition}{0}
\setcounter{remark}{0}
\setcounter{equation}{0}
\section{Main result and exact implementation}
For a differentiable current point $x$, put $g=\nabla F(x)$ and $p=-Hg$, with $H\succ0$. We use the ordinary full inverse-BFGS update
\begin{equation}\label{D:bfgs}
 s=\alpha p,\quad y=\nabla F(x+s)-g,\quad t=y^Ts,\qquad
 H_+=\left(I-\frac{sy^T}{t}\right)H\left(I-\frac{ys^T}{t}\right)+\frac{ss^T}{t}.
\end{equation}
There is no resetting, damping, rescaling of the stored matrix, memory truncation, or skipped update.

The line search is Algorithm 4.6 of Lewis--Overton~\cite{LO}: initialize a lower bracket $\ell=0$, an upper bracket $h=+\infty$, and $\alpha=1$. At each trial:
\begin{enumerate}
\item If Armijo fails, set $h=\alpha$.
\item Otherwise, if weak Wolfe fails, set $\ell=\alpha$.
\item Otherwise, accept $\alpha$ and stop.
\end{enumerate}
After a rejection, try $(\ell+h)/2$ if $h<+\infty$, and $2\ell$ otherwise. The tests are
\[
 F(x+\alpha p)<F(x)+c_1\alpha g^Tp,
 \qquad \nabla F(x+\alpha p)^Tp>c_2g^Tp.
\]
All decisions below are strict; using non-strict acceptance inequalities produces the identical orbit. Every trial point, not just every accepted point, is differentiable.

\begin{theorem}\label{D:main}
There exist four convex functions $f_i\in C^2(\R^2)$, with
\[
 0\preceq\nabla^2 f_i(x)\preceq64I \quad(i=0,1,2,3;\ x\in\R^2),
\]
such that $F=\max_i f_i$ has compact sublevel sets and the unique global minimizer $x_*=0$, where $F$ is nonsmooth. For the deterministic line search specified above, with $c_1=10^{-4}$ and $c_2=1/2$, there are starting points $x_0^{(n)}\to0$, all initialized with $H_0=I$, whose full-BFGS iterates satisfy
\[
 x_k^{(n)}\longrightarrow p_n\ne0,\qquad
 F(x_k^{(n)})\downarrow F(p_n)>0,\qquad 0\notin\partial F(p_n).
\]
Every line search terminates after at most five trials. In fact, the same function, initial points, matrices, and trial decisions work for every fixed
\[
 0<c_1\le\tfrac13,\qquad c_1<c_2<1.
\]
The entire accepted and trial trajectories can be contained in an arbitrarily small neighborhood of $0$ by choosing $n$ sufficiently large.
\end{theorem}

The accepted lengths are not assigned externally. They are generated by the deterministic search:
\[
 \boxed{\alpha_{3m}=1,\qquad\alpha_{3m+1}=\tfrac14,\qquad\alpha_{3m+2}=\tfrac1{16}.}
\]
Thus the construction resolves the implementation-selection issue left open by the previous admissible-length counterexample. It does not address strong-Wolfe interpolation searches or a positive-measure basin of failing initial points.


\setcounter{section}{1}
\setcounter{theorem}{1}
\setcounter{lemma}{0}
\setcounter{corollary}{0}
\setcounter{proposition}{0}
\setcounter{remark}{0}
\setcounter{equation}{1}
\section{An exactly specified algebraic seed}
Let $(c,d)$ be the real pair specified below. Define
\begin{align*}
 U&=127c+16d,& V&=16c-127d-16,& R&=c-8d,\\
 K&=8c^2+8d^2+d-8,& T&=c^2-c+d^2,& Y&=(c-1)^2+d^2.
\end{align*}
The defining equations are
\begin{equation}\label{D:polys}
 \boxed{264192RK-U V(U+129)=0,\qquad16VTK-RUY=0.}
\end{equation}
Choose the unique root in the closed $\ell_\infty$ box of radius $10^{-30}$ centered at the rational pair
\begin{align*}
 c_0&=0.31848142762711863064645131250290703717912205023605,\\
 d_0&=-1.02523211967574872548135087001300414284511909318444.
\end{align*}
The decimals here are exact terminating rationals defining the box, not replacements for the algebraic root. Set
\begin{equation}\label{D:rho}
 \rho=\frac{129R}{V(U+129)}
       \simeq0.060199206107543764815373666287.
\end{equation}

\begin{lemma}\label{D:root}
The designated box contains exactly one solution of~\eqref{D:polys}. All denominators in this note are nonzero there, and $0.060<\rho<0.061$.
\end{lemma}
\begin{proof}
Let $P=(P_1,P_2)^T$ denote the left sides of~\eqref{D:polys}, let $\bar r=(c_0,d_0)$, and let $M=DP(\bar r)^{-1}$, an exactly specified rational matrix. On the box $X$, exact rational interval arithmetic gives
\[
 \sup_{z\in X}\|I-MDP(z)\|_\infty<10^{-26},
 \qquad\|MP(\bar r)\|_\infty<10^{-50}.
\]
Hence $z\mapsto z-MP(z)$ is a strict contraction mapping $X$ into itself, because
\[
 10^{-50}+10^{-26}\,10^{-30}<10^{-30}.
\]
The contraction theorem gives the unique root. The interval calculation also bounds the denominators away from zero and certifies~\eqref{D:rho}. The standard-library program \texttt{verify\_certificate.py} constructs the two polynomials, differentiates them, constructs $M$, and checks these rational inequalities without floating-point arithmetic. This is an existence certificate, not a residual-tolerance assertion.
\end{proof}

Set
\[
 g_0=(1,0)^T,\qquad g_1=(-1,-1/8)^T,\qquad g_2=(c,d)^T,
 \qquad G(z)=\max_{0\le i\le2}g_i^Tz.
\]
Start with $M_0=I$, and define
\[
 M_1=\begin{pmatrix}129/256&-1/16\\-1/16&1\end{pmatrix},
 \qquad M_2=\B(M_1,g_1,g_2,1/4),
\]
where $\B$ denotes~\eqref{D:bfgs} with the stated current and next gradients. Then
\begin{equation}\label{D:closure}
 \B(I,g_0,g_1,1)=M_1,\qquad
 \boxed{\B(M_2,g_2,g_0,1/16)=\rho I.}
\end{equation}
The three displacements are
\begin{equation}\label{D:steps}
 s_0=(-1,0)^T,\quad s_1=(127/1024,1/64)^T,\quad
 s_2=\rho(1-c,-d)^T.
\end{equation}
Define the starting point by
\begin{equation}\label{D:x0}
 z_0=\frac1{1-\rho}
 \begin{pmatrix}897/1024-\rho(1-c)\\ \rho d-1/64\end{pmatrix},
 \qquad z_{j+1}=z_j+s_j.
\end{equation}
Then
\[
 \boxed{z_3=\rho z_0.}
\]
For clarity, approximate points are
\[
\begin{array}{c|rr}
 j& (z_j)_1&(z_j)_2\\\hline
0& 0.888432624149&-0.082297397686\\
1&-0.111567375851&-0.082297397686\\
2& 0.012456061649&-0.066672397686\\
3& 0.053482938654&-0.004954238006
\end{array}
\]
The uniquely active gradient at $z_j$ is $g_j$ for $j=0,1,2$, and $g_0$ at $z_3$.

\subsection{Derivation and independent checks of the matrix closure}
The identity is not inferred from near-closure. For arbitrary symbolic $(c,d)$ with the denominators nonzero, put
\[
 \beta=\frac{66048RK}{UV(U+129)},\qquad
 \gamma=\frac{VTK}{RUY},\qquad
 \widetilde\rho=\frac{2064TK}{U(U+129)Y}.
\]
Direct rational identities give
\[
 \B\bigl(\B(M_1,g_1,g_2,\beta),g_2,g_0,\gamma\bigr)
 =\widetilde\rho I.
\]
One way to derive the formulas is to write the middle update as $C+\beta D$. Requiring $(C+\beta D)g_2$ to be parallel to $g_0-g_2$ is linear in $\beta$. Its solution is the formula above. If $w\perp(g_0-g_2)$, set
\[
 \widetilde\rho=\frac{w^T(C+\beta D)w}{w^Tw}
\]
and choose $\gamma$ so that the last displacement is $\widetilde\rho(g_0-g_2)$. The last update then gives $\widetilde\rho I$: its two projection factors are the orthogonal projector onto $(g_0-g_2)^\perp$.

Equations~\eqref{D:polys} are exactly $\beta=1/4$ and $\gamma=1/16$. Also
\[
 \widetilde\rho-\rho
 =\frac{2064R}{V(U+129)}\left(\gamma-\frac1{16}\right).
\]
This proves~\eqref{D:closure} and~\eqref{D:steps}. The separate SymPy program \texttt{verify\_symbolic\_closure.py} checks these identities exactly, along with the secant equations and an independent Hessian-form update calculation. Interval Sylvester tests certify positive definiteness at the designated root.


\setcounter{section}{2}
\setcounter{theorem}{1}
\setcounter{lemma}{1}
\setcounter{corollary}{0}
\setcounter{proposition}{0}
\setcounter{remark}{0}
\setcounter{equation}{6}
\section{Every deterministic oracle decision is certified}
Use $c_1=10^{-4}$ and $c_2=1/2$ for the displayed table. The line searches below use $p_j=-M_jg_j$. At each rejected trial, the objective is strictly greater than the current value, not merely greater than the Armijo threshold.
\begin{center}\small
\begin{tabular}{cccrl}\toprule
Phase & Trial $\alpha$ & Active branch & $G(z_j+\alpha p_j)$ & Decision\\\midrule
0&1&1&0.121854550561&Accept\\\midrule
1&1&0&0.384526374149&Armijo rejection\\
1&1/2&0&0.136479499149&Armijo rejection\\
1&1/4&2&0.072321707900&Accept\\\midrule
2&1&0&0.668886093720&Armijo rejection\\
2&1/2&0&0.340671077685&Armijo rejection\\
2&1/4&0&0.176563569667&Armijo rejection\\
2&1/8&0&0.094509815658&Armijo rejection\\
2&1/16&0&0.053482938654&Accept\\\bottomrule
\end{tabular}
\end{center}
The starting values in the three phases are approximately
\[
 (0.888432624149,\;0.121854550561,\;0.072321707900).
\]
The ratios at the accepted steps are
\[
\begin{array}{c|rr}
 j & \dfrac{G(z_j)-G(z_{j+1})}{-g_j^Ts_j}
 &\dfrac{g_{j+1}^Ts_j}{g_j^Ts_j}\\\hline
0&0.766578073588&-1\\
1&0.393190937094&-0.186383158092\\
2&0.375205963572&-0.817119670486
\end{array}
\]
The displayed decimals are explanatory. The verifier encloses every value and strict comparison using exact rational intervals for the algebraic root. In particular, each accepted Armijo ratio exceeds $1/3$, and every accepted Wolfe ratio is negative. The rejected trial values exceed their respective current values. Therefore the identical trial tree works for \emph{all} $0<c_1\le1/3$ and $c_1<c_2<1$.

All nine trial points have a unique active branch. They satisfy the additional uniform bounds
\begin{equation}\label{D:trialbounds}
 |g_i^Tz_{\rm trial}|<3\quad(i=0,1,2),
 \qquad (z_{\rm trial})_2<2,
 \qquad G(z_{\rm trial})>0.
\end{equation}
The accepted points satisfy these bounds as well.

Positive homogeneity and~\eqref{D:closure}--\eqref{D:x0} prove by induction that the deterministic full-BFGS run on $G$ is
\[
 x_{3m+j}=\rho^m z_j,\qquad H_{3m+j}=\rho^m M_j.
\]
Every trial at the corresponding phase is exactly $\rho^m$ times its listed base trial. All comparison quantities scale by $\rho^m$, preserving every decision, not only acceptance at the selected length. Thus $x_k\to0$ and $G(x_k)\downarrow0$.

This limit is nonstationary: $d_* =(-1,16)^T$ satisfies
\[
 g_0^Td_*=-1,\quad g_1^Td_*=-1,\quad g_2^Td_*=-c+16d<-1.
\]
In particular $G(td_*)=-t$ for $t\ge0$ and $0\notin\operatorname{conv}\{g_0,g_1,g_2\}$.

\subsection{A simple coercive global failure}
The same trajectory is a deterministic counterexample for the finite polyhedral function
\[
 \overline G(u,v)=\max\{u,-u-v/8,cu+dv,v-17\}.
\]
All trial values of $G$ are positive while $v-17<0$, by~\eqref{D:trialbounds} and scaling. The fourth piece never affects the line search. The convex combination of branches $0,1,3$ with weights $(8/17,8/17,1/17)$ is identically $-1$. Equality fixes the unique minimizer $(-1,16)$, and branch $2$ is strictly below $-1$ there. Sublevels are bounded, because $u\le t$, $v\le t+17$, and $-u-v/8\le t$ give lower bounds on both coordinates. This supplies a coercive deterministic failure, but does not alone refute local convergence at its true minimizer. The next construction does.


\setcounter{section}{3}
\setcounter{theorem}{1}
\setcounter{lemma}{1}
\setcounter{corollary}{0}
\setcounter{proposition}{0}
\setcounter{remark}{0}
\setcounter{equation}{7}
\section{Two smooth convex scalar functions}
Set
\[
 \vartheta=\frac1{16},\qquad a_n=\vartheta^n\ (n\ge0),\qquad
 \epsilon_0=\frac1{128},\qquad b_* =2c-32d\in(33,34).
\]
We define scalar functions only for $b\in\{2,b_*\}$. Let
\[
 c_b=\frac{1-(1+\vartheta+\vartheta^2)/b}{1-\vartheta^2}.
\]
Define a monotone $C^1$ ramp $P_b$ on $[0,1]$ by
\[
P_b(\xi)=
\begin{cases}
0,&\xi\le c_b-\epsilon_0,\\
3t^2-2t^3,&c_b-\epsilon_0<\xi<c_b+\epsilon_0,
\quad t=(\xi-c_b+\epsilon_0)/(2\epsilon_0),\\
1,&\xi\ge c_b+\epsilon_0.
\end{cases}
\]
Then $0\le P_b'\le96$ and $\int_0^1P_b(\xi)\,d\xi=1-c_b$. Its primitive $J_b(\xi)=\int_0^\xi P_b(t)\,dt$ is
\[
J_b(\xi)=\begin{cases}
0,&\xi\le c_b-\epsilon_0,\\
2\epsilon_0(t^3-t^4/2),&c_b-\epsilon_0<\xi<c_b+\epsilon_0,\\
\xi-c_b,&\xi\ge c_b+\epsilon_0.
\end{cases}
\]
For $x\in[b\vartheta a_n,ba_n]$, put
\[
 \xi=\frac{x/(ba_n)-\vartheta}{1-\vartheta},
\]
and define
\begin{equation}\label{D:scalar}
 \varphi_b(x)=(\vartheta a_n)^3+
 ba_n^3(1-\vartheta)
 \left[\vartheta^2\xi+(1-\vartheta^2)J_b(\xi)\right].
\end{equation}
Complete its definition by
\[
\varphi_b(x)=\begin{cases}
1-3(x+1)+3(x+1)^2,&x<-1,\\
-x^3,&-1\le x\le0,\\
1+x-b,&x\ge b.
\end{cases}
\]
These intervals cover the real line, with agreeing formulas at shared endpoints.

\begin{lemma}\label{D:scalarlemma}
Each $\varphi_b$ is globally convex, $C^2$, nonnegative, coercive, and uniquely minimized at $0$. Moreover,
\begin{gather}
 \varphi_b(ba_n)=a_n^3,\quad
 \varphi_b'(ba_n)=a_n^2,\quad\varphi_b''(ba_n)=0,\label{D:jets}\\
 0\le\varphi_b''(x)\le\max\{6,102/b\}\quad(x\in\R),\label{D:curvature}\\
 \boxed{\varphi_b(ba_n+h)=a_n^3+a_n^2h\quad\text{for }|h|\le a_n/2.}\label{D:affineplateau}
\end{gather}
\end{lemma}
\begin{proof}
Differentiate~\eqref{D:scalar} on its positive intervals:
\[
 \varphi_b'(x)=a_n^2[\vartheta^2+(1-\vartheta^2)P_b(\xi)],
 \qquad
 \varphi_b''(x)=\frac{a_n(1+\vartheta)}bP_b'(\xi).
\]
The latter is nonnegative and at most $102a_n/b$. The identity
\[
 \vartheta^2+(1-\vartheta^2)(1-c_b)
 =\frac{1+\vartheta+\vartheta^2}{b}
\]
ensures the required integral and value matching at every node. The derivative is constant near each interval endpoint; thus the first two derivatives match adjoining intervals.

For $b\vartheta a_n\le x\le ba_n$, $a_n\le x/(b\vartheta)$. Consequently as $x\downarrow0$,
\[
 \varphi_b(x)=O(x^3),\qquad\varphi_b'(x)=O(x^2),\qquad
 \varphi_b''(x)=O(x).
\]
They match the negative formula at zero, including two derivatives. The joins at $-1$ and $b$ are also $C^2$. The derivative is negative on $x<0$ and positive on $x>0$; both outer extensions diverge to infinity. This proves convexity, the minimum assertions, coercivity, and~\eqref{D:curvature}.

The constant-derivative lengths below and above $ba_n$, divided by $a_n$, are respectively
\[
 \ell_b=b(1-\vartheta)(1-c_b-\epsilon_0),\qquad
 u_b=\frac{b(1-\vartheta)(c_b-\epsilon_0)}{\vartheta}.
\]
At $n=0$ the affine outer extension supplies the upper length. Exact rational interval bounds give
\[
\begin{array}{c|ccc}
 b&c_b&\ell_b&u_b\\\hline
2&0.468627450981&0.981675091912&13.824448529412\\
b_*&0.971910573260&0.635765118558&483.655089309488
\end{array}
\]
Both transition endpoints are strictly inside $(0,1)$ and both lengths exceed $1/2$. This proves~\eqref{D:affineplateau}.
\end{proof}

\paragraph{Additional regularity.}
The construction even gives globally Lipschitz second derivatives. Within a ramp,
\[
 |\varphi_b'''(x)|
 \le\frac{1+\vartheta}{b^2(1-\vartheta)}
       \frac{3}{2\epsilon_0^2}
 =\frac{139264}{5b^2}.
\]
The second derivative is continuous across every join and at the accumulation point $0$. The uniform bound above, together with the bound $6$ on the negative cubic part, proves global Lipschitz continuity of $\varphi_b''$ by summing over interval lengths. Thus the branches below are in $C^{2,1}$ as well. This does not make them analytic or uniformly strongly convex.


\setcounter{section}{4}
\setcounter{theorem}{1}
\setcounter{lemma}{2}
\setcounter{corollary}{0}
\setcounter{proposition}{0}
\setcounter{remark}{0}
\setcounter{equation}{8}
\section{The single objective with bad starts arbitrarily near its minimizer}
Define
\begin{equation}\label{D:fullobjective}
\begin{split}
 f_0(u,v)&=\varphi_2(u),\\
 f_1(u,v)&=\varphi_2(-u-v/8),\\
 f_2(u,v)&=\varphi_{b_*}(cu+dv),\\
 f_3(u,v)&=v,\qquad F(u,v)=\max_{0\le i\le3}f_i(u,v).
\end{split}
\end{equation}
Each branch is convex and $C^2$. For $i\le2$, its Hessian is the scalar second derivative times $g_i g_i^T$. Thus
\[
 \|\nabla^2f_0\|\le51,\qquad
 \|\nabla^2f_1\|\le51\frac{65}{64}<52,
\]
while $b_*>33$ gives $\varphi_{b_*}''\le6$ and $c^2+d^2<2$, so
\[
 \|\nabla^2f_2\|<12,\qquad\nabla^2f_3=0.
\]
A common global upper bound $64I$ therefore holds. The preceding derivative estimate and $\|g_i\|<2$ also give a common finite Lipschitz bound for the branch Hessians (for example $60000$ suffices).

The first three branches are nonnegative and vanish at zero. If $F(x)=0$, its first two scalar arguments must vanish; $g_0$ and $g_1$ are independent, so $x=0$. The same two scalar arguments lie in bounded scalar sublevels whenever $F(x)$ is bounded above, proving compactness of every sublevel of $F$.

At zero, the first three branches have zero gradient. Therefore
\[
 F'(0;h)=\max\{0,h_2\},
\]
which proves that the unique minimizer is nonsmooth.

Put
\[
 p_n=a_n(2,-32)^T.
\]
Then $g_0^Tp_n=g_1^Tp_n=2a_n$ and $g_2^Tp_n=b_*a_n$. Lemma~\ref{D:scalarlemma} yields the exact local representation
\begin{equation}\label{D:localcone}
 \boxed{F(p_n+h)=a_n^3+a_n^2G(h)}
\end{equation}
whenever $|g_i^Th|\le a_n/2$ for $i=0,1,2$. Indeed, all three transformed branches are in their affine intervals and have value at least $a_n^3/2$. The first two slab inequalities imply $|h_1|\le a_n/2$ and $|h_2|\le8a_n$, so $f_3(p_n+h)<0$ and is inactive.

At $p_n$, the direction $d_*=(-1,16)^T$ strictly decreases all three active branches. Hence
\[
 F(p_n)=a_n^3>0,\qquad
 F'(p_n;d_*)=-a_n^2<0,
 \qquad0\notin\partial F(p_n).
\]


\setcounter{section}{5}
\setcounter{theorem}{1}
\setcounter{lemma}{2}
\setcounter{corollary}{0}
\setcounter{proposition}{0}
\setcounter{remark}{0}
\setcounter{equation}{10}
\section{All rejected trials also stay in the local model}
For each $n\ge1$, initialize
\[
 \boxed{x_0^{(n)}=p_n+a_n^2z_0,\qquad H_0=I.}
\]
The claimed deterministic trajectory is
\begin{equation}\label{D:finalorbit}
 x_{3m+j}^{(n)}=p_n+a_n^2\rho^mz_j,\qquad
 H_{3m+j}^{(n)}=\rho^mM_j,\qquad j=0,1,2.
\end{equation}
At the corresponding phase, \emph{every} trial generated by the line search has the form
\[
 p_n+a_n^2\rho^m z_{\rm trial},
\]
where $z_{\rm trial}$ is one of the nine base trials already certified. By~\eqref{D:trialbounds},
\[
 |g_i^T(a_n^2\rho^mz_{\rm trial})|
 <3a_n^2\le\frac{3a_n}{16}<\frac{a_n}{2}.
\]
Thus~\eqref{D:localcone} holds at every accepted and rejected trial. Their values and gradients are
\[
 F(p_n+a_n^2\rho^mz_{\rm trial})
 =a_n^3+a_n^4\rho^mG(z_{\rm trial}),
 \qquad \nabla F=a_n^2g_{\rm active}.
\]
Every function difference and every predicted directional decrease is the corresponding seed quantity multiplied by $a_n^4\rho^m$. The additive $a_n^3$ cancels from Armijo. Therefore all rejection and acceptance decisions are the same as in the seed table. In particular the accepted lengths are actual outputs of the deterministic search.

Substituting gradients scaled by $a_n^2$ and steps scaled by $a_n^2\rho^m$ into~\eqref{D:bfgs} leaves the projection factors unchanged and scales the rank-one term by $\rho^m$. This proves the matrix part of~\eqref{D:finalorbit}. The point and matrix cycle closures complete the induction. All finite matrices are SPD, and all oracle evaluations are at differentiable points.

Consequently,
\[
 x_k^{(n)}\to p_n\ne0,\qquad
 F(x_k^{(n)})\downarrow a_n^3>0.
\]
Yet $x_0^{(n)}\to0$ and $F(x_0^{(n)})\to0$. Also the supremum of the norms of \emph{all accepted and trial points} in the $n$th trajectory is bounded by
\[
 a_n\|(2,-32)\|+a_n^2\max_{\rm base\ trials}\|z_{\rm trial}\|\to0.
\]
This proves Theorem~\ref{D:main}. The failure is not caused by leaving a local neighborhood.


\endgroup

% END original chapters/D.tex

\StartModule{L}{Failure with positive lower curvature}{Complete construction from [SLC]. Stable-manifold persistence and local lower-curvature embeddings. The local smooth branches are $C^{1,1}$, not $C^2$ at the accumulation point.}

% BEGIN original chapters/L.tex
\begingroup
\providecommand{\R}{}
\renewcommand{\R}{\mathbb R}
\providecommand{\T}{}
\renewcommand{\T}{^{\mathsf T}}
\providecommand{\norm}{}
\renewcommand{\norm}[1]{\left\lVert#1\right\rVert}
\providecommand{\ip}{}
\renewcommand{\ip}[2]{\langle #1,#2\rangle}
\setcounter{section}{0}
\setcounter{theorem}{0}
\setcounter{lemma}{0}
\setcounter{corollary}{0}
\setcounter{proposition}{0}
\setcounter{remark}{0}
\setcounter{equation}{0}
\section{Statements and retained implementation}
All runs use ordinary full BFGS, $H_0\succ0$, differentiable oracle queries, and a fixed unit-first Armijo--weak-Wolfe search. The search starts with $\alpha=1$, lower bracket zero and upper bracket infinity; an Armijo failure updates the upper bracket, a Wolfe failure the lower bracket; use the midpoint after bracketing and doubling otherwise. Fix
\[
 c_1=10^{-4},\qquad c_2=\tfrac12.
\]
The update is unmodified:
\[
 s=-\alpha Hg,\qquad y=g^+-g,\qquad
 H^+=\left(I-\frac{sy\T}{y\T s}\right)H
 \left(I-\frac{ys\T}{y\T s}\right)+\frac{ss\T}{y\T s}.
\]
The seed inequalities have strict margins, so there are no equality-convention issues.

\begin{theorem}[Global failure with constant positive branch Hessians]\label{L:thm:global}
There are three fixed slopes $g_0,g_1,g_2\in\R^2$ such that, for every $\mu>0$,
\[
 Q_\mu(x)=\tfrac\mu2\norm{x}^2+\max_{i=0,1,2}g_i\T x
\]
is $\mu$-strongly convex, coercive, and has a unique nonsmooth minimizer $x_*\ne0$, yet admits initial states $(x_0,H_0)$ producing deterministic full-BFGS iterates converging to the nonstationary origin. Initial states and every future trial may be chosen arbitrarily close to the origin in the sense made precise in Section~\ref{L:sec:blowup}. Every branch has Hessian exactly $\mu I$ globally. The accepted lengths repeat $1,1/4,1/16$.
\end{theorem}
\noindent This is \emph{global} failure: approaching the wrong vertex is not the same as approaching the actual minimizer. It does not contradict our local PLQ theorem.

\begin{theorem}[Local failure under the pointwise Hessian bounds]\label{L:thm:local}
There is a finite strongly convex objective $F:\R^2\to\R$ with compact sublevels and a unique nonsmooth minimizer at zero, and constants $0<\underline\mu\le\overline L<\infty$, such that
\[
 \underline\mu I\preceq \nabla^2F(x)\preceq\overline L I
 \quad\text{at every twice-differentiable point},
\]
yet the fixed deterministic full-BFGS implementation with $H_0=I$ has bad initial points $x_0^{(n)}\to0$ converging to nonstationary points $p_n\to0$. All trial points of the $n$th run remain uniformly near zero as $n\to\infty$.

The function can be chosen as a maximum of four globally $C^{1,1}$, uniformly strongly convex, globally gradient-Lipschitz branches. Alternatively, it can be chosen so that its Hessian is one fixed positive-definite matrix wherever it is twice differentiable; the latter version has a countable polyhedral nonsmooth part, not a finite smooth-branch representation.
\end{theorem}
\noindent The local theorem satisfies the user's literal classical-Hessian assumption, and even actual strong convexity. It is \emph{not} a local counterexample with four globally $C^2$ branches having uniform lower curvature. That stronger combination remains unproved here.


\setcounter{section}{1}
\setcounter{theorem}{2}
\setcounter{lemma}{0}
\setcounter{corollary}{0}
\setcounter{proposition}{0}
\setcounter{remark}{0}
\setcounter{equation}{0}
\section{The inherited algebraic seed}
The seed is reproduced sufficiently to specify the new calculations. Its exact symbolic closure and nine strict query decisions were certified in the preceding deterministic-bracketing report~\cite{SD}.
\[
 g_0=(1,0)\T,\qquad g_1=(-1,-1/8)\T,\qquad g_2=(c,d)\T.
\]
Put
\begin{align*}
 U&=127c+16d,&V&=16c-127d-16,&R&=c-8d,\\
 K&=8c^2+8d^2+d-8,&T&=c^2-c+d^2,&Y&=(c-1)^2+d^2.
\end{align*}
The algebraic constants solve
\[
 264192RK-UV(U+129)=0,\qquad 16VTK-RUY=0
\]
in the rational $\ell_\infty$ box of radius $10^{-30}$ about
\begin{align*}
 c_0&=0.31848142762711863064645131250290703717912205023605,\\
 d_0&=-1.02523211967574872548135087001300414284511909318444.
\end{align*}
The seed verifier proves existence and uniqueness of this root by a rational-interval contracting self-map. Define
\[
 \rho=\frac{129R}{V(U+129)}\in(0.060,0.061),\qquad
 z_0=\frac1{1-\rho}
 \begin{pmatrix}897/1024-\rho(1-c)\\ \rho d-1/64\end{pmatrix}.
\]
Let $M_0=I$ and use consecutive gradients $g_0,g_1,g_2,g_0$ with lengths $1,1/4,1/16$. Exact full-BFGS algebra gives
\[
 M_1=\begin{pmatrix}129/256&-1/16\\-1/16&1\end{pmatrix},
 \qquad M_3=\rho I,\qquad z_3=\rho z_0.
\]
Here $z_{j+1}=z_j-\alpha_jM_jg_j$. Every displayed accepted gradient is uniquely active for $G(z)=\max_i g_i\T z$.

The finite line-search query lists are
\[
 [1],\qquad[1,1/2,1/4],\qquad[1,1/2,1/4,1/8,1/16].
\]
Every rejected trial has a value larger than the current value. The accepted decrease ratios exceed $1/3$, and the accepted new-to-old directional-derivative ratios are negative. All query points have a unique active branch. Thus the query decisions and branch itinerary persist in a sufficiently small open neighborhood of the normalized seed state. This strictness is essential below.

The vector $(-1,16)\T$ has negative inner product with every $g_i$, so $0\notin\operatorname{conv}\{g_i\}$. On the other hand,
\[
 \bar g=(1,-16)\T/257
 =\frac{129}{257}g_0+\frac{128}{257}g_1
\]
is the least-norm point of this convex hull. Indeed $g_0\T\bar g=g_1\T\bar g=1/257=\norm{\bar g}^2$ and $g_2\T\bar g>1/257$. Consequently
\begin{equation}\label{L:eq:minimizer}
 x_*=-\bar g/\mu,\qquad \min Q_\mu=-\frac1{514\mu}.
\end{equation}
At $x_*$ precisely the first two branches are active; their gradients differ, so the unique minimizer is nonsmooth. At the origin $Q_\mu(0)=0$ and the subdifferential is the same nonstationary triangle as for $G$.


\setcounter{section}{2}
\setcounter{theorem}{2}
\setcounter{lemma}{0}
\setcounter{corollary}{0}
\setcounter{proposition}{0}
\setcounter{remark}{0}
\setcounter{equation}{1}
\section{Blowing up the shrinking cycle}\label{L:sec:blowup}
Adding a quadratic changes the BFGS iterates. We do not claim the old exact orbit survives unchanged. Instead we prove existence of nearby perturbed orbits.

Use the five normalized state coordinates
\[
 x=\tau z,\qquad H=\tau M(b,e),\qquad
 M(b,e)=\begin{pmatrix}1&b\\b&e\end{pmatrix},\qquad \tau>0.
\]
The limiting seed state is
\[
 P_*=(z_0,0,1,0),
\]
where the last three entries denote $(b,e,\tau)$. To define a three-iteration return map near this boundary state, keep the input scale $\tau$ fixed during the three substeps. Start with $(Z,M)=(z,M(b,e))$ and, for $j=0,1,2$, set
\begin{align}
 \gamma_j&=g_j+\mu\tau Z,\qquad
 S=-\alpha_jM\gamma_j,\qquad Z^+=Z+S,\label{L:eq:formal1}\\
 Y_j&=g_{j+1}-g_j+\mu\tau S,\label{L:eq:formal2}\\
 M^+&=\left(I-\frac{SY_j\T}{Y_j\T S}\right)M
 \left(I-\frac{Y_jS\T}{Y_j\T S}\right)
 +\frac{SS\T}{Y_j\T S}.
\end{align}
Gradient indices are cyclic and $\alpha=(1,1/4,1/16)$. Let the final quantities be $(Z_3,M_3)$ and let $r=(M_3)_{11}$. The normalized return map is
\begin{equation}\label{L:eq:return}
 \mathcal T_\mu(z,b,e,\tau)=
 \left(Z_3/r,\ (M_3)_{12}/r,\ (M_3)_{22}/r,\ \tau r\right).
\end{equation}
All denominators are nonzero at $P_*$, and $r=\rho>0$ there. This is a real-analytic map on an open neighborhood of $P_*$, allowing an artificial negative $\tau$ in that analytic extension. Only its positive-$\tau$ part represents actual BFGS.

\begin{lemma}[Agreement with the actual implementation]\label{L:lem:agreement}
There is an open neighborhood $\mathcal N$ of $P_*$ such that for every state in $\mathcal N$ with $\tau>0$, the actual three consecutive full-BFGS iterations on $Q_\mu$ have exactly the certified trial lists and return map~\eqref{L:eq:return}.
\end{lemma}
\begin{proof}
At a query $\tau Z$, the common quadratic does not affect which slope is active. Its function value divided by $\tau$ is
\[
 G(Z)+\tfrac{\mu\tau}{2}\norm Z^2,
\]
and its gradient is $g_i+\mu\tau Z$. All Armijo and weak-Wolfe comparisons, divided by their common positive scale when needed, are continuous in the normalized state. The seed has strict branch gaps and strict decisions at its finite list of nine queries. Intersecting the corresponding open neighborhoods preserves the entire decision tree. The normalized secant and update formulas are exactly~\eqref{L:eq:formal1}--\eqref{L:eq:formal2}; no reset of the actual matrix is made.
\end{proof}


\setcounter{section}{3}
\setcounter{theorem}{2}
\setcounter{lemma}{1}
\setcounter{corollary}{0}
\setcounter{proposition}{0}
\setcounter{remark}{0}
\setcounter{equation}{2}
\section{A certified hyperbolic fixed point}
At $\tau=0$ the matrix return depends only on $(b,e)$, not on $z$. Its normalized shape map is
\[
 (b,e)\longmapsto ((M_3)_{12}/(M_3)_{11},(M_3)_{22}/(M_3)_{11}).
\]
Let $J$ be its derivative at $(0,1)$. The supplied rational-interval automatic-differentiation program computes rigorous enclosures of all entries; their decimal displays are
\[
 J\simeq\begin{pmatrix}
 1.112398625073290&0.645970102731598\\
 0.751501570427499&0.490513688583175
 \end{pmatrix}.
\]
More importantly, using exact rational interval signs of
$P_J(\lambda)=\lambda^2-\operatorname{tr}(J)\lambda+\det J$, it proves
\begin{equation}\label{L:eq:eigshape}
 0.038<\lambda_s<0.039,\qquad 1.564<\lambda_u<1.565.
\end{equation}
Specifically $P_J(0.038)>0$, $P_J(0.039)<0$, $P_J(1.564)<0$, and $P_J(1.565)>0$ throughout the algebraic root box. The two disjoint sign-change intervals contain the two distinct real roots. No floating-point eigensolver is used for this assertion.

The full derivative has the analytically determined block-triangular form
\begin{equation}\label{L:eq:jacobian}
 D\mathcal T_\mu(P_*)=
 \begin{pmatrix}
 \rho^{-1}I_2&C&\mu a\\
 0&J&\mu b\\
 0&0&\rho
 \end{pmatrix}.
\end{equation}
The unneeded blocks $C,a,b$ are finite. The upper left block follows because, at $\tau=0$, $Z_3=z+v(b,e)$ and $r$ is independent of $z$. The last row follows from $\tau^+=\tau r$.

The eigenvalues are therefore
\[
 \rho^{-1},\ \rho^{-1},\ \lambda_s,\ \lambda_u,\ \rho.
\]
None lies on the unit circle and none is zero. There are exactly two stable eigenvalues. The derivative is invertible, so the inverse function theorem makes $\mathcal T_\mu$ a local diffeomorphism.

Crucially $\rho$ is different from every eigenvalue of the upper left $4\times4$ block. Thus its eigenvector can be chosen with last coordinate one: solve the top four linear equations after assigning $v_\tau=1$. The stable subspace is not contained in the artificial boundary $\{\tau=0\}$.

\begin{proof}[Proof of Theorem~\ref{L:thm:global}]
Apply the local stable-manifold theorem for a hyperbolic fixed point of a smooth local diffeomorphism; see~\cite[Sections 3--4]{DY}. It gives a smooth two-dimensional local stable manifold tangent to the stable eigenspace at $P_*$. Since this tangent space has a vector with nonzero $\tau$ component, the stable manifold contains points with arbitrarily small positive $\tau$.

Choose such a point sufficiently close that its entire forward return-map orbit remains in the neighborhood of Lemma~\ref{L:lem:agreement} and converges to $P_*$. The positive half-space is preserved because $\tau^+=\tau r$ and $r>0$ in this neighborhood. At every return, $M(b,e)$ is positive definite and close to $I$; all actual substep matrices are positive definite as well.

The return-map orbit represents the actual deterministic full-BFGS orbit, with accepted lengths $1,1/4,1/16$. Its physical states satisfy
\[
 x_{3m}=\tau_m z_m\to0,\qquad H_{3m}=\tau_mM(b_m,e_m)\to0.
\]
The intermediate accepted states and all rejected queries are $\tau_m$ times bounded continuous functions of the normalized state, so they tend to zero too. They retain strict branch gaps and strict decisions. Their objective values are positive near the seed and decrease to zero. Equation~\eqref{L:eq:minimizer} and $0\notin\partial Q_\mu(0)$ prove actual nonstationary convergence.

The same argument shows more: given any neighborhood of physical zero, the initial stable-manifold point can be chosen so that \emph{all queries of the orbit} lie inside it. This small-orbit property will be used for the embedding. The initial matrix is supplied by the stable manifold, not assumed to be $I$ in these particular coordinates.
\end{proof}
\noindent This is an existence proof of an actual trajectory, not a claim that a numerically truncated stable eigenvector produces it forever. Expanding eigenvalues make ordinary floating-point forward shooting an inappropriate certificate of membership in the stable manifold.


\setcounter{section}{4}
\setcounter{theorem}{2}
\setcounter{lemma}{1}
\setcounter{corollary}{0}
\setcounter{proposition}{0}
\setcounter{remark}{0}
\setcounter{equation}{4}
\section{A local accumulation embedding with positive lower curvature}
We now construct a single objective satisfying the literal Hessian bounds and having no guaranteed local basin. Fix
\[
 \mu=1/4096,\quad c_*=(2,-32)\T,\quad
 q_i=g_i-\mu c_*,\quad b_i=q_i\T c_*,\quad r=1/128,\quad a_n=r^n.
\]
The seed interval certificate proves
\[
 b_0=b_1=1791/1024,\qquad 1.7<b_i<34.
\]
Only two distinct values of $b_i$ occur. For either value write
\[
 \vartheta_b=\frac{1-(1+r)/b}{1-r},\qquad \epsilon=1/1024.
\]
Define a monotone $C^1$ ramp $P_b$ on $[0,1]$: it is zero for $\xi\le\vartheta_b-\epsilon$, one for $\xi\ge\vartheta_b+\epsilon$, and on the intervening interval equals
\[
 P_b(\xi)=3t^2-2t^3,\qquad t=\frac{\xi-\vartheta_b+\epsilon}{2\epsilon}.
\]
Let $J_b(\xi)=\int_0^\xi P_b(v)\,dv$. The elementary formulas give
\[
 \int_0^1P_b=1-\vartheta_b,\qquad 0\le P_b'\le768.
\]
On $[ba_{n+1},ba_n]$ put
\begin{align}
 \xi&=\frac{x/(ba_n)-r}{1-r},\\
 \varphi_b(x)&=(ra_n)^2+ba_n^2(1-r)
 \{r\xi+(1-r)J_b(\xi)\}.\label{L:eq:scalar}
\end{align}
Complete the definition by $\varphi_b(0)=0$, $\varphi_b(x)=x^2/2$ for $x<0$, and $\varphi_b(x)=1+x-b$ for $x\ge b$.

\begin{lemma}[Scalar regularity and exact affine neighborhoods]\label{L:lem:scalar}
Each $\varphi_b$ is convex, nonnegative, coercive, uniquely minimized at zero, and globally $C^{1,1}$. It is $C^2$ on $(0,\infty)$ and on $(-\infty,0)$, but is not twice differentiable at zero. At every node,
\[
 \varphi_b(ba_n)=a_n^2,\quad \varphi_b'(ba_n)=a_n,
\]
and
\begin{equation}\label{L:eq:plateau}
 \varphi_b(ba_n+h)=a_n^2+a_nh\qquad (|h|\le a_n/2).
\end{equation}
Where its second derivative exists, $0\le\varphi_b''\le\max\{1,768/b\}$.
\end{lemma}
\begin{proof}
On the positive intervals,
\[
 \varphi_b'=a_n\{r+(1-r)P_b(\xi)\},\qquad
 \varphi_b''=P_b'(\xi)/b.
\]
The defining value of $\vartheta_b$ gives
$r+(1-r)(1-\vartheta_b)=(1+r)/b$, which proves value matching. The derivative rises monotonically from $ra_n$ to $a_n$ and is constant near the endpoints. Thus values and derivatives match at every node. At zero $\varphi_b'=O(x)$ and $\varphi_b=O(x^2)$, giving differentiability with derivative zero. Uniform bounded derivative slopes on all intervals imply a global Lipschitz bound, including across the accumulation point by taking limits. The negative and outer extensions match in value and first derivative.

The constant-derivative lengths on the two sides of $ba_n$, divided by $a_n$, are
\[
 b(1-r)(1-\vartheta_b-\epsilon),\qquad
 b(1-r)(\vartheta_b-\epsilon)/r.
\]
Both exceed $1/2$ for the certified values of $b$. This proves~\eqref{L:eq:plateau}. At zero, along $x=ba_n$ the ratio $\varphi_b'(x)/x$ equals $1/b$, whereas along $x=(b+1/4)a_n$ it equals $1/(b+1/4)$. Both sequences tend to zero, so $\varphi_b''(0)$ does not exist. This is the deliberate regularity distinction from the previous $C^{2,1}$ embedding.
\end{proof}

Define four functions and their maximum:
\begin{equation}\label{L:eq:embedding}
 f_i(x)=\tfrac\mu2\norm x^2+\varphi_{b_i}(q_i\T x)\quad (i=0,1,2),\qquad
 f_3(x)=\tfrac\mu2\norm x^2+x_2,\qquad F=\max_i f_i.
\end{equation}
Each branch is globally $\mu$-strongly convex and has a globally Lipschitz gradient. The rational interval verifier proves the common gradient-Lipschitz constant $512$: the first three Hessian upper bounds are less than $439,446,27$, respectively, and the fourth has Hessian $\mu I$. Thus the maximum is $\mu$-strongly convex.

Since the first three scalar terms are nonnegative,
\[
 F(x)\ge\tfrac\mu2\norm x^2,\qquad F(0)=0.
\]
This proves compact sublevels and the unique minimum at zero. The first three gradients vanish there and the fourth is $(0,1)\T$; consequently $F'(0;h)=\max\{0,h_2\}$ and the minimizer is nonsmooth.

At any differentiable point of a finite maximum of $L$-smooth functions, all active gradients agree. Inactive pieces remain inactive locally, and the upper smooth Taylor bounds for the active pieces therefore give a local upper Taylor bound $L\norm h^2/2$ for the maximum. Together with strong convexity, this proves
\begin{equation}\label{L:eq:Hbounds}
 \mu I\preceq\nabla^2F(x)\preceq512I
 \quad\text{whenever }F\text{ is twice differentiable}.
\end{equation}
This conclusion does not assume the maximum has a globally Lipschitz gradient.


\setcounter{section}{5}
\setcounter{theorem}{2}
\setcounter{lemma}{2}
\setcounter{corollary}{0}
\setcounter{proposition}{0}
\setcounter{remark}{0}
\setcounter{equation}{7}
\section{Copying one perturbed bad orbit at every scale}
Set $p_n=a_nc_*$ and $C_*=1+\mu\norm{c_*}^2/2$. For any $y$ with $|q_i\T y|\le1/2$,
\begin{align*}
 f_i(p_n+a_ny)
 &=a_n^2\left\{\tfrac\mu2\norm{c_*+y}^2+1+q_i\T y\right\}\\
 &=a_n^2\{C_*+g_i\T y+\tfrac\mu2\norm y^2\}.
\end{align*}
The fourth branch is negative relative to its quadratic part there if $y$ is sufficiently small and $n\ge1$. Hence, on a fixed small neighborhood of $y=0$,
\begin{equation}\label{L:eq:local-model}
 F(p_n+a_ny)=a_n^2\{C_*+Q_\mu(y)\}.
\end{equation}
Use Theorem~\ref{L:thm:global} to choose one entire bad orbit $\{\bar x_k,\bar H_k\}$ of $Q_\mu$ whose accepted and rejected queries are all in that neighborhood. The same orbit can be made as small as needed, by the last part of its proof. Set
\[
 x_0^{(n)}=p_n+a_n\bar x_0,\qquad H_0^{(n)}=\bar H_0.
\]
Every query in the $n$th run is exactly $p_n+a_n$ times the corresponding seed query, and
\[
 x_k^{(n)}=p_n+a_n\bar x_k,\qquad H_k^{(n)}=\bar H_k.
\]
Indeed, function differences and directional-derivative products scale by $a_n^2$, while both gradients and steps scale by $a_n$. In the inverse-BFGS formula both scale factors cancel. The line search therefore makes the same decisions, including all rejections. Equation~\eqref{L:eq:local-model} applies to every query by the choice of the whole seed orbit. Induction proves the claim without numerical approximation.

Since $\bar x_k\to0$,
\[
 x_k^{(n)}\to p_n,\qquad F(x_k^{(n)})\downarrow a_n^2C_*>0.
\]
At $p_n$ the three active gradients are $a_ng_i$, so $0\notin\partial F(p_n)$. Meanwhile $x_0^{(n)}\to0$ and the supremum of the norms of all queries is at most
\[
 a_n\left(\norm{c_*}+\sup_{\text{seed queries}}\norm{\bar x}\right)\to0.
\]
The initial matrix is one fixed SPD matrix for every $n$.

\paragraph{Fixing identity initialization.}
Let $A=\bar H_0^{1/2}$ and define $\widetilde F(z)=F(Az)$. Full-BFGS affine invariance preserves the complete line-search decisions and gives initial matrix $I$ in the $z$ coordinates. The branch strong-convexity and gradient-Lipschitz constants become, for example,
\[
 \underline\mu=\mu\lambda_{\min}(\bar H_0)>0,\qquad
 \overline L=512\lambda_{\max}(\bar H_0)<\infty.
\]
These constants are fixed for the one transformed objective, not changed between starts. Its bad initial and limit points approach zero, and all queries remain arbitrarily close. This proves the four-$C^{1,1}$-branch part of Theorem~\ref{L:thm:local}.


\setcounter{section}{6}
\setcounter{theorem}{2}
\setcounter{lemma}{2}
\setcounter{corollary}{0}
\setcounter{proposition}{0}
\setcounter{remark}{0}
\setcounter{equation}{8}
\section{The literal Hessian condition permits an even simpler variant}
Replace the smooth ramp by the step $P_b(\xi)=\mathbf1_{\{\xi>\vartheta_b\}}$ and its primitive by $J_b(\xi)=(\xi-\vartheta_b)_+$. Use~\eqref{L:eq:scalar} on the positive intervals, the negative extension $\varphi_b(x)=-x$ for $x<0$, and the same affine positive outer extension. This scalar function is convex and piecewise affine with countably many pieces; it has the same node values, node slopes, and affine neighborhoods~\eqref{L:eq:plateau}. Thus the entire local-copy construction and the same fixed initial matrix work unchanged.

For this variant $F=\mu\norm{x}^2/2+p(x)$, with convex countably piecewise-affine $p$. At a nonzero point, any branch with $q_i\T x=0$ is strictly below the maximum: its scalar value is zero, whereas at least one of $q_0\T x,q_1\T x$ is nonzero and has positive scalar value. Here $q_0,q_1$ are independent. All potentially active scalar arguments are therefore nonzero and encounter only finitely many breakpoints in a small neighborhood. Hence $p$ is locally finite-polyhedral near every nonzero point. At any differentiable such point it is locally affine. The origin is nonsmooth. Consequently
\[
 \nabla^2F(x)=\mu I
 \quad\text{at every twice-differentiable point}.
\]
After the fixed affine transformation making $H_0=I$, that Hessian is the one fixed SPD matrix $\mu\bar H_0$. This establishes the remaining version of Theorem~\ref{L:thm:local}. No claim of four smooth branches is made for this alternative.


\endgroup

% END original chapters/L.tex

\StartModule{E}{Fresh verification and reproducibility}{New isolated reruns of the preserved source scripts. The status of finite algebra, rigorous intervals, numerical experiments and analytic proof steps is separated explicitly.}

% BEGIN original chapters/verification.tex
\textbf{Fresh audit date: 10 September 2026. All 19 requested jobs completed with exit code zero and PASS status.} These are reruns of preserved source scripts in isolated writable directories. They are not an independent mathematical review. The exact commands, script hashes, environment and output paths are recorded in \path{verification/fresh/SUMMARY.json}.

\section{What the checks mean}\label{what-the-checks-mean}

Four distinct evidential levels must not be conflated. Symbolic identities and exact rational arithmetic prove the specific finite statements that they actually compute. Rational interval enclosures additionally certify continuum assertions such as an isolated root or strict signs over a box, when paired with the stated contraction or sign argument. High-precision non-interval computations are numerical regressions, not rigorous enclosures. Infinite-trajectory, compact-cluster, geometric joining and stable-manifold conclusions require their analytic proofs; running many iterations does not replace those arguments.

The negative program and positive program have independent foundations. The deterministic negative seed does not use the principal-minor lemma. The positive-curvature negative extension adds a hyperbolicity certificate and a standard stable-manifold theorem. The positive convergence family uses the principal-minor hierarchy and comparison energies. No published theorem for a modified or limited-memory algorithm is used as a substitute for full BFGS.

\section{Fresh execution register}\label{fresh-execution-register}

\begin{longtable}[]{@{}
  >{\raggedright\arraybackslash}p{(\columnwidth - 6\tabcolsep) * \real{0.3400}}
  >{\raggedright\arraybackslash}p{(\columnwidth - 6\tabcolsep) * \real{0.4500}}
  >{\centering\arraybackslash}p{(\columnwidth - 6\tabcolsep) * \real{0.0900}}
  >{\raggedleft\arraybackslash}p{(\columnwidth - 6\tabcolsep) * \real{0.1200}}@{}}
\toprule\noalign{}
\begin{minipage}[b]{\linewidth}\raggedright
Job
\end{minipage} & \begin{minipage}[b]{\linewidth}\raggedright
Arithmetic / role
\end{minipage} & \begin{minipage}[b]{\linewidth}\centering
Result
\end{minipage} & \begin{minipage}[b]{\linewidth}\raggedleft
Seconds
\end{minipage} \\
\midrule\noalign{}
\endhead
\bottomrule\noalign{}
\endlastfoot
\path{D-root} & Exact rational intervals & PASS & 0.723 \\
\path{D-symbolic} & Symbolic identities & PASS & 6.196 \\
\path{D-full-oracle} & Exact rational intervals & PASS & 1.041 \\
\path{L-hyperbolicity} & Exact rational intervals & PASS & 0.797 \\
\path{L-embedding} & Exact rational intervals & PASS & 1.215 \\
\path{A-minors} & Exact rationals & PASS & 2.981 \\
\path{H-lifting} & Exact rationals & PASS & 0.917 \\
\path{M-mixed} & Exact rationals & PASS & 2.070 \\
\path{P-polyhedral} & Exact rationals & PASS & 1.224 \\
\path{R-secant} & Exact rationals & PASS & 1.129 \\
\path{R-interpretation} & Exact rationals & PASS & 0.780 \\
\path{G-obstruction} & Exact rationals & PASS & 0.640 \\
\path{V-ridges} & Exact rationals; auxiliary proof & PASS & 2.677 \\
\path{T-certificate} & Rationals plus decimal logs & PASS & 1.386 \\
\path{C-planar} & Symbolic / exact rationals & PASS & 1.217 \\
\path{U-multidimensional} & Exact rationals & PASS & 0.935 \\
\path{C-planar-numerical} & 140-digit, not interval & PASS & 3.229 \\
\path{U-multidimensional-numerical} & 100-digit, not interval & PASS & 1.993 \\
\path{G-oscillatory-numerical} & 220-digit, not interval; auxiliary & PASS & 0.712 \\
\end{longtable}

\section{What each suite actually checks}\label{what-each-suite-actually-checks}

\subsection{D-root}\label{d-root}

Existence/uniqueness of the algebraic seed; every accepted and rejected trial; matrix positivity; embedding constants. Raw transcript: \path{verification/fresh/D-root/stdout.txt}; the job entry in \texttt{SUMMARY.json} identifies every generated JSON result and the executed command.

\subsection{D-symbolic}\label{d-symbolic}

Independent inverse/Hessian update closure and secant identities. Raw transcript: \path{verification/fresh/D-symbolic/stdout.txt}; the job entry in \texttt{SUMMARY.json} identifies every generated JSON result and the executed command.

\subsection{D-full-oracle}\label{d-full-oracle}

45 accepted updates and 135 actual trials across five shrinking identity-initialized starts; complete global oracle. Raw transcript: \path{verification/fresh/D-full-oracle/stdout.txt}; the job entry in \texttt{SUMMARY.json} identifies every generated JSON result and the executed command.

\subsection{L-hyperbolicity}\label{l-hyperbolicity}

Normalized Jacobian, characteristic-polynomial signs, stable/unstable eigenvalue separation, embedding bounds. Raw transcript: \path{verification/fresh/L-hyperbolicity/stdout.txt}; the job entry in \texttt{SUMMARY.json} identifies every generated JSON result and the executed command.

\subsection{L-embedding}\label{l-embedding}

336 affine-plateau checks, 228 ramp checks, 72 full local value/gradient identities. Raw transcript: \path{verification/fresh/L-embedding/stdout.txt}; the job entry in \texttt{SUMMARY.json} identifies every generated JSON result and the executed command.

\subsection{A-minors}\label{a-minors}

125 update cases, 1,425 compressed identities and 375 minor-hierarchy inequalities. Raw transcript: \path{verification/fresh/A-minors/stdout.txt}; the job entry in \texttt{SUMMARY.json} identifies every generated JSON result and the executed command.

\subsection{H-lifting}\label{h-lifting}

100 projected-update and invariant-subspace cases. Raw transcript: \path{verification/fresh/H-lifting/stdout.txt}; the job entry in \texttt{SUMMARY.json} identifies every generated JSON result and the executed command.

\subsection{M-mixed}\label{m-mixed}

425 accepted-step cases; 1,030 minor inequalities; 425 weighted-potential checks; 684 compression identities; matrix-limit example. Raw transcript: \path{verification/fresh/M-mixed/stdout.txt}; the job entry in \texttt{SUMMARY.json} identifies every generated JSON result and the executed command.

\subsection{P-polyhedral}\label{p-polyhedral}

144 accepted steps, including 72 with a C1,1 but non-C2 smooth part; 40 lifting checks. Raw transcript: \path{verification/fresh/P-polyhedral/stdout.txt}; the job entry in \texttt{SUMMARY.json} identifies every generated JSON result and the executed command.

\subsection{R-secant}\label{r-secant}

1,500 finite-max pair checks, sum/composition checks, 825 powered-gauge cases including 165 fractional-power cases, 90 accepted steps. Raw transcript: \path{verification/fresh/R-secant/stdout.txt}; the job entry in \texttt{SUMMARY.json} identifies every generated JSON result and the executed command.

\subsection{R-interpretation}\label{r-interpretation}

600 endpoint-pair cases and 1,800 inward Taylor inequalities; supporting-plane, boundary and strictness examples. Raw transcript: \path{verification/fresh/R-interpretation/stdout.txt}; the job entry in \texttt{SUMMARY.json} identifies every generated JSON result and the executed command.

\subsection{G-obstruction}\label{g-obstruction}

24 common-target obstruction cases, 117 SPD step realizations, 60 accumulating-kink cases, 150 rank-one feasibility checks. Raw transcript: \path{verification/fresh/G-obstruction/stdout.txt}; the job entry in \texttt{SUMMARY.json} identifies every generated JSON result and the executed command.

\subsection{V-ridges}\label{v-ridges}

370 accepted steps, 263 crossings, 107 same-mode steps, 270 asymmetric-slope cases, 90 degenerate cases, 80 matrix-frame checks. Preserved as corroboration of a now-subsumed theorem. Raw transcript: \path{verification/fresh/V-ridges/stdout.txt}; the job entry in \texttt{SUMMARY.json} identifies every generated JSON result and the executed command.

\subsection{T-certificate}\label{t-certificate}

96 projected-PSD energy cases, 384 accepted steps, 36 strictness examples. 208 logarithmic rescaling checks use 70-digit Decimal arithmetic and are not exact interval certificates. Raw transcript: \path{verification/fresh/T-certificate/stdout.txt}; the job entry in \texttt{SUMMARY.json} identifies every generated JSON result and the executed command.

\subsection{C-planar}\label{c-planar}

Angular and determinant identities; 90 full-BFGS moving-energy cases; explicit quadratic barycentric/basin checks; 3D field obstruction. Raw transcript: \path{verification/fresh/C-planar/stdout.txt}; the job entry in \texttt{SUMMARY.json} identifies every generated JSON result and the executed command.

\subsection{U-multidimensional}\label{u-multidimensional}

72 full-matrix residual-energy cases, 150 cutoff transports, 87 minor inequalities, 18 finite-horizon tail checks, 23 spike-node identities, 161 infinite-series oracle cases. Raw transcript: \path{verification/fresh/U-multidimensional/stdout.txt}; the job entry in \texttt{SUMMARY.json} identifies every generated JSON result and the executed command.

\subsection{C-planar-numerical}\label{c-planar-numerical}

18 actual full-BFGS runs; 3,084 updates, 7,381 trials; all reached F \textless{} 1e-45. Its repeated exact stage is not counted twice. Raw transcript: \path{verification/fresh/C-planar-numerical/stdout.txt}; the job entry in \texttt{SUMMARY.json} identifies every generated JSON result and the executed command.

\subsection{U-multidimensional-numerical}\label{u-multidimensional-numerical}

Nine actual full-BFGS runs in dimensions 3, 5 and 8; 495 updates and 1,050 trials; all reached F \textless{} 1e-28. Complete infinite-series oracle; repeated exact stage not counted twice. Raw transcript: \path{verification/fresh/U-multidimensional-numerical/stdout.txt}; the job entry in \texttt{SUMMARY.json} identifies every generated JSON result and the executed command.

\subsection{G-oscillatory-numerical}\label{g-oscillatory-numerical}

24 finite incoming-pair checks for the source oscillatory example. Historical auxiliary geometry, not a main BFGS counterexample. Raw transcript: \path{verification/fresh/G-oscillatory-numerical/stdout.txt}; the job entry in \texttt{SUMMARY.json} identifies every generated JSON result and the executed command.

\section{Analytic steps that still require independent review}\label{analytic-steps-that-still-require-independent-review}

\begin{longtable}[]{@{}
  >{\raggedright\arraybackslash}p{(\columnwidth - 4\tabcolsep) * \real{0.3333}}
  >{\raggedright\arraybackslash}p{(\columnwidth - 4\tabcolsep) * \real{0.3333}}
  >{\raggedright\arraybackslash}p{(\columnwidth - 4\tabcolsep) * \real{0.3333}}@{}}
\toprule\noalign{}
\begin{minipage}[b]{\linewidth}\raggedright
Result or proof component
\end{minipage} & \begin{minipage}[b]{\linewidth}\raggedright
Supplied evidence
\end{minipage} & \begin{minipage}[b]{\linewidth}\raggedright
What was not established by finite reruns
\end{minipage} \\
\midrule\noalign{}
\endhead
\bottomrule\noalign{}
\endlastfoot
Principal-minor hierarchy & Complete Gram-determinant proof; exact algebra cases & A third-party audit of the general proof \\
Reciprocal growth and the superquadratic tail theorem & Finite-horizon algebra and explicit ordered-limit proof & Any empirical inference from a finite horizon to an infinite limit \\
Transported certificate and radial equivalence & Full proof and identities & Uniform hypotheses for an arbitrary unlisted objective \\
Planar affine-critical theorem & Geometric joining proof; exact local identities; numerical example & Independent validation of the compact-neighborhood, joining and all-segment arguments \\
Fixed-normal multidimensional theorem & Residual-energy and cutoff proof; exact identities; full-series numerical oracle & A theorem for rotating normals or arbitrary curved active manifolds \\
Deterministic C2,1 local counterexample & Isolated algebraic root, symbolic closure, strict all-trial decisions, exact interpolation and containment induction & A positive-measure basin, robustness to floating-point perturbations, or a different line search \\
Positive-curvature counterexamples & Certified hyperbolicity plus the analytic return map and stable-manifold theorem & A numerical starting vector proved to lie exactly on the stable manifold \\
Semialgebraic affine-critical gap & Source-derived geometric proof using standard tame-geometric facts & A blanket semialgebraic BFGS convergence theorem \\
\end{longtable}

All retained results remain working mathematical claims pending independent review. A successful source script can validate a mistaken formula if the same mistake occurs in its construction and checks; independent inverse/Hessian forms and separate symbolic/interval routes reduce, but do not eliminate, that concern. No publication-priority certification is made.

\section{Reproduce the audit}\label{reproduce-the-audit}

From the package root, with Python 3.13 (or a compatible Python), SymPy 1.14.0 and mpmath 1.3.0:

\begin{verbatim}
python -m pip install -r requirements.txt
python verification/run_all.py --workers 3 --timeout 600
python verification/run_all.py --numerical --workers 3 --timeout 600
\end{verbatim}

The first command runs the non-optional algebraic suites; note that the transported source suite includes explicitly labeled decimal-log checks. The second includes the optional high-precision runs. Each selected job receives a new working copy of its module, so scripts may write their normal outputs without modifying \path{verification/modules/}. Existing fresh logs are output artifacts, not the source inputs; preserve them before a new audit when comparison is desired.

The exact environment for this run was:

\begin{verbatim}
python: 3.13.5 (main, Jul 15 2026, 20:25:40) [GCC 14.2.0]
platform: Linux-6.18.35-x86_64-with-glibc2.41
sympy: 1.14.0
mpmath: 1.3.0
\end{verbatim}

No verification job requires network access. Scripts under \path{verification/modules/} preserve the source versions; \path{verification/historical/} contains the old logs, while \path{verification/fresh/} contains this execution. Counts from old files are never represented as fresh checks, and exact stages repeated inside numerical scripts are not double-counted. The per-job script hash and the complete package manifest support later comparisons.

\section{Rebuild the report}\label{rebuild-the-report}

\begin{verbatim}
cd report
latexmk -pdf -interaction=nonstopmode -halt-on-error master.tex
\end{verbatim}

The editable report sources are complete. \path{report/tools/assemble.py} reconstructs selected proof modules from the preserved source snapshots; it is an editorial tool, not part of a theorem verifier. The authored common foundation, structural corollaries, affine-critical geometry, inventory and audit are separate editable chapters. Details of intentional changes appear in \texttt{EDITORIAL\_NOTES.md} and \path{report/EDIT_MANIFEST.json}.

% END original chapters/verification.tex

\backmatter
\renewcommand{\ModuleTag}{B}
\cleardoublepage
\phantomsection
\addcontentsline{toc}{chapter}{References and source provenance}

% BEGIN original bibliography.tex
\begin{thebibliography}{SCURVE}
\bibitem[XW]{XW} Y. Xie and A. W\"achter. \emph{On the convergence of BFGS on a class of piecewise linear non-smooth functions}. arXiv:1712.08571v1 (2017). User-supplied paper; pre-existing algorithmic formulation and two-piece finite-escape theorem. Complete original PDF preserved in \path{provenance/Xie_Waechter_2017.pdf}.
\bibitem[LO]{LO} A. S. Lewis and M. L. Overton. \emph{Nonsmooth optimization via quasi-Newton methods}. Mathematical Programming 141 (2013), 135--163. DOI 10.1007/s10107-012-0514-2. Algorithm 4.6 is the fixed weak-Wolfe line search used by the negative construction. Source-cited background; no new literature audit was performed.
\bibitem[DY]{DY} S. Dyatlov. \emph{Notes on hyperbolic dynamics}. arXiv:1805.11660 (2018), Sections 3--4. Standard local stable-manifold theorem used by Module L.
\bibitem[WW]{WW} D. Wachsmuth and G. Wachsmuth. \emph{A simple proof of the Baillon--Haddad theorem on open subsets of Hilbert spaces}. arXiv:2204.00282 (2022). Convex cocoercivity on an open convex domain.
\bibitem[BG]{BG} B. Gebken. \emph{Technical Results on the Convergence of Quasi-Newton Methods for Nonsmooth Optimization}. JOTA 209, article 57 (2026), DOI 10.1007/s10957-026-03008-z; arXiv:2511.03296. Source-cited related work with iterate/spectral hypotheses, not a proof of the new working claims.
\bibitem[LW]{LW} A. S. Lewis and C. Wylie. \emph{A simple Newton method for local nonsmooth optimization}. arXiv:1907.11742. Related smooth-max nondegeneracy for a different algorithm.
\bibitem[BDLS]{BDLS} J. Bolte, A. Daniilidis, A. S. Lewis and M. Shiota. \emph{Clarke subgradients of stratifiable functions}. arXiv:math/0601530. Source-cited projection and definable-stratification background.
\bibitem[LL]{LL} A. S. Lewis and J. Liang. \emph{Partial smoothness and constant rank}. arXiv:1807.03134. Source-cited tangent/normal geometry; not an ordinary-BFGS convergence theorem.
\bibitem[BB]{BB} A. Banerjee, S. Merugu, I. S. Dhillon and J. Ghosh. \emph{Clustering with Bregman divergences}. Journal of Machine Learning Research 6 (2005), 1705--1749. Source-cited supporting-plane-error terminology only.
\bibitem[SM]{SM} Working report, \emph{Full BFGS on mixed smooth--nonsmooth convex objectives}. Source snapshot: \path{provenance/mixed_ridge/}. Common principal-minor and fixed-reference energy foundation.
\bibitem[SH]{SH} Working report, \emph{Full BFGS on homogeneous convex objectives}. Source: \path{provenance/general_homogeneous/}. Unconditional value dichotomy, affine-hull escape and low-rank iterates. Earlier rational counterexample is superseded by Module D.
\bibitem[ST]{ST} Working report, \emph{A weaker objective-level secant condition for full BFGS}. Source: \path{provenance/transported_secant/}. Objective certificate and arbitrary-many-branch parabolic proof.
\bibitem[SG]{SG} Working report, \emph{What does a transported secant certificate mean?} Source: \path{provenance/certificate_geometry/}. Exact normalization, pointwise feasibility, obstructions and scalar nonnecessity.
\bibitem[SR]{SR} Working report, \emph{Radial smoothness, supporting-plane consistency, and local full-BFGS convergence}. Source: \path{provenance/radial_interpretation/}. RS equivalence and finite-polyhedral-partition proof.
\bibitem[SL]{SL} Working report, \emph{Local convergence of full BFGS from objective geometry}. Source: \path{provenance/local_general/}. Earlier radial-secant and powered-gauge deductions, retained through stronger results.
\bibitem[SP]{SP} Working report, \emph{Local convergence of full BFGS for smooth convex plus polyhedral objectives}. Source: \path{provenance/local_polyhedral/}. Local reduction, certified basin, affine-fiber lifting and well-definedness distinctions.
\bibitem[SQ]{SQ} Working note, \emph{Next BFGS research directions, with finite-horizon rate corollaries}. Source: \path{provenance/rates/}. Only proved rate and finite-horizon deductions are operative; old open-target lists are historical.
\bibitem[SE]{SE} Working note, \emph{BFGS after the convex nonsmooth counterexample}, Section 3. Source: \path{provenance/gauge_transform/}. Increasing smooth transforms of coercive gauges.
\bibitem[SV]{SV} Working report, \emph{Moving-frame energies for full BFGS}. Source: \path{provenance/moving_frames/}. Its general scalar-ridge function theorem is now a corollary of Module C, as proved in the inventory. Independent proof and verifier retained as provenance.
\bibitem[SA]{SA} Working note, \emph{Affine criticality as a possible basis for local full-BFGS convergence}. Source: \path{provenance/affine_criticality/}. Geometry, semialgebraic affine-critical values and explicitly unresolved general-dimensional conjecture.
\bibitem[SC]{SC} Working report, \emph{Forced affine-critical cluster points for planar full BFGS}. Source: \path{provenance/affine_cluster/}. Complete planar joining, transport, cluster and local smooth-max arguments.
\bibitem[SU]{SU} Working report, \emph{Full BFGS in arbitrary dimension: a smooth-direction transport budget and parallel-interface convergence}. Source: \path{provenance/multidimensional/}. Superquadratic reciprocal theorem, imperfect frame transport, fixed-normal function theorems and exact spike oracle.
\bibitem[SD]{SD} Working report, \emph{Failure of local convergence for full BFGS with a unit-first Armijo--weak-Wolfe line search}. Source: \path{provenance/deterministic_bracketing/}. Canonical algebraic seed, all-query decisions and globally $C^{2,1}$ upper-regularity local embedding.
\bibitem[SLC]{SLC} Working report, \emph{Positive classical Hessian bounds do not eliminate nonsmooth full-BFGS counterexamples}. Source: \path{provenance/positive_lower_curvature/}. Hyperbolicity, stable-manifold persistence and lower-curvature local embeddings.
\end{thebibliography}

% END original bibliography.tex

\end{document}
